\documentclass[11pt, letterpaper]{article}

\usepackage[top=.8in, bottom=.8in, left=0.5in, right=0.5in]{geometry}
\usepackage{amsmath, amssymb, amsthm, mathtools}
\usepackage{longtable}
\usepackage{booktabs}
\usepackage{array}
\usepackage{tabularx}
\usepackage{xcolor}
\usepackage{fancyhdr}
\usepackage{titlesec}
\usepackage{setspace}
\usepackage[T1]{fontenc}
\usepackage{enumitem}
\usepackage{tocloft}
\usepackage{tcolorbox}
\tcbuselibrary{breakable, skins}

\definecolor{RSblue}{RGB}{25, 60, 120}
\definecolor{RSgray}{RGB}{90, 90, 90}
\definecolor{RSgreen}{RGB}{20, 110, 60}
\definecolor{RSred}{RGB}{160, 30, 30}
\definecolor{RSamber}{RGB}{180, 110, 0}
\definecolor{RSpurple}{RGB}{90, 30, 120}

\newtcolorbox{rsbox}{
  enhanced, breakable,
  colback=RSblue!4, colframe=RSblue!60,
  title={\textbf{RS Ontological Statement}},
  fonttitle=\small\bfseries\color{white},
  attach boxed title to top left={yshift=-2mm, xshift=4mm},
  boxed title style={colback=RSblue},
  left=4pt, right=4pt, top=6pt, bottom=4pt
}

\newtcolorbox{verifybox}{
  enhanced, breakable,
  colback=RSgreen!4, colframe=RSgreen!50,
  title={\textbf{External Verification (Continuous Lens)}},
  fonttitle=\small\bfseries\color{white},
  attach boxed title to top left={yshift=-2mm, xshift=4mm},
  boxed title style={colback=RSgreen!70!black},
  left=4pt, right=4pt, top=6pt, bottom=4pt
}

\newtcolorbox{observebox}{
  enhanced, breakable,
  colback=RSamber!5, colframe=RSamber!60,
  title={\textbf{Structural Observation}},
  fonttitle=\small\bfseries\color{white},
  attach boxed title to top left={yshift=-2mm, xshift=4mm},
  boxed title style={colback=RSamber!80!black},
  left=4pt, right=4pt, top=6pt, bottom=4pt
}


\pagestyle{fancy}
\fancyhf{}
\fancyhead[L]{\small\textcolor{RSgray}{RigbySpace Discrete Dynamics}}
\fancyhead[R]{\small\textcolor{RSgray}{Master Reference - Page \thepage}}
\fancyfoot[L]{\small\textcolor{RSgray}{D. Veneziano}}
\fancyfoot[R]{\small\textcolor{RSgray}{Discrete Unification Framework}}
\renewcommand{\headrulewidth}{0.4pt}

\titleformat{\section}
  {\large\bfseries\color{RSblue}}{\thesection}{1em}{}[\titlerule]
\titleformat{\subsection}
  {\normalsize\bfseries\color{RSblue}}{\thesubsection}{1em}{}
\titleformat{\subsubsection}
  {\normalsize\bfseries\color{RSgray}}{\thesubsubsection}{1em}{}

\onehalfspacing
\setlength{\parskip}{4pt}
\setlength{\parindent}{0pt}

\newcolumntype{L}[1]{>{\raggedright\arraybackslash}p{#1}}

\newcommand{\Z}{\mathbb{Z}}
\newcommand{\Sinter}{S_{\mathrm{int}}}
\newcommand{\SL}{S_{\mathrm{L}}}
\newcommand{\SLe}{S_{\mathrm{L}}^{\mathrm{ext}}}
\newcommand{\SP}{S_{\mathrm{P}}}
\newcommand{\SPs}{S_{\mathrm{P}}^{*}}
\newcommand{\SPe}{S_{\mathrm{P}}^{\mathrm{ext}}}
\newcommand{\IP}{I_{\mathrm{P}}}
\newcommand{\Vac}{\mathcal{V}}
\newcommand{\SLpos}{S_{\mathrm{L}}^{+}}
\newcommand{\GK}{\Gamma_{K}}
\newcommand{\Ethree}{E_{3+1}}
\newcommand{\SRS}{S_{\mathrm{RS}}}
\newcommand{\TRS}{\mathcal{T}_{\mathrm{RS}}}
\newcommand{\Ens}{\mathcal{E}}

\newcommand{\lop}{\lambda}
\newcommand{\aop}{\eta}

\newcommand{\bplus}{\boxplus}
\newcommand{\med}{\oplus}
\newcommand{\ecplus}{\oplus_{\mathrm{ec}}}
\renewcommand{\succ}{\operatorname{succ}}

\newcommand{\rank}{\rho}
\newcommand{\hfn}{h}
\newcommand{\Dcross}{\Delta_{\mathrm{cross}}}
\newcommand{\sgn}{\operatorname{sgn}}
\newcommand{\rem}{\operatorname{rem}}

\newcommand{\koppa}{\mathbin{\text{\rm\k{o}}}}

\newcommand{\Tcyc}{T_{11}}
\newcommand{\PhiSet}{\Phi}
\newcommand{\Tbary}{T_{\mathrm{bary}}}
\newcommand{\Psistate}{\Psi}
\newcommand{\Tcharge}{\Gamma}

\newcommand{\Dg}{D_{\mathrm{gauge}}}
\newcommand{\Sint}{S_{\mathrm{int}}}
\newcommand{\Bridge}{\mathcal{B}}
\newcommand{\Lalpha}{\alpha_{\mathrm{lat}}}
\newcommand{\LQ}{\mathcal{L}_{\mathbb{Z}}}
\newcommand{\Sigimb}{\Sigma_{\mathrm{imb}}}
\newcommand{\VS}{\mathcal{V}_S}
\newcommand{\PG}{\mathrm{PG}}

% --- New macros: RS Sommerfeld and Pi-Analog ---
\newcommand{\Gmg}{G_{\mathrm{mg}}}
\newcommand{\Ksat}{K_{\mathrm{sat}}}
\newcommand{\dslip}{\delta_{\mathrm{slip}}}
\newcommand{\Tvac}{T_{\mathrm{vac}}}
\newcommand{\tauX}{\tau_{10}}
\newcommand{\Pithree}{\Pi_3}
\newcommand{\nG}{n_G}
\newcommand{\aSC}{\alpha^{-1}_{\mathrm{int}}}
\newcommand{\piRS}{\pi_{\mathrm{RS}}}
\newcommand{\pa}{\mathfrak{p}_\alpha}
\newcommand{\qa}{\mathfrak{q}_\alpha}
\newcommand{\Da}{\mathcal{D}_\alpha}
\newcommand{\qosc}{\mathfrak{q}_{\mathrm{osc}}}

\theoremstyle{definition}
\newtheorem{definition}{Definition}[section]
\newtheorem{axiom}{Axiom}[section]
\newtheorem{remark}{Remark}[section]
\newtheorem{postulate}{Postulate}[section]
\newtheorem{law}{Law}[section]
\newtheorem{observation}{Observation}[section]

\theoremstyle{plain}
\newtheorem{theorem}{Theorem}[section]
\newtheorem{lemma}{Lemma}[section]
\newtheorem{proposition}{Proposition}[section]
\newtheorem{corollary}{Corollary}[section]


\begin{document}

\begin{titlepage}
  \centering
  \vspace*{2cm}
  {\Huge\bfseries\color{RSblue} RigbySpace Discrete Dynamics\par}
  \vspace{0.5cm}
  {\LARGE Master Reference \par}
  \vspace{0.5cm}
  {\large Notation $\cdot$ Symbols $\cdot$ Definitions $\cdot$ Glossary\par}
  \vspace{2cm}
  \rule{0.6\textwidth}{0.4pt}\par
  \vspace{1cm}
  {\large D.\ Veneziano\par}
  \vspace{0.8cm}
  {\normalsize Reference and Standards\par}
  \vspace{1.5cm}
  \rule{0.6\textwidth}{0.4pt}\par
  \vspace{0.5cm}
  {\small Propagation Scaffolding, Operators and Structural Formation\par}
  \vfill
  {\small\color{RSgray}
  RigbySpace is a formally consistent candidate theory internal to its own
  axioms.  No empirical equivalence to external physical theory is asserted
  unless an explicit mapping is provided and derived under that mapping.\par}
\end{titlepage}

\tableofcontents
\newpage


\section{Preface: Ontological Ground Rules}
\addcontentsline{toc}{section}{Preface: Ontological Ground Rules}

RigbySpace is a foundational physical framework whose substrate consists
exclusively of integers.  The following prohibitions are axiomatic and apply
uniformly across all documents in the corpus.

\medskip
\noindent\textbf{Forbidden primitives.}
Real numbers $\mathbb{R}$; limits to infinity; infinitesimals; continuous
manifolds; modular quotient arithmetic as an ontological primitive; the absolute
value operator $|\cdot|$ as an ontological primitive; the logarithm; Greatest Common Divisor (GCD); 
Least Common Multiple (LCM); scalars; floor or ceiling operators; exponents (excluding square and cube).

\medskip
\noindent\textbf{Permitted objects.}
Elements of $\Z$; finite sets, sequences, and directed graphs with
integer-labelled vertices; finite compositions of the atomic operators defined
in Section~\ref{sec:defs}; finite integer comparisons.

\medskip
\noindent\textbf{Separation Principle.}
Every document in the RS corpus maintains a strict separation between
integer-arithmetic derivations (presented in RS Ontological Statement
boxes) and external verification using real-number or continuous-mathematics
tools (presented in External Verification boxes).  Content in
External Verification boxes carries no ontological weight and is never
reintroduced into RS evolution.

\medskip
\noindent\textbf{Verification lens.}
Real-number computations (eigenvalues, decimal approximations of constants)
are admissible solely as external verification of integer dynamics.  They are
presented in External Verification boxes and never inserted back into
the evolution.

\newpage


\section{Symbol Table}\label{sec:symboltable}

\subsection{Integer and Set Primitives}

\begin{longtable}{L{2.0cm} L{2.6cm} L{8.5cm}}
\toprule
\textbf{Symbol} & \textbf{Command} & \textbf{Meaning} \\
\midrule\endhead\bottomrule\endfoot

$\Z$ & \verb|\Z| & Set of all integers; the sole ontological substrate \\[3pt]
$\Z_{\geq 0}$ & & Non-negative integers \\[3pt]
$\Z_{>0}$ & & Positive integers \\[3pt]
$n$ & & Numerator component of an ERP; magnitude or count \\[3pt]
$d$ & & Denominator component of an ERP; geometric capacity \\[3pt]
$(n,d)$ & & Explicit Rational Pair; the fundamental state object \\[3pt]
$\tau$ & \verb|\tau| & Microtick index; element of $\Tcyc$ \\[3pt]
$t$ & & Causal time (integer step counter) \\[3pt]
$k$ & & Iteration index or rank value (context-dependent) \\[3pt]
$p,q$ & & Generic ERP components in proof contexts \\[3pt]
$m$ & & Generic integer argument to rank/height functions \\[3pt]
\end{longtable}

\subsection{State Domains}

\begin{longtable}{L{2.0cm} L{2.6cm} L{8.5cm}}
\toprule
\textbf{Symbol} & \textbf{Command} & \textbf{Meaning} \\
\midrule\endhead\bottomrule\endfoot

$\SL$ & \verb|\SL| & Linear ERP state space:
  $\{(n,d)\in\Z^2:d\neq 0\}$; primary vacuum domain \\[3pt]
$\SLe$ & \verb|\SLe| & Extended linear domain:
  $\SL\cup\{(n,0):n\in\Z\}$; includes suspension states \\[3pt]
  $\SLpos$ & \verb|\SLpos| & Positive-denominator sublattice:
  $\{(n,d)\in\Z^2:d\geq 1\}$; canonical representative of each
  observational equivalence class in $\SL$; domain of all
  velocity ERPs and ensemble states \\[3pt]
$\SPs$ & \verb|\SPs| & Active matter domain:
  $\{(X,Y,Z)\in\Z^3:X\neq 0,\,Y\neq 0,\,Z\neq 0\}$ \\[3pt]
$\SP$ & \verb|\SP| & Full projective domain: $\SPs\cup\IP$ \\[3pt]
$\SPe$ & \verb|\SPe| & Extended projective domain \\[3pt]
$\IP$ & \verb|\IP| & Projective identity element: $\{(0,1,0)\}$ \\[3pt]
$V$ & & Unit entry state: $(0,1)\in\SL$; generative seed of the $\aop$-orbit;
  not a physical state.  The zero numerator means no magnitude is present;
  serves as the structural origin from which the Lucas and Fibonacci sequences
  are generated, not as an active ERP in the ontological chain \\[3pt]
$\Vac$ & \verb|\Vac| & Vacuum line: $\{(p,1):p\in\Z\}\subset\SL$ \\[3pt]
$\Tcyc$ & \verb|\Tcyc| & Prime-11 time cycle: $\{1,2,\dots,11\}$;
  also $\mathcal{T}_{11}$ \\[3pt]
$\PhiSet$ & \verb|\PhiSet| & Triadic phase set: $\{E,M,R\}$ \\[3pt]
$\Phi_N$ & & Phase space modulus-$N$:
  $\{0,\dots,N-1\}$ with successor \\[3pt]
$E$ & & Spacetime event:
  $((n_t,d_t),(n_x,d_x))\in\SL\times\SL$ \\[3pt]
$\upsilon$ & \verb|\upsilon| & Upper oscillator:
  first linear ERP in a coupled vacuum pair \\[3pt]
$\beta_{\mathrm{osc}}$ & & Lower oscillator:
  second linear ERP in a coupled vacuum pair \\[3pt]
\end{longtable}

\subsection{Atomic Operators}

\begin{longtable}{L{2.0cm} L{2.6cm} L{8.5cm}}
\toprule
\textbf{Symbol} & \textbf{Command} & \textbf{Meaning / Action} \\
\midrule\endhead\bottomrule\endfoot

$\lop$ & \verb|\lop| & \textbf{Linear Transform:}
  $\lop(n,d):=(n+d,d)$; vacuum propagation;
  matrix $\mathbf{L}=\bigl(\begin{smallmatrix}1&1\\0&1\end{smallmatrix}\bigr)$
  \\[3pt]

$\aop$ & \verb|\aop| & \textbf{Aggregate Transform:}
  $\aop(n,d):=(n+d,n)$;
  matrix $\mathbf{H}=\bigl(\begin{smallmatrix}1&1\\1&0\end{smallmatrix}\bigr)$;
  Fibonacci generator; Lucas generator from seed $(T_{\mathrm{vac}},\Sigimb)$
  \\[3pt]

$\psi$ & \verb|\psi| & \textbf{Transformative Reciprocal (coupled):}
  $\psi\!\left(\tfrac{a}{b},\tfrac{c}{d}\right):=\left(\tfrac{d}{a},\tfrac{b}{c}\right)$;
  denominators cross to become numerators of the other state;
  numerators become new denominators;
  coupled rational pairs. They cannot act on one state in isolation;
  order~4; $\psi^4=\mathrm{Id}$; $\psi^2$ is charge conjugation;
  suspension resolution, spin, and chirality are structural consequences \\[3pt]


$\bplus$ & \verb|\bplus| & \textbf{Barycentric/Matter Addition:}
  $(n_1,d_1)\bplus(n_2,d_2):=(n_1 d_2+n_2 d_1,d_1 d_2)$ \\[3pt]

$\med$ & \verb|\med| & \textbf{Mediant Sum:}
  $(a,b)\med(c,d):=(a+c,b+d)$; width-preserving \\[3pt]

$\succ$ & \verb|\succ| & \textbf{Successor with Reset:}
  $\succ(\tau)=\tau+1$ if $\tau<11$; $\succ(11)=1$ \\[3pt]

$G$ & & \textbf{Projective Matter Generator:}
  $G:\SP\to\SP$; polynomial map; three cases (Section~\ref{sec:mattergenerator})
  \\[3pt]

$\koppa$ & \verb|\koppa| & \textbf{Imbalance Operator/Ledger:}
  $x\koppa L$ returns the constructive remainder of $x$ relative to $L$
  via iterative subtraction (no division); as ledger: accumulated structural
  history at $\Omega$; rest mass $m_{\mathrm{rest}}\equiv\rank(\koppa_{\mathrm{ledger}})$
  \\[3pt]

$L(\upsilon,\beta)$ & & \textbf{Cross-Synchronisation Lift:}
  $(X,Y,Z):=(n_\upsilon d_\beta,\,n_\beta d_\upsilon,\,d_\upsilon d_\beta)$;
  embeds $\SL\times\SL\to\SPe$ \\[3pt]

$\PG(k)$ & \verb|\PG| & \textbf{Prime Gate:}
  the $k$-th prime in natural order; admissible RS primitive
  (Section~\ref{sec:primegate}) \\[3pt]
\end{longtable}

\subsection{Measurement and Structural Functions}

\begin{longtable}{L{2.2cm} L{2.6cm} L{8.5cm}}
\toprule
\textbf{Symbol} & \textbf{Command} & \textbf{Meaning} \\
\midrule\endhead\bottomrule\endfoot

$\rank(m)$ & \verb|\rank(m)| & \textbf{Structural Rank:}
  unique $k\geq 0$ s.t.\ $G_{\mathrm{mg}}^k\leq m<G_{\mathrm{mg}}^{k+1}$
  ($m>0$); computed by counting $G_{\mathrm{mg}}$-doublings (integer
  multiplication and comparison only); $\rank(0):=0$ \\[3pt]
  
  $\rho_\Omega(m)$ & & \textbf{Prime-Event Depth:}
  total prime factors of $m$ with multiplicity; computed via
  $\koppa$-primality and iterative subtraction; $\rho_\Omega(0):=\rho_\Omega(1):=0$;
  $\rho_\Omega(\mathcal{V}_S)=N$; $\rho_\Omega(\Lambda_\mu)=N$ \\[3pt]

$\hfn(m)$ & \verb|\hfn(m)| & \textbf{Height Function:}
  unique $k\geq 0$ s.t.\ $2^{2k}\leq m^2<2^{2k+2}$ ($m\neq 0$) \\[3pt]

$\Dcross(a,b)$ & \verb|\Dcross| & \textbf{Cross-Determinant:}
  $\Dcross(a,b):=p_b q_a-p_a q_b$; skew-symmetric \\[3pt]

$W(L,U)$ & & \textbf{Interval Width:}
  $W(L,U):=\Dcross(L,U)$ for $L\prec U$; saturation at $W=1$ \\[3pt]

$\sgn(m)$ & \verb|\sgn| & \textbf{Sign Function:}
  $\{-1,0,+1\}$; integer-only \\[3pt]

$\rem(a,n)$ & \verb|\rem| & \textbf{Constructive Remainder:}
  iterative subtraction; replaces modulo \\[3pt]

$\hat{h}(P)$ & \verb|\hat{h}(P)| & \textbf{Unreduced Height on Elliptic Curve:}
  $\max(\rank(n),\rank(d))$ for coordinate $n/d$ \\[3pt]

$\VS$ & \verb|\VS| & \textbf{Viscosity Sum:}
  accumulated phase drift over a closed causal orbit;
  $\VS = 18$ for the standard vacuum orbit \\[3pt]

$\Phi(t)$ & \verb|\Phi(t)| & \textbf{Structural Friction:}
  cumulative rank accumulation over $t$ steps \\[3pt]

$\LQ(N)$ & \verb|\LQ(N)| & \textbf{Discrete L-Function:}
  $(R_N,T_N)\in\SL$; orientation-bias summary \\[3pt]

$H(s)$ & & \textbf{Structural Entropy:}
  $H(s):=\rank(d)$ for $s=(n,d)$; causal depth of the metric basis \\[3pt]
  
  $\SRS(\Ens)$ & \verb|\SRS| & RS Thermodynamic Entropy:
  $\sum_{j}\rank(d_j)$ over an ensemble $\Ens\subset\SLpos$ \\[3pt]
$\TRS(\Ens)$ & \verb|\TRS| & RS Thermal Charge:
  $\sum_{j}(n_j\koppa d_j)$ over an ensemble $\Ens\subset\SLpos$ \\[3pt]
$\GK(U)$ & \verb|\GK| & RS Kinetic Charge:
  $n_u^2$; complement of $\Tcharge$ in the partition
  $\Tcharge(U)+\GK(U)=d_u^2$ \\[3pt]
$\Ens$ & \verb|\Ens| & RS Ensemble: finite ordered collection
  $\{s_1,\ldots,s_M\}\subset\SLpos$ \\[3pt]

$L(n)$ & & \textbf{Lucas Number:}
  $L(0)=2$, $L(1)=1$, $L(n)=L(n-1)+L(n-2)$;
  generated by $\aop$ from seed $(T_{\mathrm{vac}},\Sigimb)$ \\[3pt]

$\#(p)$ & & \textbf{Primorial:}
  product of all primes $\leq p$;
  $\#(T_{\mathrm{vac}})=2\cdot3\cdot5\cdot7\cdot11=2310$ \\[3pt]
\end{longtable}

\subsection{Equivalence and Ordering Relations}

\begin{longtable}{L{2.2cm} L{2.6cm} L{8.5cm}}
\toprule
\textbf{Symbol} & \textbf{Command} & \textbf{Meaning} \\
\midrule\endhead\bottomrule\endfoot

$a\sim b$ & \verb|a\sim b| & Observational equivalence:
  $n_1 d_2=n_2 d_1$; GCD reduction forbidden \\[3pt]
$a\prec b$ & \verb|a\prec b| & Cross-ordering: $\Dcross(a,b)>0$ \\[3pt]
$\operatorname{Adj}(a,b)$ & & Adjacency: $\Dcross(a,b)=\pm 1$ \\[3pt]
\end{longtable}

\subsection{Temporal and Phase Symbols}

\begin{longtable}{L{2.2cm} L{2.6cm} L{8.5cm}}
\toprule
\textbf{Symbol} & \textbf{Command} & \textbf{Meaning} \\
\midrule\endhead\bottomrule\endfoot

$\Tcyc$ & \verb|\Tcyc| & Prime-11 time cycle $\{1,\dots,11\}$
  as finite directed graph \\[3pt]
$\phi(t)$ & & Phase evolution:
  $\phi(0):=1$; $\phi(t+1):=\succ(\phi(t))$ \\[3pt]
$\Omega$ & & Cycle boundary event: the reset $\succ(11)=1$ \\[3pt]
$\Omega_{\mathrm{vac}}$ & & Vacuum Resolution Frequency:
  integer count of vacuum generator applications \\[3pt]
$N_E,N_M,N_R$ & & Phase cardinalities:
  $N_E=4$, $N_M=4$, $N_R=3$ (4-4-3 partition) \\[3pt]
$\Tbary$ & \verb|\Tbary| & Barycentric period: $11\times 3=33$ ticks \\[3pt]
$\tau_{10}$ & & Inertial closure tick:
  $\tau_{10}:=T_{\mathrm{vac}}-1=10$; the unique microtick at which
  the temporal Arrow becomes irreversible within the cycle
  (Theorem~\ref{thm:tau10}) \\[3pt]
$f:\Tcyc\to\PhiSet$ & & Sequential phase mapping;
  the pointer function assigning each tick a phase \\[3pt]
$B=(L_E,L_M,L_R)$ & & Barycentric structure:
  cyclic set of ERPs at triadic stabilisation \\[3pt]
\end{longtable}

\subsection{Physical State Vector and Temporal Charge}

\begin{longtable}{L{2.5cm} L{2.6cm} L{8.5cm}}
\toprule
\textbf{Symbol} & \textbf{Command} & \textbf{Meaning} \\
\midrule\endhead\bottomrule\endfoot

$\Psistate(s)=(M,H,\Theta)$ & \verb|\Psistate| &
  Topological state vector;
  projection $\Pi:\SL\to\mathcal{V}_{\mathrm{phys}}$ \\[3pt]
$M(s)$ & & Magnitude component: $M(s):=n$ \\[3pt]
$H(s)$ & & History component: $H(s):=\hfn(d)$ \\[3pt]
$\Theta(s)$ & & Orientation component: $\Theta(s):=\tau\in\Tcyc$ \\[3pt]
$\Tcharge(n,d)$ & \verb|\Tcharge| & Temporal Charge:
  $\Tcharge(n,d):=d^2-n^2=\Dcross((n,d),(d,n))$ \\[3pt]
$\Gamma_{\min}=3$ & & Minimum temporal charge:
  at mass-gap state $(1,2)$ \\[3pt]
\end{longtable}

\subsection{Rational Special Relativity Symbols}

\begin{longtable}{L{2.5cm} L{2.6cm} L{8.5cm}}
\toprule
\textbf{Symbol} & \textbf{Command} & \textbf{Meaning} \\
\midrule\endhead\bottomrule\endfoot

$U=(n_u,d_u)$ & & Boost parameter ERP; velocity $u=n_u/d_u$ \\[3pt]
$P_i$ & & RS 3-Momentum: $m_{\mathrm{rest}}\cdot n_i$ for
  $i=1,2,3$; integer product of rest mass and velocity numerator \\[3pt]
$\Ethree$ & \verb|\Ethree| & RS 3+1D Spacetime Event:
  $((n_t,d_t),(X,Y,Z))\in\SLpos\times\SPs$ \\[3pt]
$\alpha_{\mathrm{boost}}$ & & Hyperbolic cosine term: $d_u^2+n_u^2$ \\[3pt]
$\beta_{\mathrm{boost}}$ & & Hyperbolic sine term: $2n_u d_u$ \\[3pt]
$\Delta_{\mathrm{boost}}$ & & Metric denominator: $d_u^2-n_u^2$;
  vanishes at luminal boundary $n_u=d_u$ \\[3pt]
$A_Q$ & & Rational Lorentz boost matrix:
  $A_Q=\bigl(\begin{smallmatrix}\alpha_{\mathrm{boost}}&\beta_{\mathrm{boost}}
  \\\beta_{\mathrm{boost}}&\alpha_{\mathrm{boost}}\end{smallmatrix}\bigr)$
  \\[3pt]
$I(E)$ & & Spacetime interval ERP:
  $(n_t^2 d_x^2-n_x^2 d_t^2,\,d_t^2 d_x^2)$ \\[3pt]
$U_{\mathrm{new}}$ & & Composed boost:
  $(n_1 d_2+n_2 d_1,\,d_1 d_2+n_1 n_2)$ \\[3pt]
$\Lambda_U$ & & Boost transformation on events \\[3pt]
\end{longtable}

\subsection{Structural Constants and Derived Integers}

\begin{longtable}{L{2.8cm} L{2.3cm} L{8.0cm}}
\toprule
\textbf{Symbol} & \textbf{Value} & \textbf{Meaning / Derivation} \\
\midrule\endhead\bottomrule\endfoot

$T_{\mathrm{vac}}$ & $11$ & Vacuum period; prime \\[3pt]
$N$ & $3$ & Triadic interaction logic period \\[3pt]
$\delta_{\mathrm{slip}}$ & $2$ & Phase slip:
  net $+2$ states after 11 ticks on 3-cycle \\[3pt]
$\Delta$ & $5$ & Flux imbalance: $(N_E+N_M)-N_R=5$ \\[3pt]
$\varepsilon$ & $2/3$ & Phase Slip Efficiency:
  $\varepsilon:=\delta_{\mathrm{slip}}/N=2/3$ \\[3pt]
$\Dg$ & $12$ & Gauge manifold dimension: $1+3+8=12$;
  also $\Delta+\Sigimb$ and $L(2)\times L(3)$ \\[3pt]
$\Sint$ & $132$ & Interaction surface:
  $11\times 12=132$ \\[3pt]
$\Tbary$ & $33$ & Barycentric period: $11\times 3=33$ \\[3pt]
$\aSC$ & $137$ & Sommerfeld Constant:
  $\Sint+\Delta=132+5=137=\PG(\Tbary)$ \\[3pt]
$\delta_A=\delta_B$ & $2/55$ & Dual fractional anomaly;
  $\delta_A=\delta_{\mathrm{slip}}/(\Delta\cdot T_{\mathrm{vac}})
  =(N^3-\Delta^2)/(\Delta\cdot T_{\mathrm{vac}})$ \\[3pt]
$\alpha^{-1}_{RS}$ & $137+\tfrac{2}{55}$ &
  RS geometric bare inverse fine-structure constant \\[3pt]
$K_{\mathrm{sat}}$ & $25$ & Saturation limit: $\Delta^2=25$ \\[3pt]
$S_{\mathrm{vac}}$ & $(1,11)$ & Vacuum seed ERP \\[3pt]
$S_{\mathrm{gap}}$ & $(2,5)$ & Imbalance seed ERP:
  mass gap on flux $\Delta=5$ \\[3pt]
$G_{\mathrm{mg}}$ & $2$ & Mass gap integer:
  minimal non-zero capacity rank; $=\PG(1)$ \\[3pt]
$\Lambda_\mu$ & $207$ & Secondary stability integer:
  $\Sint+\mathcal{R}_D+T_{\mathrm{vac}}=207$ \\[3pt]
$\Lambda_\tau$ & $3481$ & Tau mass integer:
  $\PG(\Dg+\Delta)^2=59^2$ \\[3pt]
$Q$ & $2/3$ & Koide ratio; equals $\varepsilon$ \\[3pt]
$\Bridge$ & $1303/22$ & Generation bridge: $\delta_\alpha/\delta_Q$ \\[3pt]
$\mu$ & $1836+\tfrac{8}{55}$ & Proton-electron mass ratio:
  $\Dg^3+(\Sint-2\Dg)+4\delta_A$ \\[3pt]
$m_n-m_p$ & $(\tfrac{5}{2}+\tfrac{2}{55})m_e$ &
  Neutron-proton mass difference \\[3pt]
$\Lalpha$ & ERP & Lattice coupling constant:
  $(\mathrm{count}(\Omega_\psi),\mathrm{count}(\Omega_{33}))$ \\[3pt]
$N_{\mathrm{crit}}$ & $>2^{128}$ & Critical phase space modulus;
  proton stability bound \\[3pt]
$\Sigimb$ & $7$ & Boundary-Active Phase Sum:
  $N_E+N_R=4+3=7$; also $\Delta+\delta_{\mathrm{slip}}$ \\[3pt]
$\tau_{10}$ & $10$ & Inertial closure tick:
  $T_{\mathrm{vac}}-1$; Arrow irreversibly fixed \\[3pt]
$45$ & $\Delta\cdot N^2$ & PMNS suppressor:
  $\Tbary+\Dg=33+12$ \\[3pt]
$13$ & $\PG(\Sigimb)$ & 4th Fibonacci prime;
  $T_{\mathrm{vac}}+G_{\mathrm{mg}}$; $F_7$ \\[3pt]
$29$ & $\PG(\tau_{10})$ & $\PG(10)=\Tbary-N_E$;
  phi-oscillator descent prime \\[3pt]
$\Lambda_{q1}$ & $207$ & Gen-1 quark mass integer:
  $\Sint+64+T_{\mathrm{vac}}$ (shared with $\Lambda_\mu$) \\[3pt]
$\Lambda_{q2}$ & $303$ & Gen-2 quark mass integer:
  $\Sint+160+T_{\mathrm{vac}}$ \\[3pt]
$\Lambda_{q3}$ & $192$ & Gen-3 quark mass integer:
  phase-locked reference; $192\koppa 3=0$ \\[3pt]
$4621$ & & $\phi$-oscillator denominator:
  $2\times\#(T_{\mathrm{vac}})+1$; minimal uniform structural tension \\[3pt]
$7477$ & & $\phi$-oscillator numerator:
  prime; $\rank=\Dg$; $(7477+N)\koppa T_{\mathrm{vac}}=0$ \\[3pt]
$1309$ & $\Lambda_{\mathrm{thresh}}$ & Threshold sphenic integer:
  $\Sigimb\times T_{\mathrm{vac}}\times(\Dg+\Delta)=7\times11\times17$ \\[3pt]
$\pa$ & $115331$ & RS Sommerfeld prime: numerator of the RS Sommerfeld ERP;
  generated by the CF convergent recurrence from $\Da$;
  $\pa\koppa\Tbary=29=\PG(\tauX)$; $\pa\koppa\Tvac=7=\Sigimb$ \\[3pt]
$\qa$ & $15804499$ & RS Sommerfeld prime: denominator of the RS Sommerfeld ERP;
  generated by the CF convergent recurrence from $\Da$;
  $\qa\koppa\Tvac=7=\Sigimb$; $\qa\koppa\Tbary=7=\Sigimb$ \\[3pt]
$\piRS$ & $(355,113)$ & RS $\pi$-Analog ERP: structural representative of the
  $W=1$ bracket reached by four mediant-tree legs consuming $N$, $\Sigimb$,
  $N\times\Delta$, and the vacuum unit; total koppa load
  $\Gmg\times\Pithree=26$ \\[3pt]
\end{longtable}

\subsection{Particle Classification Symbols}

\begin{longtable}{L{2.5cm} L{2.3cm} L{8.3cm}}
\toprule
\textbf{Symbol} & \textbf{Value} & \textbf{Meaning} \\
\midrule\endhead\bottomrule\endfoot

$\mathcal{R}_D$ & $C=64$ & Dyadic closure class; quark sector \\[3pt]
$\mathcal{R}_T$ & $C=72$ & Triadic closure class; Electron \\[3pt]
$\mathcal{R}_E$ & $C=128$ & Extended dyadic class; Tau \\[3pt]
$C$ & & Geometric capacity at stabilisation; equals cycle discriminant
  $\mathrm{disc}(M)$ for minimal-orbit cycles;
  equals observable mass $m_{\mathrm{obs}}=C$ \\[3pt]
$A_\mu=64$ & & Effective capacity for muon anomalous magnetic moment \\[3pt]
$A_\tau=128$ & & Effective capacity for tau anomalous magnetic moment \\[3pt]
\end{longtable}

\subsection{Matrix and Group Symbols}

\begin{longtable}{L{2.5cm} L{2.6cm} L{8.5cm}}
\toprule
\textbf{Symbol} & \textbf{Command} & \textbf{Meaning} \\
\midrule\endhead\bottomrule\endfoot

$\mathbf{L}=\bigl(\begin{smallmatrix}1&1\\0&1\end{smallmatrix}\bigr)$ & &
  Matrix of $\lop$ \\[3pt]
$\mathbf{H}=\bigl(\begin{smallmatrix}1&1\\1&0\end{smallmatrix}\bigr)$ & &
  Matrix of $\aop$ \\[3pt]
$\mathbf{S}=\bigl(\begin{smallmatrix}0&1\\1&0\end{smallmatrix}\bigr)$ & &
  Sign-change matrix; external verification path only;
  completes $GL_2(\Z)$ generator set for cycle discriminant computation \\[3pt]
$M\in GL_2(\Z)$ & & Cycle matrix:
  finite product of $\mathbf{L}$, $\mathbf{H}$, $\mathbf{S}$;
  external verification path only \\[3pt]
$\kappa_1,\kappa_2$ & & Eigenvalues of $M$; external verification only \\[3pt]
$GL_2(\Z)$ & & Group generated by $\mathbf{L}$, $\mathbf{H}$, $\mathbf{S}$;
  external verification path only \\[3pt]
\end{longtable}

\subsection{Mediant, Geodesic, Elliptic Curve, and L-Function Symbols}

\begin{longtable}{L{2.5cm} L{2.6cm} L{8.5cm}}
\toprule
\textbf{Symbol} & \textbf{Command} & \textbf{Meaning} \\
\midrule\endhead\bottomrule\endfoot

$\mathcal{T}$ & & Mediant Tree: binary tree of integer pairs \\[3pt]
$\gamma(A,B)$ & & Mediant geodesic from $A$ to $B$ \\[3pt]
$\{A,B\}$ & & Structural symbol:
  formal sum of geodesic edges \\[3pt]
$P\ecplus Q$ & & Projective elliptic group law on $\SP$;
  polynomial; no division \\[3pt]
$G_{\mathrm{ec}}\in\SPs$ & & Generator point of infinite order
  on elliptic curve \\[3pt]
$\{S_t\}$ & & Evolution sequence: $S_t=t\cdot G_{\mathrm{ec}}$ \\[3pt]
$T_N$ & & Orbit period in $\LQ(N)$ \\[3pt]
$R_N$ & & Orientation bias sum:
  $\sum_t\sgn(\Dcross(S_t,S_{t+1}))$ \\[3pt]
$\tau_t$ & & Cross-determinant at step $t$:
  $\Dcross(S_t,S_{t+1})$ \\[3pt]
$\sigma(\tau_t)$ & & Orientation sign: $\sgn(\tau_t)\in\{-1,0,+1\}$ \\[3pt]
\end{longtable}

\newpage


\section{Core Axioms}\label{sec:axioms}

\begin{axiom}[Integer-only substrate]\label{ax:integer-substrate}
All admissible objects are defined over $\Z$ using finite comparisons and
finite compositions.  No real-number limits are used.
\end{axiom}

\begin{axiom}[No division]\label{ax:no-division}
Division is not an admissible primitive.  Any expression resembling a
rational number must be represented as an ERP $(n,d)$ without GCD reduction.
\end{axiom}

\begin{axiom}[No logarithm; no absolute value as primitive]\label{ax:no-log}
The absolute value $|m|$ implies a square-root construction and is not a
primitive.  Magnitude is expressed via $\rank(m)$, $\hfn(m)$, or $m^2$.
\end{axiom}

\begin{axiom}[Observational non-reduction]\label{ax:non-reduction}
GCD reduction is forbidden as an ontological rule.  Two states $(n_1,d_1)$
and $(n_2,d_2)$ are observationally equivalent iff $n_1 d_2=n_2 d_1$ and
are nonetheless distinct states with distinct structural histories.
\end{axiom}

\begin{axiom}[ZERO-logic suspension]\label{ax:zero-logic}
A Constraint Vacuum state $n/0$ with $n\neq 0$ is a suspension state, not
a division error.  It carries structural history $n$ in the numerator with
no metric basis.  It is resolved by $\psi$ acting on it paired with an
active state, distributing history to the kinematic ledger and capacity
to the structural entropy ledger at $\Omega$ (Definition~\ref{def:zero-resolution}).
\end{axiom}

\begin{axiom}[Forbidden continuum]\label{ax:forbidden-continuum}
$\mathbb{R}$, continuous manifolds, limits to infinity, and infinitesimal
operators are prohibited.  Existence is restricted to integers, finite
sets/sequences of integers, and finite directed graphs with integer-labelled
vertices.
\end{axiom}

\begin{axiom}[Mask Postulate]\label{ax:mask}
Continuous constants ($\pi$, $e$, $\phi$, $\sqrt{n}$, \ldots) are not
ontological primitives.  They are phenomenological masks produced by
high-frequency rational oscillations between ERPs; the barycentric average of
a convergent integer orbit.  All such constants appear exclusively in
External Verification boxes and carry no ontological weight.
\end{axiom}

\begin{axiom}[Triadic Incommensurability]\label{ax:triadic}
The interaction logic period is $N=3$; the temporal geometry period is
$T_{\mathrm{vac}}=11$.  After 11 successor operations on a 3-state cycle,
the net phase displacement is $+2$ states ($\delta_{\mathrm{slip}}=2$).
This residual is the engine of the Arrow of Time and of all non-equilibrium
evolution.
\end{axiom}

\begin{axiom}[Mass-Gap Postulate]\label{ax:mass-gap}
Vacuum phase: evolution in $\SL$ via $\{\lop,\aop,\psi\}$; rank growth
additive.  Matter phase: evolution in $\SP$ via $\bplus$; rank growth
multiplicative (cubic for the doubling generator).  No admissible state
satisfies $0<m_L(s)<1$.
\end{axiom}

\begin{axiom}[Arrow of Time]\label{ax:arrow}
Any physically valid matter generator must enforce strict structural potential
monotonicity: $\sum\hfn(P_{t+1})>\sum\hfn(P_t)$ for every valid matter
transition.
\end{axiom}

\begin{axiom}[Invariance of the Interval]\label{ax:interval}
The interval $I(E)=(n_t^2 d_x^2-n_x^2 d_t^2,\,d_t^2 d_x^2)$ is invariant
under the rational Lorentz boost $\Lambda_U$; proved by the integer identity
$\alpha_{\mathrm{boost}}^2-\beta_{\mathrm{boost}}^2=\Delta_{\mathrm{boost}}^2$.
\end{axiom}

\subsection{The Prime Gate: Admissible Primitive}\label{sec:primegate}

\begin{definition}[Prime Gate $\PG$]\label{def:primegate}
\begin{rsbox}
For any positive integer $k$, the Prime Gate $\PG(k)$ is the $k$-th
prime in natural order: $\PG(1)=2$, $\PG(2)=3$, $\PG(3)=5$, $\PG(4)=7$,
$\ldots$

Primality of a candidate integer $p>1$ is detected by $\koppa$ alone: $p$ is
prime if and only if $p\koppa m\neq 0$ for all integers $m$ with $1<m<p$.
Finding $\PG(k)$ requires initialising a counter at zero, starting at $p=2$,
applying the $\koppa$-primality test, and advancing the counter each time a
prime is found until the counter equals $k$.  This terminates for any finite
$k$ using only $\koppa$, comparison, and $\succ$.
\end{rsbox}
\end{definition}

\begin{theorem}[Ontological Admissibility of $\PG$]\label{thm:pg-admissible}
$\PG(k)$ is an admissible RS primitive.  It is a finite integer determination
using exclusively RS-native operations.
\begin{proof}
Primality testing requires only $\koppa$ (iterative subtraction), integer
comparison, and $\succ$.  None of these involve real numbers, limits, or any
operation prohibited by Axiom~\ref{ax:forbidden-continuum}.  The determination
terminates at $k$ prime-finds, each of which is a finite integer bound.
\end{proof}
\end{theorem}

\begin{postulate}[Frictionless Propagation Principle]\label{post:frictionless}
\begin{rsbox}
A prime-capacity ERP $(n,p)$ is structurally irreducible.  Because
$p$ has no proper divisors, the state cannot bleed partial structural load
into any sub-period oscillator: no sub-capacity $m$ with $1<m<p$ divides $p$,
so no intermediate resonance is available.  The state must engage the full
vacuum architecture across $\Sint$ or not at all.

This all-or-nothing vacuum engagement is the structural mechanism by which
structure precipitates at prime-capacity states.  A composite-capacity state
$(n,ab)$ can shed structural load into the $a$-period and $b$-period
oscillators separately; a prime-capacity state cannot.

The same principle governs $T_{\mathrm{vac}}=11$: the vacuum period is prime
so that the 11-tick cycle cannot decompose into sub-periods that allow the
triadic pointer to resynchronise prematurely.  The Frictionless Propagation
Principle is the general statement of this structural fact.
\end{rsbox}
\end{postulate}

\newpage


\section{Definitions: Substrate and Operators}\label{sec:defs}

\begin{definition}[Explicit Rational Pair (ERP)]\label{def:ERP}
An ordered pair $(n,d)\in\Z\times(\Z\setminus\{0\})$.  The numerator $n$
encodes event magnitude; the denominator $d$ encodes geometric capacity.
ERPs are never reduced to lowest terms.
\end{definition}

\begin{definition}[Observational equivalence]\label{def:obs-equiv}
$(n_1,d_1)\sim(n_2,d_2)$ iff $n_1 d_2=n_2 d_1$.  The sole admissible
equivalence; GCD reduction is not a rule of the ontology.
\end{definition}

\begin{definition}[Cross-determinant deviation $\Dcross$]\label{def:Dcross}
For $a=(p_a,q_a)$, $b=(p_b,q_b)\in\SL$:
\[
\Dcross(a,b):=p_b q_a-p_a q_b\in\Z.
\]
Vanishes iff $a\sim b$.  Skew-symmetric: $\Dcross(a,b)=-\Dcross(b,a)$.
\end{definition}

\begin{definition}[Interval width and adjacency]\label{def:width}
For $L\prec U$: $W(L,U):=\Dcross(L,U)\in\Z_{>0}$.
Adjacency (saturation): $W(L,U)=1$.
\end{definition}

\begin{definition}[Mediant sum]\label{def:mediant}
$(a,b)\med(c,d):=(a+c,b+d)$.  Property:
$\Dcross(L,L\med U)=\Dcross(L\med U,U)=\Dcross(L,U)$ (width preserved).
\end{definition}

\begin{definition}[Discrete Convergence]\label{def:discrete-convergence}
A sequence $\{x_n\}\subset\SL$ is said to converge to a
Barycentric Average if it enters a stable Discrete Bracket
$[L,U]$ of width $W(L,U)=1$ in finite iterations.  Once $W=1$ is achieved,
the orbit is Barycentrically Truncated; the ``value'' of the
attractor is the oscillation between the adjacent integer states $L$ and $U$,
not a point on a real-number line.
\end{definition}

\begin{definition}[Barycentric Truncation Operator $\mathcal{T}_B$]
\label{def:bary-trunc}
For a system in a state of Discrete Convergence within a bracket $[L,U]$
where $W(L,U)=1$, the truncation operator $\mathcal{T}_B$ selects the
representative state $S_{\mathrm{rep}}$ based on the parity of the structural
history.  Let
$S_{\mathrm{rep}}=\mathcal{T}_B([L,U],\koppa_{\mathrm{ledger}})$ be defined
by the selection rule:
\[
S_{\mathrm{rep}}:=\begin{cases}
  L & \text{if }\rank(\koppa_{\mathrm{ledger}})\text{ is even}\\
  U & \text{if }\rank(\koppa_{\mathrm{ledger}})\text{ is odd}
\end{cases}
\]
Here $\koppa_{\mathrm{ledger}}$ denotes the accumulated Imbalance Ledger
value at the cycle boundary $\Omega$. The integer $r$ accumulated via the
iterative subtraction procedure of Definition~\ref{def:koppa} over the
current orbit (see the operator/ledger distinction therein).  This selection
resolves the oscillation between adjacent integer states into a single ERP
state for matter generation $G$ while preserving the underlying parity of the
imbalance ledger.
\end{definition}

\begin{definition}[Linear Transform $\lop$]\label{def:lop}
$\lop(n,d):=(n+d,d)$.
Matrix: $\mathbf{L}=\bigl(\begin{smallmatrix}1&1\\0&1\end{smallmatrix}\bigr)$.
Increments magnitude; preserves denominator.
\end{definition}

\begin{definition}[Aggregate Transform $\aop$]\label{def:aop}
$\aop(n,d):=(n+d,n)$.
Matrix: $\mathbf{H}=\bigl(\begin{smallmatrix}1&1\\1&0\end{smallmatrix}\bigr)$.
The Fibonacci generator.  Applied to the vacuum seed
$(T_{\mathrm{vac}},\Sigimb)=(11,7)$, it generates the Lucas sequence
(Section~\ref{sec:lucas}).  When $n=0$, maps into $\SLe\setminus\SL$.
\end{definition}


\begin{definition}[Transformative Reciprocal $\psi$]\label{def:rop}
The \emph{Transformative Reciprocal} acts on a coupled pair of rational states.
It is the primary structural mechanism of RS\@.  Spin, chirality, charge
conjugation, and suspended-state resolution are direct consequences of its
cross-reciprocal form, not assigned properties.  The operator period
$\psi^4=\mathrm{Id}$ equals the Emission cardinality $N_E=4$
(Proposition~\ref{prop:psi-ne}).

\[
\psi\!\left(\frac{a}{b},\,\frac{c}{d}\right) := \left(\frac{d}{a},\,\frac{b}{c}\right)
\]

The denominator of each state crosses to become the numerator of the other.
The original numerators become the new denominators.
$\psi$ is a coupled operation and cannot act on one state in isolation.

The four iterates:
\begin{align*}
\psi^1\!\left(\tfrac{a}{b},\tfrac{c}{d}\right)
  &= \left(\tfrac{d}{a},\,\tfrac{b}{c}\right)
  & \text{(cross-reciprocal)} \\
\psi^2\!\left(\tfrac{a}{b},\tfrac{c}{d}\right)
  &= \left(\tfrac{c}{d},\,\tfrac{a}{b}\right)
  & \text{(full pair exchange; Charge Conjugation - Theorem~\ref{thm:cc})} \\
\psi^3\!\left(\tfrac{a}{b},\tfrac{c}{d}\right)
  &= \left(\tfrac{b}{c},\,\tfrac{d}{a}\right)
  & \text{(inverse cross-reciprocal)} \\
\psi^4\!\left(\tfrac{a}{b},\tfrac{c}{d}\right)
  &= \left(\tfrac{a}{b},\,\tfrac{c}{d}\right)
  & \text{(identity)}
\end{align*}

Physical roles: spin (period~4 equals $N_E$), chirality, resolution of
Constraint Vacuum states (Definition~\ref{def:zero-resolution}), and
charge conjugation ($\psi^2$).
\end{definition}

\begin{definition}[Barycentric addition $\bplus$ (linear form)]
\label{def:bplus}
$(n_1,d_1)\bplus(n_2,d_2):=(n_1 d_2+n_2 d_1,\,d_1 d_2)$.
Denominator becomes $d_1 d_2$; rank grows multiplicatively.  Engine of mass
generation in the linear sector.
\end{definition}

\begin{definition}[Imbalance Operator and Ledger $\koppa$]\label{def:koppa}
\begin{rsbox}
\textit{As operator (constructive remainder):}

For integers $x$ and positive integer $L>0$, initialise $r:=x$.  Repeatedly
subtract $L$ from $r$ while $r\geq L$, and repeatedly add $L$ to $r$ while
$r<0$.  When the process terminates, the resulting integer $r$ satisfies
$0\leq r<L$.  Define $x\koppa L:=r$.  This procedure uses only addition and
subtraction; no division or floor function appears.

\textit{As quantity (Imbalance Ledger):}

The notation $\koppa_{\mathrm{ledger}}$ denotes the accumulated structural
history unresolved within the current cycle, stored at the cycle boundary
$\Omega$.  The operator use $x\koppa L$ (with two explicit operands) and the
ledger use $\koppa_{\mathrm{ledger}}$ (a standalone accumulated integer) are
always distinguishable by context; the subscript is used where disambiguation
is required.

The rest mass is defined by $m_{\mathrm{rest}}\equiv\rank(\koppa_{\mathrm{ledger}})$.
\end{rsbox}
\end{definition}

\begin{definition}[Two-Ledger Structure at $\Omega$]\label{def:two-ledgers}
\begin{rsbox}
Two distinct accumulation ledgers operate across each vacuum cycle and are
read simultaneously at the cycle boundary $\Omega$.  They are not redundant:
a state may contribute to one and not the other.  Together they provide the
complete accounting of structural history at the cycle boundary.

\medskip
\noindent\textbf{1.\ Kinematic ledger ($\koppa_{\mathrm{ledger}}$):}
Accumulates $\koppa(n_j, d_j)$ at each step $j$ in the cycle.  Read at
$\Omega$: $m_{\mathrm{rest}} = \rank(\koppa_{\mathrm{ledger}})$.  Active
massive states ($\Tcharge > 0$) contribute non-zero values.  Luminal states
($\Tcharge = 0$, i.e.\ $n=d$) contribute zero ($\koppa(n,n)=0$).  States
of the form $0/c$ contribute zero ($\koppa(0,c)=0$ for all $c>0$).
Superluminal states ($\Tcharge<0$) contribute zero and are structurally
virtual (see Remark~\ref{rem:virtual-carriers}).

\medskip
\noindent\textbf{2.\ Structural entropy ledger:}
Accumulates $\rank(d_j)$ at each step $j$ in the cycle.  Read at $\Omega$
as causal depth.  States of the form $0/c$ contribute $\rank(c)\neq 0$:
capacity $c$ is deposited as causal depth, not kinematic imbalance.  This
is how suspended state resolution propagates structural history into 3D
expression without generating rest mass at the kinematic ledger.

\medskip
\noindent\textbf{Remark on two uses of koppa accumulation.}
The same ledger mechanism operates at two distinct structural scales.
(a)~\emph{Per-state rest mass:} the $\koppa_{\mathrm{ledger}}$ of a specific
matter state pair accumulated over one $\psi$-period gives
$m_{\mathrm{rest}} = \rank(\koppa_{\mathrm{ledger}})$ for that state.
(b)~\emph{Scale parameter for mixing channels:} cumulative
$\koppa$-accumulation from a specific seed evolved over a characteristic
number of steps gives the denominator scale at which a mixing channel
operates (Section~\ref{sec:dressing}).  These are the same mechanism
applied at different structural levels; they are not in conflict.
\end{rsbox}
\end{definition}

\begin{definition}[Cross-synchronisation lift $L(\upsilon,\beta)$]
\label{def:cross-sync}
For $\upsilon=(n_\upsilon,d_\upsilon)\in\SL$ and
$\beta=(n_\beta,d_\beta)\in\SL$:
\[
L(\upsilon,\beta):=(n_\upsilon d_\beta,\;n_\beta d_\upsilon,\;d_\upsilon d_\beta).
\]
If $n_\upsilon,n_\beta\neq 0$ and $d_\upsilon d_\beta\neq 0$, the output lies
in $\SPs$.
\end{definition}

\newpage


\section{Mass and Charge Taxonomy}\label{sec:mass-taxonomy}

\begin{definition}[Hierarchy of mass and charge]\label{def:mass-hierarchy}
\begin{rsbox}
\emph{Rest mass and inertial mass.}  A single origin defines both rest and
inertial mass.  For any state, the rest mass $m_{\mathrm{rest}}$ and inertial
mass $m_I$ are equal: $m_I\equiv m_{\mathrm{rest}}\equiv
\rank(\koppa_{\mathrm{ledger}})$.  There is no distinct notion of inertial
versus rest mass; the structural rank of the imbalance ledger uniquely
determines the mass.

\emph{Temporal charge.}  For a linear ERP $(n,d)$, the temporal charge is
$\Tcharge(n,d)=d^2-n^2$.  Positive temporal charge ($\Tcharge>0$)
characterises subluminal, massive states; $\Tcharge=0$ corresponds to luminal
states where $n=d$; negative values are excluded.

\emph{Mass gap.}  The mass gap $G_{\mathrm{mg}}=2=\PG(1)$ is the minimal
non-zero denominator $d$ such that $\Tcharge(n,d)$ can be positive.  At
denominator $d=1$ the system is luminal and carries no mass; the smallest
positive mass arises when $d=2$, opening the gate for the accumulation of
temporal charge.

\emph{Observable mass.}  When a periodic operator sequence stabilises a
particle state, its cycle matrix $M\in GL_2(\Z)$ has discriminant
$C:=\mathrm{tr}(M)^2-4\det(M)$.  This integer discriminant equals the
geometric capacity of the stabilised ERP and is the inertial/observable mass:
$m_{\mathrm{obs}}=C=d_{\mathrm{stable}}$.  In any mapping to empirical
physics, this integer is identified with the measured mass scale.
\end{rsbox}
\end{definition}

\begin{definition}[Luminal Keep-Alive State]\label{def:photon-keepalive}
\begin{rsbox}
A luminal state $n/n$ (where $n=d$, $\Tcharge=0$) is invariant under $\psi$
when self-paired:
\[
\psi\!\left(\frac{n}{n},\,\frac{n}{n}\right) = \left(\frac{n}{n},\,\frac{n}{n}\right).
\]
Luminal states contribute zero to the $\koppa$-ledger ($\koppa(n,n)=0$) and
zero net structural entropy change per step.  They maintain lattice coherence
between $\psi$ events without accumulating rest mass or driving structural
evolution.  This is the RS-native ontic role of the luminal carrier: the
keep-alive signal that prevents lattice decoherence between interactions.
\end{rsbox}
\end{definition}

\begin{proposition}[Reciprocal Pair Photon Emission]\label{prop:reciprocal-photon}
\begin{rsbox}
For any active state $a/b$ with $a,b\neq 0$, the mutual reciprocal pair
$(a/b,\,b/a)$ produces two luminal outputs under $\psi$:
\[
\psi\!\left(\frac{a}{b},\,\frac{b}{a}\right) = \left(\frac{a}{a},\,\frac{b}{b}\right).
\]
Both outputs are luminal ($\Tcharge = 0$) with $\koppa = 0$.
Photon emission is the structural consequence of $\psi$ acting on a
mutual reciprocal pair; it is not assigned.

The mass-gap seed $(1/2,\,2/1)$ is a mutual reciprocal pair.  Its
$\psi$-orbit confirms this directly:
\[
\psi\!\left(\tfrac{1}{2},\tfrac{2}{1}\right) = \left(\tfrac{1}{1},\tfrac{2}{2}\right),
\qquad
\psi\!\left(\tfrac{2}{1},\tfrac{1}{2}\right) = \left(\tfrac{2}{2},\tfrac{1}{1}\right).
\]
At steps~1 and~3 of the mass-gap cycle the pair is in the reciprocal
configuration and $\psi$ emits a luminal pair.  The 4-step mass-gap orbit
therefore alternates: \emph{active} (subluminal $+$ superluminal) $\to$
\emph{luminal pair} $\to$ \emph{active} (reversed) $\to$ \emph{luminal
pair} $\to$ return.  Photons are structurally emitted and reabsorbed every
half-cycle of the mass-gap orbit.
\end{rsbox}
\end{proposition}

\begin{remark}[Superluminal states as virtual carriers]\label{rem:virtual-carriers}
States with $\Tcharge < 0$ (superluminal, $n > d$) contribute zero to both
the $\koppa$-ledger and the structural entropy ledger.  When the denominator
is $d=1$: $\koppa(n,1)=0$ and $\rank(1)=0$.  They are transparent to both
ledgers. They are present in the $\psi$ dynamics, driving the cycle forward, but
leaving no rest mass or causal depth record at $\Omega$.  Their structural
role is that of virtual carriers: they carry the cross-reciprocal operation
forward without themselves becoming observable states.
\end{remark}

\begin{theorem}[Three-Moment Cycle Structure]\label{thm:three-moment}
\begin{rsbox}
The $\psi$-orbit of the mass-gap seed $(1/2,\, 2/1)$ has period 4 and produces
the following fixed temporal charge pattern across one complete cycle:
\begin{align*}
s_0 &= (1/2,\, 2/1): & \Tcharge = (+3, -3) & \quad \text{particle + virtual carrier} \\
s_1 &= (1/1,\, 2/2): & \Tcharge = (0,\; 0) & \quad \text{wave pair (photon emission)} \\
s_2 &= (2/1,\, 1/2): & \Tcharge = (-3, +3) & \quad \text{virtual carrier + particle} \\
s_3 &= (2/2,\, 1/1): & \Tcharge = (0,\; 0) & \quad \text{wave pair (photon emission)} \\
s_4 &= s_0 & & \quad \text{return}
\end{align*}
The three structural roles are particle (subluminal, $\Tcharge > 0$), wave
(luminal, $\Tcharge = 0$), and virtual carrier (superluminal, $\Tcharge < 0$).
Time accumulates only during the particle steps.  The wave steps are timeless:
they propagate structural information without contributing to any ledger and
without advancing the proper time count.  The virtual carrier drives the cycle
forward without becoming observable.
\end{rsbox}
\begin{proof}
The four states $s_0$ through $s_3$ are computed by direct application of
Definition~\ref{def:rop}.  The temporal charge at each component is computed by
Definition~\ref{def:temporal-charge}: $\Tcharge(1,2) = 4-1 = 3$,
$\Tcharge(2,1) = 1-4 = -3$, $\Tcharge(1,1) = 1-1 = 0$, $\Tcharge(2,2) = 4-4
= 0$.  The period equals 4 because $s_4 = s_0$ by direct computation, and $s_k
\neq s_0$ for $k \in \{1, 2, 3\}$ by inspection of the state table.  No step of
this proof uses any arithmetic beyond integer addition and subtraction.
\end{proof}
\end{theorem}

\begin{lemma}[Mass-Gap Cycle Koppa Accumulation]\label{lem:koppa-cycle}
\begin{rsbox}
The koppa accumulated by the mass-gap cycle $(1/2,\, 2/1)$ over one complete
four-step $\psi$-period is $G_{\mathrm{mg}} = 2$.  This equals the mass gap integer.
\end{rsbox}
\begin{proof}
We apply Definition~\ref{def:koppa} at each component of each step.

Step 0, first component $1/2$: initialise $r = 1$.  Since $r < d = 2$, no
subtraction is performed.  Result: $1 \koppa 2 = 1$.

Step 0, second component $2/1$: initialise $r = 2$.  Subtract $d = 1$: $r = 1$.
Subtract again: $r = 0$.  Result: $2 \koppa 1 = 0$.  The superluminal component
contributes zero.

Step 1, first component $1/1$: $r = 1$.  Subtract $d = 1$: $r = 0$.  Result:
$1 \koppa 1 = 0$.  The luminal state contributes zero.

Step 1, second component $2/2$: $r = 2$.  Subtract $d = 2$: $r = 0$.  Result:
$2 \koppa 2 = 0$.

Step 2, first component $2/1$: $2 \koppa 1 = 0$.

Step 2, second component $1/2$: $1 \koppa 2 = 1$.  The subluminal component
(roles exchanged under $\psi^2$) contributes 1.

Step 3: both components luminal; both koppa values are 0.

Total koppa over steps 0 through 3: $1 + 0 + 0 + 0 + 0 + 1 + 0 + 0 = 2 =
G_{\mathrm{mg}}$.
\end{proof}
\end{lemma}

\begin{proposition}[Temporal Surplus Localised at Particle Steps]\label{prop:time}
\begin{rsbox}
In the mass-gap $\psi$-cycle, positive temporal charge, the condition for
proper time accumulation, occurs only at the subluminal (particle) steps.
The luminal steps carry $\Tcharge = 0$ and contribute zero to any time
accumulation.  The superluminal component at each particle step carries
$\Tcharge < 0$ but contributes zero to both ledgers and is structurally virtual.
Proper time accumulates in discrete units, one per particle step, and is absent
at wave steps.  Time is not a background parameter.  It is the count of
particle-phase steps.
\end{rsbox}
\begin{proof}
From Theorem~\ref{thm:three-moment}: at steps 1 and 3, both components have
$\Tcharge = 0$.  At steps 0 and 2, one component has $\Tcharge > 0$ (subluminal)
and one has $\Tcharge < 0$ (superluminal, virtual carrier).  By
Definition~\ref{def:temporal-charge} and Definition~\ref{def:two-ledgers}, the
superluminal component contributes zero to both ledgers.  Positive temporal
charge, and therefore proper time accumulation, occurs at exactly two steps
out of four.
\end{proof}
\end{proposition}

\begin{proposition}[Relativistic Compression as Temporal Charge Reduction]\label{prop:relcomp}
\begin{rsbox}
In RS, a particle gains energy when its state $n/d$ has $n$ increasing toward $d$
with $d$ fixed, reducing the temporal charge $\Tcharge = d^2 - n^2$.  As
$\Tcharge \to 0$, the fraction of the $\psi$-cycle spent in the subluminal phase
decreases.  The net temporal surplus accumulated per vacuum cycle decreases
proportionally.  The particle's proper time accumulates more slowly relative to a
lower-energy reference.  This is the RS structural origin of relativistic time
dilation.  No external metric, no curved spacetime, and no real-valued velocity
parameter enters the derivation.

At the luminal limit $\Tcharge = 0$, no proper time accumulates.  The photon is
the realisation of this limit: a luminal state $n/n$ that propagates without
experiencing time (Definition~\ref{def:photon-keepalive}).  Space-like separation
arises from the compression of the mediant geodesic between lattice positions as
the state approaches the luminal boundary.
\end{rsbox}
\begin{proof}
From Definition~\ref{def:temporal-charge}, $\Tcharge(n,d) = d^2 - n^2$, which
is strictly decreasing in $n$ for fixed $d > 0$ and $0 < n < d$.  At $n = d$,
$\Tcharge = 0$.  From Proposition~\ref{prop:time}, proper time accumulates only
at subluminal steps.  As $n$ increases toward $d$, the state spends more of its
extended vacuum orbit in the near-luminal regime, reducing the time-accumulation
rate.  At $\Tcharge = 0$ the state is luminal and no time accumulates.
\end{proof}
\end{proposition}

\newpage


\section{Definitions: Matter Generator}\label{sec:mattergenerator}

The projective matter generator $G:\SP\to\SP$ is specified by three cases.
All outputs involve only additions and multiplications in $\Z$; no division is
performed.  Weierstrass coefficients $a,b\in\Z_{\geq 0}$ are system
parameters.  The underlying curve is $Y^2Z=X^3+aXZ^2+bZ^3$ in projective
coordinates.

\begin{definition}[Case A: Identity]\label{def:caseA}
If $P_2=\IP=(0,1,0)$, return $P_1$.
\end{definition}

\begin{definition}[Case B: Doubling]\label{def:caseB}
\begin{rsbox}
If $P_1=P_2=(X_1,Y_1,Z_1)\in\SPs$, compute the following intermediate
integers:
\begin{align*}
W_B &:= 3X_1^2 + a Z_1^2 \\
S_B &:= Y_1 Z_1 \\
B_B &:= X_1 Y_1 S_B \\
H_B &:= W_B^2 - 8B_B
\end{align*}
Output coordinates:
\begin{align*}
X_3 &:= 2 H_B S_B \\
Y_3 &:= W_B(4B_B - H_B) - 8Y_1^2 S_B^2 \\
Z_3 &:= 8 S_B^3
\end{align*}
Every operation is integer multiplication or addition; no division appears.
Note that $Z_3=8S_B^3=8Y_1^3Z_1^3$, which is distinct from the Case C
intermediate $W:=Z_1Z_2$.  For any input $(X_1,Y_1,Z_1)\in\SPs$ on the
curve, $(X_3,Y_3,Z_3)\in\SPs$ lies on the same curve.

Output rank: $\rank(Z_3)=3+3\rank(Y_1)+3\rank(Z_1)\;>\;\rank(Z_1)$ for all
$Y_1\neq 0$, enforcing the Arrow of Time axiom (Axiom~\ref{ax:arrow}).
\end{rsbox}
\end{definition}

\begin{definition}[Case C: Addition]\label{def:caseC}
If $P_1,P_2\in\SPs$ are projectively distinct ($P_1\neq P_2$):
\begin{align*}
U &:= Y_2 Z_1-Y_1 Z_2, \quad V := X_2 Z_1-X_1 Z_2, \quad W := Z_1 Z_2 \\
A &:= U^2 W-V^3-2V^2 X_1 Z_2 \\
X_3 &:= VA \\
Y_3 &:= U(V^2 X_1 Z_2-A)-V^3 Y_1 Z_2 \\
Z_3 &:= V^3 W
\end{align*}
Non-degeneracy: for projectively distinct $P_1,P_2\in\SPs$, $V\neq 0$,
so $Z_3\neq 0$ and $P_3\in\SPs$.
\end{definition}

\begin{definition}[Projective ZERO Logic]\label{def:proj-zero}
For any matter generation step $G:\SP\times\SP\to\SP$, if any component of
the resulting triple $(X,Y,Z)$ is $0$, the state enters
Projective Suspension.  The system suspends to the Projective
Identity $\IP=(0,1,0)$.  Evolution is frozen until the suspension is
resolved by $\psi$ acting on the suspended projective state paired with
an active partner.  The cross-reciprocal form moves the structural history
into an active configuration; capacity information is distributed to the
structural entropy ledger at $\Omega$.  Information is preserved, not lost.
\end{definition}

\newpage


\section{Definitions: Time Substrate and Phase Logic}\label{sec:time}

\begin{definition}[Prime-11 time cycle $\Tcyc$]\label{def:T11}
$\Tcyc:=\{1,2,\dots,11\}$, organised as a finite directed graph with
successor function $\succ$; the temporal container.
\end{definition}

\begin{definition}[Successor with reset]\label{def:succ}
\[
\succ(\tau):=\begin{cases}\tau+1 & \tau<11 \\ 1 & \tau=11\end{cases}
\]
The transition $11\to 1$ is the Cycle Boundary Event $\Omega$.
\end{definition}

\begin{definition}[Phase evolution]\label{def:phase-evol}
$\phi(0):=1$; $\phi(t+1):=\succ(\phi(t))$.  Integer-valued; no division.
The initial value $\tau_0:=1$ is the canonical convention; the Viscosity Sum
$\VS$ is invariant under cyclic permutation of the starting point
(Definition~\ref{def:viscosity}).
\end{definition}

\begin{definition}[Triadic phase set and 4-4-3 partition]\label{def:443}
$\PhiSet:=\{E,M,R\}$.  The pointer map $f:\Tcyc\to\PhiSet$ cycles
$E,M,R,E,M,\dots$ over ticks $1,\dots,11$, yielding:

\medskip
\begin{center}
\begin{tabular}{lcccccccccccc}
\toprule
Tick $\tau$ & 1 & 2 & 3 & 4 & 5 & 6 & 7 & 8 & 9 & 10 & 11 \\
\midrule
Phase $f(\tau)$ & $E$ & $M$ & $R$ & $E$ & $M$ & $R$ & $E$ & $M$
  & $R$ & $E$ & $M$ \\
\bottomrule
\end{tabular}
\end{center}
\medskip

Derived partition: $N_E=4$, $N_M=4$, $N_R=3$.
\end{definition}

\begin{definition}[Boundary-Active Phase Sum $\Sigimb$]\label{def:sigimb}
\begin{rsbox}
\[
\Sigimb := N_E + N_R = 4 + 3 = 7.
\]
$\Sigimb$ is the sum of the two boundary-active phase cardinalities: the
Emission phase ($N_E=4$, ticks $\{1,4,7,10\}$, temporal contribution $+1$
each) and the Return phase ($N_R=3$, ticks $\{3,6,9\}$, temporal
contribution $-1$ each).  The Modulation phase ($N_M=4$) is excluded: it
carries structural history between cycle boundaries, contributes zero temporal
charge per tick, and generates no boundary stress.

An equivalent path: $\Sigimb=\Delta+\delta_{\mathrm{slip}}=5+2=7$.  The path
through $N_E+N_R$ is more primitive, requiring only the phase cardinalities.
\end{rsbox}
\end{definition}

\begin{theorem}[Inertial Closure at $\tau_{10}$]\label{thm:tau10}
\begin{rsbox}
Tracing the running totals of Emission and Return ticks through the cycle:

\medskip
\begin{center}
\small
\begin{tabular}{cccccc}
\toprule
$\tau$ & Phase & Cumul.\ $E$ & Cumul.\ $R$ & Net surplus & Arrow status \\
\midrule
1  & $E$ & 1 & 0 & $+1$ & \\
2  & $M$ & 1 & 0 & $+1$ & \\
3  & $R$ & 1 & 1 & $0$  & \\
4  & $E$ & 2 & 1 & $+1$ & \\
5  & $M$ & 2 & 1 & $+1$ & \\
6  & $R$ & 2 & 2 & $0$  & \\
7  & $E$ & 3 & 2 & $+1$ & \\
8  & $M$ & 3 & 2 & $+1$ & \\
9  & $R$ & 3 & 3 & $0$  & \textit{last zero-surplus moment} \\
10 & $E$ & 4 & 3 & $+1$ & \textbf{Arrow irreversibly fixed} \\
11 & $M$ & 4 & 3 & $+1$ & \\
\bottomrule
\end{tabular}
\end{center}
\medskip

At $\tau=9$ the accumulated emission and return counts are $E=3$, $R=3$: the
cycle is in perfect temporal balance, net surplus zero.  No directional bias
has accumulated.

At $\tau=10$ the fourth emission fires: $E=4$, $R=3$, net surplus $=+1$.
The remaining tick $\tau=11$ is a Modulation tick contributing zero temporal
charge, and $\Omega$ then resets.  No subsequent tick within the same cycle
can cancel this surplus.  The temporal Arrow is locked.

$\tau_{10}:=T_{\mathrm{vac}}-1=10$ is therefore the inertial closure
tick: the unique microtick at which the temporal surplus becomes permanent.
The count of zero-surplus instants before this moment is exactly $N_R=3$,
at $\tau\in\{3,6,9\}$: the arrowless instants.

Two structural identities follow immediately:
\[
\tau_{10} = T_{\mathrm{vac}}-1 = 10, \qquad
\tau_{10}+\delta_{\mathrm{slip}} = 10+2 = 12 = \Dg.
\]
The gauge manifold dimension equals the inertial closure index extended by
exactly one phase slip.
\end{rsbox}
\end{theorem}

\begin{definition}[Barycentric structure $B$]\label{def:bary-struct}
A stabilised triadic system precipitates a barycentric structure
$B:=(L_E,L_M,L_R)$ where $L_\phi\in\SL$; stabilisation is defined by entry
into a discrete bracket with $W(L,U)=1$ maintained for one full triadic
period.  No real-valued limit is invoked.
\end{definition}

\newpage


\section{Definitions: Measurement Functions}\label{sec:measurement}

\begin{definition}[Viscosity Sum Operator $\VS$]\label{def:viscosity}
\begin{rsbox}
The Viscosity Sum $\VS$ is the accumulated phase drift over a closed causal
orbit, quantifying discrete gravitational curvature.  For an orbit of period
$T_N$:
\[
\VS := \sum_{t=0}^{T_N-1}(\delta_{\mathrm{slip}}\koppa\phi(t))
\]
where the result is an integer in $\Z$ representing the net displacement from
the ideal vacuum geodesic.  $\VS=0$ defines a flat Minkowski-equivalent
integer path.

For the standard vacuum orbit traversing all eleven microticks of $\Tcyc$ in
sequence with $\tau_0:=1$, the constructive remainder
$\delta_{\mathrm{slip}}\koppa\tau$ evaluates to $0$ at $\tau\in\{1,2\}$ and
to $2$ at $\tau\in\{3,\dots,11\}$, giving $\VS=18$ as the vacuum baseline
curvature measure.

The value $\VS=18$ is invariant under cyclic permutation of the starting
tick: since the 11-tick orbit is a complete cycle through all elements of
$\Tcyc$, the sum $\sum_{\tau=1}^{11}\delta_{\mathrm{slip}}\koppa\tau=18$ is
independent of the initial phase $\tau_0$.  By convention, $\tau_0:=1$.

\begin{observebox}
Three independent structural paths converge to the vacuum baseline $\VS=18$:
(i) the direct sum above; (ii) $\Sigimb+T_{\mathrm{vac}}=7+11=18$; (iii)
$\delta_{\mathrm{slip}}\cdot N^2=2\cdot 9=18$.  These three paths, causal
orbit measurement, boundary-active phase sum extended by vacuum period, and
phase slip scaled by triadic square converge to the same integer by a
structural identity.
\end{observebox}
\end{rsbox}
\end{definition}

\begin{definition}[Mass Gap $G_{\mathrm{mg}}$]\label{def:Gmg}
\begin{rsbox}
\[
G_{\mathrm{mg}} \;:=\; T_{\mathrm{vac}} \koppa N \;=\; 11 \koppa 3 \;=\; 2.
\]
The mass gap is derived, not primitive.  It is the irresolvable per-cycle
phase slip: the constructive remainder when the prime vacuum period
$T_{\mathrm{vac}}=11$ is measured against the triadic period $N=3$.
Because $T_{\mathrm{vac}}$ is prime and $\gcd(T_{\mathrm{vac}},N)=1$, this
residual is permanent and accumulates without cancellation.  $G_{\mathrm{mg}}$
is the deposit made by $\Omega$ at each cycle boundary
(Definition~\ref{def:Omega}) and the base unit of the structural rank.
\end{rsbox}
\end{definition}

\begin{definition}[Structural Rank $\rank$]\label{def:rank}
\begin{rsbox}
For a positive integer $m$, $\rank(m)$ is the unique $k\geq 0$ satisfying
\[
G_{\mathrm{mg}}^k \;\leq\; m \;<\; G_{\mathrm{mg}}^{k+1}.
\]
\textbf{RS-native computation:} initialise $\mathit{power}:=1$ and
$k:=0$.  While $\mathit{power}\times G_{\mathrm{mg}}\leq m$, set
$\mathit{power}:=\mathit{power}\times G_{\mathrm{mg}}$ and
$k:=k+1$.  Return $k$.  Only integer multiplication and comparison appear;
no floor function, logarithm, or division is used.  Define $\rank(0):=0$.

The rank measures causal depth in units of the mass gap: the number of
$G_{\mathrm{mg}}$-doublings required to reach depth $m$.  It is the
RS-native discrete logarithm base $G_{\mathrm{mg}}$.
\end{rsbox}
\end{definition}

\begin{definition}[Prime-Event Depth $\rho_\Omega$]\label{def:primeventdepth}
\begin{rsbox}
For a positive integer $m$, $\rho_\Omega(m)$ is the total number of prime
factors of $m$ counted with multiplicity.

\textbf{RS-native computation:} initialise $\mathit{remaining}:=m$ and
$\mathit{count}:=0$.  For each prime $p$ (detected by $\koppa$-primality
testing): while $\mathit{remaining}\koppa p = 0$, extract one factor of $p$
by iterative subtraction ($\mathit{remaining}:=\mathit{remaining}/p$ via
repeated subtraction) and increment $\mathit{count}$.  Advance to the next
prime.  Return $\mathit{count}$.  No division operator appears.
Define $\rho_\Omega(0):=\rho_\Omega(1):=0$.

The prime-event depth measures the number of prime-gated $\psi$-events
required to fully decompose state $m$ into irreducible components.
Key structural values: $\rho_\Omega(\mathcal{V}_S\!=\!18)=N=3$
(generation count); $\rho_\Omega(\Lambda_\mu\!=\!207)=N=3$;
$\rho_\Omega(\Lambda_\tau\!=\!3481)=2$.

\textbf{Distinction from $\rank$:} $\rank(m)$ measures binary causal depth
(doublings of $G_{\mathrm{mg}}$); $\rho_\Omega(m)$ measures prime-event
decomposition depth.  Both are RS-native; they measure different structural
properties and are used in different contexts.
\end{rsbox}
\end{definition}

\begin{definition}[Height function $\hfn(m)$]\label{def:height}
For $m\neq 0$: unique $k\geq 0$ s.t.\ $2^{2k}\leq m^2<2^{2k+2}$.
Sign-independent: $\hfn(m)=\hfn(-m)$.  $\hfn(0):=0$.
\end{definition}

\begin{definition}[Structural entropy $H(s)$]\label{def:entropy}
For $s=(n,d)\in\SL$: $H(s):=\rank(d)$.  Quantifies causal depth of the
metric basis.
\end{definition}

\begin{definition}[Constructive remainder $\rem(a,n)$]\label{def:rem}
This operator is structurally identical to the Imbalance Operator $\koppa$. The symbol $\rem$ is redundant. All constructive remainder operations are executed exclusively via $\koppa$.
\end{definition}

\begin{definition}[Inertial mass and rest mass]\label{def:inertial-mass}
$m_I\equiv m_{\mathrm{rest}}\equiv\rank(\koppa_{\mathrm{ledger}})$.  Equality
is a definition from single origin, not a measured coincidence.
\end{definition}

\begin{definition}[Temporal Charge $\Tcharge$]\label{def:temporal-charge}
For $(n,d)\in\SL$: $\Tcharge(n,d):=d^2-n^2$.  Temporal regimes:
$\Tcharge>0$ (subluminal/massive); $\Tcharge=0$ (luminal; $n=d$);
$\Tcharge<0$ (superluminal; excluded).
\end{definition}

\begin{theorem}[Neutrino Refresh Period]\label{thm:neutrino-refresh}
\begin{rsbox}
The $\koppa$-ledger of the mass-gap seed $(1/2, 2/1)$ evolving under $\psi$
accumulates over successive $T_{\mathrm{vac}}$-step intervals in the
alternating pattern:
\[
\kappa_1 = \Delta = 5,\qquad
\kappa_2 = \Delta+1 = 6,\qquad
\kappa_3 = 5,\qquad
\kappa_4 = 6,\quad\ldots
\]
Each consecutive pair sums to $T_{\mathrm{vac}} = 11$.  The ledger reaches
$T_{\mathrm{vac}}$ for the first time after exactly two vacuum cycles:
\[
\text{Neutrino refresh period}
= 2 \cdot T_{\mathrm{vac}}
= G_{\mathrm{mg}} \cdot T_{\mathrm{vac}}
= 22.
\]
At the refresh boundary, $\koppa_{\mathrm{ledger}} = T_{\mathrm{vac}} = 11$,
and $\rank(11) = 1$: the neutrino injects exactly one unit of structural
rest mass into the lattice at each refresh.  The alternating
$\Delta,\,\Delta+1$ pattern is $\delta_{\mathrm{slip}}$ operating at the
ledger boundary.
\end{rsbox}
\end{theorem}

\newpage


\section{Definitions: ZERO Logic}\label{sec:zero}

Zero in RS is not a number. It is absence. A lattice position either contains
structural content or it does not, and the labels $0/0$ and $0/1$ both denote the
same thing: an empty position.  The difference between them is not ontological, it
is contextual.  $0/0$ is what an empty position looks like from a mediant propagation
context; $0/1$ is what the same emptiness looks like from a fractional or
multiplicative context.  Because GCD reduction is axiomatically forbidden, these two
labels are never collapsed into each other by a rule. They are the same absence
encountered from different angles.  When any active state arrives at an empty
position it occupies that position unchanged, because there is nothing present to
interact with, resist, or alter it.  This is not a convention imposed on the
framework; it is what absence means.

The structurally important case is $0/d$ where $d \neq 0$.  Here the numerator is
absent, there is no magnitude, no kinematic imbalance. But the denominator
carries genuine structural history. The position is empty in the physical sense yet
non-trivial in the causal sense.  The koppa operator reads this denominator and
accumulates it into the structural entropy ledger at $\Omega$, where it registers as
causal depth rather than rest mass.  A state $0/7$, for example, contributes nothing
to the kinematic ledger but deposits $\rank(7)$ of causal depth: it could represent a
residual gravitational field, a phase-locked inertial depth, or a vacuum tension
resolved but not kinematically expressed.  Information is not lost; it changes form.
There is one edge case: if the arriving state has denominator 1, as $0/1$ does, 
then $\rank(1) = 0$ and the koppa balance receives no increment.  The ledger continues
to accumulate silently through subsequent steps and is expressed in the lattice when
the next causal event provides the opportunity.

The opposite configuration, $n/0$ with $n \neq 0$, is how RS handles what continuum
mathematics calls a singularity or an infinity.  The numerator preserves structural
history but the metric basis is absent: the state is suspended, not destroyed.  It
cannot evolve until $\psi$ acts on it paired with an active partner, at which point the
history is redistributed into both ledgers and the state returns to active evolution.
No information is discarded; the suspension is a temporary form-change, not a
termination.  Two particle roles emerge directly from this zero-logic structure.  The
photon is a luminal state $n/n$ with temporal charge $\Tcharge = 0$: it contributes
zero to the koppa ledger, is invariant under self-paired $\psi$, and maintains lattice
coherence between interactions without accumulating rest mass.  It is the structural
keep-alive signal.  The neutrino governs the refresh rate of the lattice: the mass-gap
koppa ledger accumulates $\Delta = 5$ in one vacuum cycle and $\Delta + 1 = 6$ in the
next, each consecutive pair summing to $T_{\mathrm{vac}} = 11$, so that after exactly
$G_{\mathrm{mg}} \cdot T_{\mathrm{vac}} = 22$ steps the ledger reaches $T_{\mathrm{vac}}$
and $\rank(T_{\mathrm{vac}}) = 1$: one unit of structural rest mass injected into the
lattice at each refresh.

\begin{definition}[Collapse Rule for Zero Magnitude]\label{def:num-zero}
The state $0/d$ where $d\neq 0$ is not a valid physical state. Zero magnitude means the lattice position is empty. The state $0/d$ remains $0/d$ to preserve its geometric capacity for the structural entropy ledger at $\Omega$, but acts as an empty position ($0/0$ or $0/1$) for kinematic propagation. The label belongs to the propagation context, not to the point. Both are the same emptiness encountered from different propagation contexts.
\end{definition}

\begin{definition}[Constraint Vacuum]\label{def:constraint-vac}
State $n/0$, $n\neq 0$.  Suspension state: structural history $n$ preserved
in numerator; metric basis absent.  Resolved by $\psi$ paired with an active
partner; history redistributed to the kinematic and structural entropy
ledgers at $\Omega$ (Definition~\ref{def:zero-resolution}).
\end{definition}

\begin{definition}[Suspension resolution via $\psi$]\label{def:zero-resolution}
For a suspended state $n/0$ paired with active state $c/d$:
\[
\psi\!\left(\frac{n}{0},\,\frac{c}{d}\right) = \left(\frac{d}{n},\,\frac{0}{c}\right).
\]
The first output $d/n$ is resolved and active: structural history $n$ is
preserved as the new denominator (metric basis), and the partner's denominator
$d$ becomes the new numerator.  The state returns to active evolution.

The second output $0/c$ carries the partner's magnitude $c$ into
the structural ledger.  It contributes zero to the $\koppa$-ledger
($\koppa(0,c)=0$ for all $c>0$) and contributes $\rank(c)$ to the structural
entropy ledger at $\Omega$.  Information is not lost; it is redistributed
between the kinematic and causal-depth ledgers.
This is the mechanism by which the 1D $\psi$ operation propagates
structural information into 3D expression via the koppa/omega ledger system
(Definition~\ref{def:two-ledgers}).
\end{definition}

\begin{definition}[$\Omega$: Cycle-Boundary Resolution Operator]
\label{def:Omega}
\begin{rsbox}
$\Omega$ acts at the cycle boundary after each $T_{\mathrm{vac}}$-step
vacuum period.  It is a zero-logic resolution event: the phase slip
$G_{\mathrm{mg}} = T_{\mathrm{vac}}\koppa N$ has accumulated in the koppa
ledger as structural content without a temporal basis. A suspension state
carrying history but no metric ground.  $\Omega$ resolves this suspension.

\medskip
\textbf{Inputs.}
The particle ERP state $(n,d)\in\SLpos$ at the cycle boundary;
the accumulated koppa ledger $K\in\Z_{\geq 0}$ over the preceding
$T_{\mathrm{vac}}$ steps.

\medskip
\textbf{Outputs.}
\begin{itemize}
\item Temporal deposit: $T_{\mathrm{dep}} := \Tcharge(n,d) = d^2 - n^2$,
  deposited into the structural time ledger.
\item Kinematic carry: $K_{\mathrm{fwd}} := K$, carried forward into the
  next cycle's koppa ledger.
\item Temporal counter advance: $t \mapsto t+1$ (one unit of structural
  time per $\Omega$ firing, unconditional).
\end{itemize}

\medskip
\textbf{Time dilation.}  The temporal deposit $T_{\mathrm{dep}} =
d^2-n^2$ is the amount of proper time registered at this $\Omega$ event.
For a luminal state $n=d$: $T_{\mathrm{dep}}=0$; no proper time
accumulates.  For a subluminal state $n<d$: $T_{\mathrm{dep}}>0$,
strictly decreasing as $n\to d$.  Time dilation is the reduction of
$T_{\mathrm{dep}}$ as the state approaches the luminal boundary; it is
a property of the ERP integer structure, not an externally applied
transformation.

\medskip
\textbf{Arrow of Time.}  The temporal counter advances by $+1$ at every
$\Omega$ firing unconditionally.  The direction is fixed because
$G_{\mathrm{mg}} = T_{\mathrm{vac}}\koppa N > 0$ always; the phase slip
is always positive and the engine always runs forward.

\medskip
\textbf{Structural identity.}  For all $(n,d)\in\SLpos$:
$d^2 = \Tcharge(n,d) + n^2$.  The geometric capacity squared partitions
exactly into the temporal deposit and the kinetic square.  No remainder.
This is the RS-native integer metric signature.
\end{rsbox}
\end{definition}

\begin{definition}[Projective ZERO Logic]\label{def:proj-zero-logic}
See Definition~\ref{def:proj-zero} in Section~\ref{sec:mattergenerator}.
\end{definition}

\begin{definition}[Empty State]\label{def:null}
The label $0/0$ denotes an empty lattice position as seen from a mediant
propagation context.  There is no structural content at this position.
It is not a special persistent algebraic object; it is the same emptiness
as $0/1$, encountered from a different propagation context.
Nothing is there to interact with, resist, or alter an incoming state.
An incoming state occupies an empty position unchanged.
This is not a rule imposed on the framework; it is what emptiness means.
\end{definition}

\newpage


\section{Definitions: Rational Special Relativity}\label{sec:RSR}

\begin{postulate}[Dimensional Necessity]\label{post:dim-necessity}
The gauge manifold dimensionality $\Dg=12$ is the minimal integer space
required to support a non-trivial Phase Slip ($\delta_{\mathrm{slip}}=2$)
across a Triadic Logic ($N=3$) while maintaining vacuum synchronisation.
It is the unique solution to the requirement that the interaction surface
$\Sint$ must scan all logical phase states without topological overlap in a
single barycentric period.
\end{postulate}

\begin{theorem}[Exhaustive Uniqueness of $\Dg$]\label{thm:Dg-unique}
The gauge dimension $\Dg=12$ is the unique minimal integer solution
satisfying the synchronisation condition that $\delta_{\mathrm{slip}}$
divides $\Sint$ within a triadic container.
\begin{proof}
Five conditions jointly constrain $\Dg$.

Condition~(1): $\Dg\equiv 0\pmod{N}$.  The gauge manifold must be
triadic-compatible; $N=3$ divides $\Dg$.

Condition~(2): $\Dg\equiv 0\pmod{2}$.  The interaction surface must
synchronise with the phase slip $\delta_{\mathrm{slip}}=2$, requiring
$\delta_{\mathrm{slip}}$ to divide $\Sint$, which forces
$\Dg$ to be even.

Condition~(3): $T_{\mathrm{vac}}=11$ is prime and does not divide $\Dg$.
Since 11 and $\Dg$ share no common factor, the minimal integer scanned
by both periods is their product: $\Sint=T_{\mathrm{vac}}\cdot\Dg=11\Dg$.

Condition~(4): $\Sint$ is a whole multiple of $\Tbary=33$.  Requiring
$33\mid 11\Dg$ gives $3\mid\Dg$, already enforced by Condition~(1).

From Conditions~(1)--(4), admissible candidates are
$\Dg\in\{6,12,18,24,\ldots\}$, yielding
$\Sint\in\{66,132,198,264,\ldots\}$.

Condition~(5): $\Sint+\Delta=\PG(T_{\mathrm{bary}})$; the interaction
surface offset by the flux imbalance must be the prime at position
$T_{\mathrm{bary}}=33$.  Evaluating: $\Dg=6$ gives
$\Sint+\Delta=71=\PG(20)$; $\Dg=12$ gives
$\Sint+\Delta=137=\PG(33)=\PG(T_{\mathrm{bary}})$.  No candidate with
$\Dg<12$ satisfies this condition.  Therefore $\Dg=12$ is the unique minimal
solution.
\end{proof}
\end{theorem}

\begin{definition}[Rational event vector]\label{def:event}
$E=\bigl((n_t,d_t),(n_x,d_x)\bigr)\in\SL\times\SL$.
\end{definition}

\begin{definition}[Boost parameter and integer components]\label{def:boost}
$U=(n_u,d_u)\in\SL$.  Integer components:
\[
\alpha_{\mathrm{boost}}:=d_u^2+n_u^2, \quad
\beta_{\mathrm{boost}}:=2n_u d_u, \quad
\Delta_{\mathrm{boost}}:=d_u^2-n_u^2.
\]
$\Delta_{\mathrm{boost}}=0$ at the luminal boundary.
\end{definition}

\begin{definition}[Rational Lorentz boost]\label{def:Lorentz}
$A_Q:=\bigl(\begin{smallmatrix}\alpha_{\mathrm{boost}}&\beta_{\mathrm{boost}}
\\\beta_{\mathrm{boost}}&\alpha_{\mathrm{boost}}\end{smallmatrix}\bigr)$.
For an event vector $E=((n_t,d_t),(n_x,d_x))\in\SL\times\SL$, the boosted
event $E'=((n_t',d_t'),(n_x',d_x'))$ is defined by:
\[
n_t' := \alpha_{\mathrm{boost}}\,n_t\,d_x
       +\beta_{\mathrm{boost}}\,n_x\,d_t, \quad
d_t' := \Delta_{\mathrm{boost}}\,d_t\,d_x,
\]
\[
n_x' := \beta_{\mathrm{boost}}\,n_t\,d_x
       +\alpha_{\mathrm{boost}}\,n_x\,d_t, \quad
d_x' := \Delta_{\mathrm{boost}}\,d_t\,d_x.
\]
The output $E'\in\SL\times\SL$ is produced by integer multiplication and
addition alone; no division appears.
\end{definition}

\begin{definition}[Spacetime interval]\label{def:interval}
$I(E):=(n_t^2 d_x^2-n_x^2 d_t^2,\;d_t^2 d_x^2)\in\SL$.  Proved invariant
under $A_Q$ via $\alpha_{\mathrm{boost}}^2-\beta_{\mathrm{boost}}^2
=\Delta_{\mathrm{boost}}^2$ (Theorem~\ref{thm:interval-inv}).
\end{definition}

\begin{definition}[Rational velocity addition]\label{def:vel-add}
$(U_1,U_2)\mapsto U_{\mathrm{new}}:=(n_1 d_2+n_2 d_1,\;d_1 d_2+n_1 n_2)$.
\end{definition}

\newpage

\newpage


\section{Definitions: Spatial Dynamics and Thermodynamics}\label{sec:dynamics}

\begin{definition}[Positive-Denominator Sublattice $\SLpos$]\label{def:slplus}
$\SLpos:=\{(n,d)\in\Z^2:d\geq 1\}\subset\SL$.  Every observational
equivalence class of $\SL$ contains exactly one representative in $\SLpos$,
obtained by negating both components when $d<0$.  The unit entry state
$(0,1)$ lies in $\SLpos$ as the generative seed of the $\aop$-orbit; it is
not an active physical ERP.  All velocity ERPs and ensemble states are drawn
from $\SLpos$; this ensures $\rank(d)$ and $n\koppa d$ are defined for every
such object.
\end{definition}

\begin{theorem}[Three Spatial Registers]\label{thm:registers}
The active matter domain $\SPs$ provides exactly three structurally equivalent
independent integer registers.  Their count equals the triadic logic period
$N=3$.  No register carries temporal charge.
\begin{proof}
$\SPs\subset\Z^3$ has exactly three component positions.  The condition
$XYZ\neq 0$ (Definition of $\SPs$) treats each component identically; no
additional condition distinguishes $X$ from $Y$ or $Z$.  The count of
components equals $N=3$ (Axiom~\ref{ax:triadic}).  The temporal charge
$\Tcharge(n,d)=d^2-n^2$ is defined on pairs $(n,d)\in\SL$; a bare integer
is not such a pair and carries no $\Tcharge$ value.
\end{proof}
\end{theorem}

\begin{definition}[RS Kinetic Charge $\GK$]\label{def:GK}
For any $U=(n_u,d_u)\in\SLpos$:
\[
\GK(U)\;:=\;d_u^2-\Tcharge(U)\;=\;n_u^2.
\]
\end{definition}

\begin{theorem}[Kinetic-Temporal Complementarity]\label{thm:complementarity}
For all $U=(n_u,d_u)\in\SLpos$: $\;\Tcharge(U)+\GK(U)=d_u^2$.
\begin{proof}
$(d_u^2-n_u^2)+n_u^2=d_u^2$.
\end{proof}
\end{theorem}

\begin{definition}[RS 3-Momentum]\label{def:momentum3}
Let $m_{\mathrm{rest}}:=\rank(\koppa_{\mathrm{ledger}})\in\Z_{\geq 0}$ and
$U_i=(n_i,d_i)\in\SLpos$ for $i=1,2,3$.  The RS 3-Momentum is
\[
P_i\;:=\;m_{\mathrm{rest}}\cdot n_i,\qquad i=1,2,3.
\]
Each component is an integer product; no division appears.
\end{definition}

\begin{definition}[RS 3+1D Spacetime Event $\Ethree$]\label{def:event3p1}
\[
\Ethree\;:=\;\bigl((n_t,d_t),\,(X,Y,Z)\bigr)\;\in\;\SLpos\times\SPs.
\]
The temporal component $(n_t,d_t)\in\SLpos$ carries temporal charge
$\Tcharge(n_t,d_t)=d_t^2-n_t^2$.  The spatial component
$(X,Y,Z)\in\SPs$ encodes three independent non-zero integer coordinates,
one per spatial register.  No velocity ERP appears.
\end{definition}

\begin{theorem}[RS First Law]\label{thm:firstlaw}
Let $U_i=(n_i,d_i)\in\SLpos$ be the velocity ERP of spatial register $i$ at
cycle boundary $\Omega$.  If no coupling deposit acts on register $i$ at
$\Omega$, then $U_i|_{\Omega+1}=U_i|_\Omega$.
\begin{proof}
The three vacuum operators $\lop$, $\aop$, $\psi$
(Definitions~\ref{def:lop}--\ref{def:rop}) each take explicit
ERP arguments and return ERPs.  None accepts a velocity register as a
separate input or modifies $U_i$ unless $U_i$ is the explicit argument
supplied.  By Definition~\ref{def:koppa}, the velocity ERP is modified only
by an explicit coupling deposit via $\med$.  Absent such deposit on register
$i$, $U_i$ is unchanged.
\end{proof}
\end{theorem}

\begin{definition}[RS Force Deposit]\label{def:forcedeposit}
For $U_i=(n_i,d_i)\in\SLpos$ and structural depth $k\geq 0$:
\[
U_i'\;:=\;U_i\med C(k)
\;=\;\bigl(n_i+\Delta,\;\;d_i+L(k+4)\cdot L(k+5)\bigr),
\]
where $\Delta=5$ and $C(k)=(\Delta,\,L(k+4)\cdot L(k+5))$ is the Lucas coupling ERP at depth $k$. Both components are integer additions. The mediant sum $\med$ is used because a force deposit is a width-preserving momentum transfer, not a mass-generating barycentric addition ($\bplus$).
\end{definition}

\begin{theorem}[Force Deposit Closure on $\SLpos$]\label{thm:closure}
For all $U_i\in\SLpos$ and $k\geq 0$: $\;U_i\med C(k)\in\SLpos$.
\begin{proof}
$L(k+4)>0$ and $L(k+5)>0$ (Lucas positivity; direct from
Definition~\ref{def:aop} by induction), so $L(k+4)\cdot L(k+5)\geq 1$.
With $d_i\geq 1$: the output denominator $d_i+L(k+4)\cdot L(k+5)\geq 2\geq 1$.
\end{proof}
\end{theorem}

\begin{theorem}[RS Third Law]\label{thm:thirdlaw}
Let states $A$ and $B$ be in mutual coupling at depth $k\geq 0$, each
receiving the deposit $C(k)$.  Then
\[
\Dcross(U_B,C(k))\;=\;-\,\Dcross(U_A,C(k)).
\]
\begin{proof}
By Definition~\ref{def:Dcross}, $\Dcross$ is skew-symmetric:
$\Dcross(b,a)=-\Dcross(a,b)$ for all $a,b\in\SL$.  In a mutual coupling,
$C(k)$ is the same ERP in both interactions.  Swapping the argument position
of the receiving velocity ERP negates the cross-determinant.
\end{proof}
\end{theorem}

\begin{definition}[RS Ensemble $\Ens$]\label{def:ensemble}
A finite ordered collection $\Ens=\{s_1,\ldots,s_M\}$ with each
$s_j=(n_j,d_j)\in\SLpos$ and $M\in\Z_{>0}$.  The symbol $M$ denotes
ensemble cardinality; it is unrelated to the triadic period constant $N=3$.
\end{definition}

\begin{definition}[RS Thermodynamic Entropy $\SRS$]\label{def:entropy}
For an ensemble $\Ens\subset\SLpos$:
\[
\SRS(\Ens)\;:=\;\sum_{j=1}^{M}\rank(d_j).
\]
\end{definition}

\begin{theorem}[RS Second Law]\label{thm:secondlaw}
$\SRS(\Ens_{t+1})\geq\SRS(\Ens_t)$ under any valid RS evolution step,
with strict inequality whenever the step invokes the doubling generator
on any member of $\Ens$.
\begin{proof}
By Axiom~\ref{ax:arrow}, no valid RS evolution step decreases $\rank(d_j)$
for any state, and the doubling generator strictly increases it.  A finite
sum of non-decreasing integer terms is non-decreasing; a strict increase at
any term makes the sum strictly increase at that step.
\end{proof}
\end{theorem}

\begin{definition}[RS Thermal Charge $\TRS$]\label{def:thermalcharge}
For an ensemble $\Ens\subset\SLpos$:
\[
\TRS(\Ens)\;:=\;\sum_{j=1}^{M}(n_j\koppa d_j).
\]
\end{definition}

\begin{theorem}[Thermal Charge Non-Negativity and Phase-Lock]
\label{thm:thermalcharge}
$\TRS(\Ens)\geq 0$, with equality if and only if $d_j\mid n_j$ for all $j$.
A state satisfying $n_j\koppa d_j=0$ is phase-locked.
\begin{proof}
Each term $n_j\koppa d_j\in[0,d_j)$ is non-negative by
Definition~\ref{def:koppa}.  A finite sum of non-negative integers is
non-negative.  The sum equals zero iff every term equals zero iff
$d_j\mid n_j$ for all $j$.
\end{proof}
\end{theorem}


\section{Definitions: Mediant Tree, Geodesics, and Discrete L-Function}
\label{sec:mediant}

\begin{definition}[Mediant Tree $\mathcal{T}$]\label{def:mediant-tree}
Binary tree of integer pairs organised by mediant insertion; roots $(0,1)$
and $(1,0)$.  Every $(p,q)\in\Z^2$ is reachable.
\end{definition}

\begin{definition}[Mediant geodesic $\gamma(A,B)$]\label{def:geodesic}
For $A\prec B$: the unique path in $\mathcal{T}$ from $A$ to $B$ generated
by recursive mediant insertion.  Length equals the sum of partial quotients
of the continued fraction of $(bc-ad)/(bd)$.
\end{definition}

\begin{definition}[Structural symbol $\{A,B\}$]\label{def:struct-symbol}
Formal sum of mediant geodesic edges from $A$ to $B$.  Manin-Drinfeld
relations: $\{A,B\}+\{B,C\}=\{A,C\}$; $\{A,A\}=0$; $\{A,B\}+\{B,A\}=0$.
\end{definition}

\begin{definition}[Evolution sequence $\{S_t\}$]\label{def:evol-seq}
Generator point $G_{\mathrm{ec}}\in\SPs$ of infinite order.
$S_t:=t\cdot G_{\mathrm{ec}}$ via the polynomial group law; no division.
\end{definition}

\begin{definition}[Discrete L-function $\LQ(N)$]\label{def:L-function}
Let $T_N$ be the orbit period modulo $N$;
$R_N:=\sum_{t=0}^{T_N-1}\sigma(\tau_t)$:
\[
\LQ(N):=(R_N,\,T_N)\in\SL.
\]
Orientation-bias summary; integer analogue of the Hasse-Weil L-function
at $s=1$.
\end{definition}

\newpage


\section{Lucas Architecture and Structural Families}\label{sec:lucas}

\subsection{The Lucas Sequence as the RS Vacuum Architecture}

\begin{definition}[Lucas Number $L(n)$]\label{def:lucas}
\begin{rsbox}
The Lucas sequence is defined by $L(0):=2$, $L(1):=1$,
$L(n):=L(n-1)+L(n-2)$ for $n\geq 2$.  The sequence begins:
$2,1,3,4,7,11,18,29,47,76,\ldots$

It is generated by $\aop$ from the vacuum seed
$(T_{\mathrm{vac}},\Sigimb)=(11,7)=(L(5),L(4))$: repeated $\aop$
application produces the orbit
$(11,7),(18,11),(29,18),(47,29),\ldots$, with second components descending
and first components ascending through successive Lucas numbers.
\end{rsbox}
\end{definition}

\begin{theorem}[Lucas-Vacuum Encoding]\label{thm:lucas-encoding}
\begin{rsbox}
The foundational RS constants are Lucas numbers at structurally motivated
positions:
\[
\begin{array}{lll}
G_{\mathrm{mg}} = 2 = L(0), & N = 3 = L(2), &
  N_E = N_M = 4 = L(3),\\[4pt]
\Sigimb = 7 = L(4), & T_{\mathrm{vac}} = 11 = L(5). &
\end{array}
\]
The vacuum seed is $(L(5),L(4))$ and the $\aop$-orbit produces $\{L(n)\}$
in order.  The gauge dimension is the product of two consecutive Lucas
numbers:
\[
\Dg = 12 = L(2)\times L(3) = N\times N_E.
\]
\end{rsbox}
\end{theorem}

\begin{theorem}[Continued Fraction Triadic Closure of Lucas Ratios]
\label{thm:lucas-cf}
\begin{rsbox}
For adjacent Lucas numbers $L(n+1)$ and $L(n)$, the continued fraction of
$L(n+1)/L(n)$ terminates with the triadic period $N=3$ as the mandatory
final partial quotient:
\[
\mathrm{CF}\!\left(\frac{L(n+1)}{L(n)}\right)
= [\underbrace{1;\,1,\ldots,1}_{n-1\text{ ones}},\,N].
\]
The triadic period is the universal attractor of the Lucas convergent path.
\end{rsbox}
\begin{proof}
Lucas numbers satisfy $L(n)=F(n-1)+F(n+1)$, which introduces a deviation
from the pure Fibonacci path at the base case.  The terminal partial quotient
equals $N=3$ at the first deviation point and propagates by the continued
fraction recursion; the exact derivation follows from the Cassini-type
identity for Lucas numbers.
\end{proof}
\end{theorem}

\subsection{The Boundary-Active Phase Sum and Structural Integer Families}

The canonical RS constants belong to three overlapping structural families.
These families are recognitions of organisation already present in the
constant table, not new axioms.

\begin{theorem}[$T_{\mathrm{vac}}$-Semiprime Generation of Scale Constants]
\label{thm:semiprime}
\begin{rsbox}
\[
\{k\times T_{\mathrm{vac}}:k\in\{G_{\mathrm{mg}},N,\Delta,\Sigimb\}\}
= \{22,\,33,\,55,\,77\}.
\]
These are the Bridge denominator, the Barycentric Period, the anomaly
denominator, and the $\Sigimb$-vacuum semiprime respectively.  Every
member appears as a canonical denominator or period in this reference.

The Barycentric Addition of any two prime-capacity ERPs $(n_1,p_1)$ and
$(n_2,p_2)$ produces a semiprime-capacity state:
$(n_1,p_1)\bplus(n_2,p_2)=(n_1 p_2+n_2 p_1,\;p_1 p_2)$,
with denominator $p_1 p_2$.  Semiprimes are therefore the native output of
$\bplus$ applied to prime-capacity inputs.
\end{rsbox}
\end{theorem}

\begin{theorem}[Goldbach Derivation of the Gauge Dimension]
\label{thm:goldbach-dg}
\begin{rsbox}
\[
\Dg = \Delta+\Sigimb = 5+7 = 12.
\]
The gauge manifold dimension equals the sum of the flux imbalance and the
boundary-active phase sum.  This is a derivation of $\Dg$ that is entirely
RS-internal and requires no reference to the Lie algebra decomposition
$1+3+8=12$.  Both derivations arrive at the same integer; the Goldbach path
is the RS-native foundation.
\end{rsbox}
\end{theorem}

\begin{observebox}
\textbf{Further Goldbach alignments.}
Taking $\{G_{\mathrm{mg}},N,\Delta,\Sigimb,T_{\mathrm{vac}}\}=\{2,3,5,7,11\}$:
\[
G_{\mathrm{mg}}+N=5=\Delta;\qquad
\Delta+G_{\mathrm{mg}}=7=\Sigimb;\qquad
G_{\mathrm{mg}}+T_{\mathrm{vac}}=13=\PG(\Sigimb).
\]
Additionally: $\Delta+N=8=2^3$ (rank-3 pure dyadic boundary) and
$\Sigimb+T_{\mathrm{vac}}=18=\VS=\delta_{\mathrm{slip}}\cdot N^2$ (the
vacuum Viscosity Sum; three independent derivation paths converge here,
cf.\ Definition~\ref{def:viscosity}).
\end{observebox}

\subsection{The Fibonacci Prime Class}\label{sec:fib-primes}

\begin{definition}[Fibonacci Prime Class]\label{def:fib-prime}
\begin{rsbox}
An integer $q$ is a Fibonacci prime if it is both prime and a
member of the denominator sequence of the $\aop$-orbit from the unit
active seed $(1,1)$.  The denominator sequence of that orbit is the
Fibonacci sequence $1,1,2,3,5,8,13,21,34,55,89,\ldots$; the prime members
are $2,3,5,13,89,233,1597,\ldots$

Note: $(0,1)$ is the left boundary label of the mediant tree
(Definition~\ref{def:mediant-tree}), not an active ERP.
$\aop(0,1)=(1,0)$ immediately produces a constraint vacuum, so $(0,1)$
cannot serve as a generative orbit seed under $\aop$.
The correct structural origin is the unit active seed $(1,1)$, whose
$\aop$-orbit has denominator sequence beginning at $1$ and generating
the full Fibonacci sequence.

Fibonacci primes simultaneously satisfy: (i)~structural irreducibility in
the sense of Postulate~\ref{post:frictionless}, and (ii)~native generation
by the fundamental vacuum propagator $\aop$ from the unit active seed.
\end{rsbox}
\end{definition}

\begin{theorem}[Foundational Constants as First Fibonacci Primes]
\label{thm:fib-primes}
\begin{rsbox}
The first three Fibonacci primes are $\{2,3,5\}$, which equal
$\{G_{\mathrm{mg}},N,\Delta\}$ exactly.  The three integers that define the
entire non-equilibrium character of the RS vacuum, the mass gap, the
triadic logic period, and the flux imbalance, are simultaneously the first
three prime outputs of $\aop$.  The vacuum is self-consistent: its generator
produces the very integers that define its own incommensurability structure.
\end{rsbox}
\end{theorem}

\begin{observebox}
\textbf{Higher Fibonacci primes and vacuum signatures.}

$F_7=13$: Fibonacci index $7=\Sigimb$; satisfies $13=T_{\mathrm{vac}}
+G_{\mathrm{mg}}$, $13=\PG(6)$, and $13\koppa 11=2=\delta_{\mathrm{slip}}$.
Sits at the intersection of the Fibonacci prime, Goldbach, and Prime Gate
families simultaneously.

$F_{11}=89$: Fibonacci index $T_{\mathrm{vac}}=11$.  The $T_{\mathrm{vac}}$-th
Fibonacci number is itself a Fibonacci prime, generated by exactly
$T_{\mathrm{vac}}$ applications of $\aop$ from the unit active seed $(1,1)$.

$F_{13}=233$: satisfies $233\koppa 11=2$, $233\koppa 3=2$, and
$233\koppa 33=2$ simultaneously. The most vacuum-coherent Fibonacci prime
in the accessible range.  It carries the phase-slip signature against
$T_{\mathrm{vac}}$, $N$, and $\Tbary$ simultaneously.  Its role in the
generation sector is addressed in Section~\ref{sec:phi-oscillator}.
\end{observebox}

\subsection{The Matter-Gate Fibonacci Prime}

The Structural Observation of Section~12.3 records the RS roles of the first six Fibonacci
primes through $F_{13} = 233$. Each of those primes has its Fibonacci index identified with a
canonical RS structural constant: the indices $3, 4, 5, 7, 11, 13$ correspond to $N$, $N_E$,
$\Delta$, $\Sigma_{\mathrm{imb}}$, $T_{\mathrm{vac}}$, and $\Pi_3$ respectively. The present
section establishes the seventh Fibonacci prime, $F_{17} = 1597$, whose index $17 =
PG(\Sigma_{\mathrm{imb}})$ was already introduced as the boundary-active matter gate prime in
the Resolution of Open Structure~1 (Section~16). Its establishment here completes the Fibonacci
Prime Index Alignment and provides the ontological foundation required for subsequent
multi-sector derivations.

\begin{definition}[Matter-Gate Fibonacci Prime]
\begin{rsbox}
The seventh Fibonacci prime is $F(PG(\Sigma_{\mathrm{imb}})) = F(17) = 1597$.

Primality is confirmed by the $\koppa$-primality test (Definition~3.1):
$1597 \koppa\, m \neq 0$ for all integers $m$ with $1 < m < 1597$. The index
$17 = PG(\Sigma_{\mathrm{imb}}) = PG(7)$ is the prime at the boundary-active phase sum
index, established as the matter gate prime of the composite Higgs mass integer
$\Lambda_H = 244800$ (Section~16). $F_{17} = 1597$ is generated by exactly $17$
applications of $\eta$ from the unit active seed $(1,1)$, the same structural mechanism by
which $F(T_{\mathrm{vac}}) = F(11) = 89 = n_G$ is generated by $T_{\mathrm{vac}}$
applications (Structural Observation, Section~12.3).
\end{rsbox}
\end{definition}

\begin{theorem}[Three-Path Convergence to $F_{17} = 1597$]
\begin{rsbox}
The matter-gate Fibonacci prime $F_{17} = 1597$ is determined by three structurally
independent RS expressions.

\medskip
\noindent\textbf{Path A - Gauge-Cube Interaction Residual.}
\[
1597 = D_{\mathrm{gauge}}^3 - S_{\mathrm{int}} + 1 = 1728 - 132 + 1.
\]
The factor $D_{\mathrm{gauge}}^3 = 1728$ appears as the leading term of the
proton-electron mass ratio (Theorem~14.13:
$\mu = D_{\mathrm{gauge}}^3 + (S_{\mathrm{int}} - 2D_{\mathrm{gauge}}) + 4\delta_A$).
The interaction surface $S_{\mathrm{int}} = 132$ is the established synchronisation count
(Theorem~14.8). $F_{17}$ is the unique integer one unit above the difference between the
gauge cube and the interaction surface. No operation beyond integer subtraction and
addition appears.

\medskip
\noindent\textbf{Path B - Fibonacci Generation at the Matter-Gate Index.}
\[
1597 = F(PG(\Sigma_{\mathrm{imb}})) = F(17).
\]
The $\eta$-orbit from the unit active seed $(1,1)$ reaches $F(17) = 1597$ after exactly $17$
applications of $\eta$. The index $17 = PG(\Sigma_{\mathrm{imb}})$ gates the composite
Higgs mass integer $\Lambda_H = 120^2 \times PG(\Sigma_{\mathrm{imb}})$ (Section~16).
Path B accesses $1597$ through the Fibonacci generation mechanism without reference to
$D_{\mathrm{gauge}}^3$ or $S_{\mathrm{int}}$.

\medskip
\noindent\textbf{Path C -- Gravitational-Depth Matter-Gate Product.}
\[
1597 = n_G \times PG(\Sigma_{\mathrm{imb}}) + \Sigma_{\mathrm{imb}} \times D_{\mathrm{gauge}}
     = 89 \times 17 + 7 \times 12 = 1513 + 84.
\]
The gravitational coupling depth $n_G = 89$ (Resolution of Open Structure~4) scaled by
the matter-gate prime $PG(\Sigma_{\mathrm{imb}}) = 17$ yields $1513$. The boundary-active
gauge product $\Sigma_{\mathrm{imb}} \times D_{\mathrm{gauge}} = 7 \times 12 = 84$ is the
additive completion. No two of the three paths share a derivation chain or a central
organising constant.
\end{rsbox}

\noindent\textit{Structural Rank Property.}
$\rho(F_{17}) = \rho(1597) = 10 = \tau_{10}$, since $2^{10} = 1024 \leq 1597 < 2048 = 2^{11}$.
The structural rank of the matter-gate Fibonacci prime equals the inertial closure tick.
This extends the rank sequence of sector coupling constants established in the Resolution
of Open Structure~5:

\begin{center}
\begin{tabular}{lcc}
\hline
Constant & Value & Structural rank \\
\hline
$\Pi_3$ (Gen-3 quarks)      & 13   & $3 = N$              \\
$PG(\Sigma_{\mathrm{imb}})$ (Higgs/solar) & 17 & $4 = N_E$ \\
$n_G$ (gravity)             & 89   & $6$                  \\
$G$ (generation sector)     & 233  & $7 = \Sigma_{\mathrm{imb}}$ \\
$F_{17}$ (matter gate)      & 1597 & $10 = \tau_{10}$     \\
\hline
\end{tabular}
\end{center}

The ranks $\{N, N_E, 6, \Sigma_{\mathrm{imb}}, \tau_{10}\} = \{3, 4, 6, 7, 10\}$ are the
first five canonical RS structural integers above the mass gap, in strict ascending order.
\end{theorem}

\begin{rsbox}
\noindent\textbf{Koppa Signatures.}

The constructive remainders of $F_{17}$ against the fundamental RS periods:
\begin{align*}
1597 \koppa\, N      &= 1597 \koppa\, 3  = 1 \\
1597 \koppa\, T_{\mathrm{vac}} &= 1597 \koppa\, 11 = 2 = \delta_{\mathrm{slip}} \\
1597 \koppa\, \Delta &= 1597 \koppa\, 5  = 2 = \delta_{\mathrm{slip}} \\
1597 \koppa\, T_{\mathrm{bary}} &= 1597 \koppa\, 33 = 13 = \Pi_3 \\
1597 \koppa\, D_{\mathrm{gauge}} &= 1597 \koppa\, 12 = 1 \\
1597 \koppa\, S_{\mathrm{int}} &= 1597 \koppa\, 132 = 13 = \Pi_3 \\
1597 \koppa\, n_G   &= 1597 \koppa\, 89  = 84 = \Sigma_{\mathrm{imb}} \times D_{\mathrm{gauge}}
\end{align*}

Verification of $1597 \koppa\, T_{\mathrm{bary}}$: initialise $r := 1597$; subtract $33$
iteratively. $48 \times 33 = 1584$; $r := 1597 - 1584 = 13$. Since $13 < 33$, the process
terminates. $1597 \koppa\, 33 = 13 = \Pi_3$. No division appears.

The extended koppa signature table across sector gate constants:

\begin{center}
\begin{tabular}{lcccc}
\hline
Integer $x$ & $x \koppa\, N$ & $x \koppa\, T_{\mathrm{vac}}$ &
$x \koppa\, T_{\mathrm{bary}}$ & $x \koppa\, D_{\mathrm{gauge}}$ \\
\hline
$\Pi_3 = 13$ (Gen-3 quarks)      & 1 & 2 & 13 & 1 \\
$PG(\Sigma_{\mathrm{imb}}) = 17$ (Higgs) & 2 & 6 & 17 & 5 \\
$n_G = 89$ (gravity)             & 2 & 1 & 23 & 5 \\
$G = 233$ (generation)           & 2 & 2 &  2 & 5 \\
$F_{17} = 1597$ (matter gate)    & 1 & 2 & 13 & 1 \\
\hline
\end{tabular}
\end{center}

\noindent\textit{Signature Coincidence.} $F_{17} = 1597$ and $\Pi_3 = 13$ share the
identical koppa signature profile $(1, 2, 13, 1)$ against $(N, T_{\mathrm{vac}},
T_{\mathrm{bary}}, D_{\mathrm{gauge}})$. This follows from
$F_{17} - \Pi_3 = 1597 - 13 = 1584 = 48 \times T_{\mathrm{bary}}$: since $1584$ is
exactly divisible by $T_{\mathrm{bary}}$, both integers lie in the same residue class
modulo every period that divides $T_{\mathrm{bary}}$.

\noindent\textit{Gravitational residue.} $1597 \koppa\, n_G = 84 = \Sigma_{\mathrm{imb}}
\times D_{\mathrm{gauge}}$: the matter-gate Fibonacci prime carries the boundary-active
gauge product as its gravitational depth residue, consistent with Path~C of Theorem~12.6.
\end{rsbox}

\begin{theorem}[Fibonacci Prime Index Alignment]
\begin{rsbox}
The first seven Fibonacci primes have indices equal to the first seven canonical RS
structural constants in ascending order:

\begin{center}
\begin{tabular}{ccc}
\hline
Fibonacci index $k$ & Fibonacci prime $F(k)$ & RS identification of $k$ \\
\hline
3  & $2 = G_{\mathrm{mg}}$   & $N$ (triadic period) \\
4  & $3 = N$                 & $N_E$ (emission count) \\
5  & $5 = \Delta$            & $\Delta$ (flux imbalance) \\
7  & $13 = \Pi_3$            & $\Sigma_{\mathrm{imb}}$ (boundary-active sum) \\
11 & $89 = n_G$              & $T_{\mathrm{vac}}$ (vacuum period) \\
13 & $233 = G$               & $\Pi_3$ (Gen-3 Gate Prime) \\
17 & $1597$                  & $PG(\Sigma_{\mathrm{imb}})$ (matter-gate prime) \\
\hline
\end{tabular}
\end{center}

Each of the seven indices is a canonical RS structural constant established independently
of the Fibonacci prime sequence. No index in the table requires reference to any integer
outside the established corpus constant set
$\{G_{\mathrm{mg}}, N, N_E, \Delta, \Sigma_{\mathrm{imb}}, T_{\mathrm{vac}}, \Pi_3,
PG(\Sigma_{\mathrm{imb}})\}$.
\end{rsbox}
\end{theorem}

\begin{observebox}
\noindent\textbf{Cross-document connections.}

Path~A of Theorem~12.6 gives $F_{17} = D_{\mathrm{gauge}}^3 - S_{\mathrm{int}} + 1$,
equivalently:
\[
D_{\mathrm{gauge}}^3 = F_{17} + S_{\mathrm{int}} - 1 = 1597 + 132 - 1 = 1728.
\]
The three dominant RS integers in the proton mass derivation are
$D_{\mathrm{gauge}}^3$, $S_{\mathrm{int}}$, and the vacuum unit. They are exactly the three
components of the Path~A derivation of $F_{17}$.

\medskip
\noindent\textbf{Consecutive Fibonacci structure.} Section~21 establishes that the
fractional anomaly denominator and the gravitational coupling depth are consecutive
Fibonacci numbers: $F(\tau_{10}) = 55$ and $F(T_{\mathrm{vac}}) = 89$. The matter-gate
Fibonacci prime extends this to prime-indexed consecutive steps:
\[
F(11) = 89, \quad F(13) = 233, \quad F(17) = 1597.
\]
The indices $11, 13, 17$ are consecutive primes and their associated Fibonacci integers $F_{11},F_{13},F_{17}$ serve jointly as the gravitational depth, the Generation Sector Invariant, and the matter-gate Fibonacci prime, respectively. Technically these are literal integer outputs of the RS Fibonacci operator at prime indices that, by their consecutive prime indexing, encode the RS sector hierarchy from gravity through generation to the composite scalar boundary. Any mapping of these RS integers to continuum quantities or numerical constants is performed exclusively inside a clearly labeled Verification Path calculation that converts the RS integers into real number form for external assessment, thereby preserving the ontological layer as strictly integer based.


\medskip
\noindent\textbf{Denominator identity.} The integer
$S_{\mathrm{int}} - \tau_{10} = 132 - 10 = 122 = G_{\mathrm{mg}} \times PG(VS)$,
where $PG(VS) = PG(18) = 61$. Consequently:
\[
T_{\mathrm{vac}} \times (S_{\mathrm{int}} - \tau_{10}) = 11 \times 122 = 1342
= G_{\mathrm{mg}} \times T_{\mathrm{vac}} \times PG(VS).
\]
This is the RS-native structural reading of the Arrow-contracted interaction surface
scaled by the vacuum period, requiring no constants beyond those already established
in the corpus.
\end{observebox}

\begin{theorem}[Full Lattice Synchronisation Period]\label{thm:sync-period}
\begin{rsbox}
The integer $23760$ is the smallest positive integer divisible by every
RS structural period.  Divisibility is verified by constructive remainder:
\[
23760 \;\koppa\; N_E^2 = 23760 \;\koppa\; 16 = 0,\quad
23760 \;\koppa\; (\Delta \cdot T_{\mathrm{vac}}) = 23760 \;\koppa\; 55 = 0,
\]
\[
23760 \;\koppa\; (G_{\mathrm{mg}} \cdot N^3 \cdot T_{\mathrm{vac}}) = 23760 \;\koppa\; 594 = 0,\quad
23760 \;\koppa\; 4 = 0,\quad
23760 \;\koppa\; T_{\mathrm{vac}} = 0.
\]
\[
23760 = 180 \cdot \Sint.
\]
All lattice processes, $\psi$-cycles, vacuum cycles, and all three neutrino
mixing channels, return simultaneously to their initial phase at
$23760$ steps.  The lattice is fully coherent on this timescale.
\end{rsbox}
\end{theorem}

% ============================================================
% RS SOMMERFELD DERIVATION
% ============================================================

\begin{definition}[RS Sommerfeld Prime Pair]\label{def:sommerfeld-pair}
The RS Sommerfeld Prime Pair is the ordered pair of integers $(\pa,\qa)$ with
$\pa=115331$ and $\qa=15804499$.  Both integers are prime, established by the
$\koppa$-primality test of Definition~\ref{def:primegate}: $\pa\koppa m\neq 0$
for all integers $m$ with $1<m<\pa$, and $\qa\koppa m\neq 0$ for all integers
$m$ with $1<m<\qa$.
\end{definition}

\begin{definition}[RS Sommerfeld ERP]\label{def:sommerfeld-erp}
The RS Sommerfeld ERP is the explicit rational pair $(\pa,\qa)\in\SL$ with
numerator $\pa=115331$ and denominator $\qa=15804499$.  It is never reduced.
\end{definition}

\begin{definition}[CF Descent Sequence $\Da$]\label{def:Da}
The continued fraction descent sequence of the RS Sommerfeld ERP is the
ordered tuple of RS structural constants
\begin{equation}
\Da := \bigl[\aSC;\;N^3,\;1,\;N,\;\dslip,\;N\times\Sigimb,\;
1,\;1,\;\tauX\bigr]
= \bigl[137;\;27,\;1,\;3,\;2,\;21,\;1,\;1,\;10\bigr].
\label{eq:Da}
\end{equation}
Every element of $\Da$ is an RS structural constant established in
Section~\ref{sec:symboltable}.  No element is introduced for the purpose of
this definition.
\end{definition}

\begin{definition}[CF Convergent Recurrence]\label{def:cf-recurrence}
For a continued fraction $[a_0;\,a_1,\ldots,a_n]$, the convergent denominator
recurrence has seed $k_{-1}:=0$, $k_0:=1$ and step rule
\[
k_j := a_j\times k_{j-1}+k_{j-2},\qquad j\geq 1.
\]
The convergent numerator recurrence has seed $h_{-1}:=1$, $h_0:=a_0$ and the
same step rule
\[
h_j := a_j\times h_{j-1}+h_{j-2},\qquad j\geq 1.
\]
Every operation is integer multiplication and integer addition.  No division
appears.
\end{definition}

\begin{theorem}[Forward Generation of the RS Sommerfeld Prime Pair]
\label{thm:sommerfeld-forward}
\begin{rsbox}
The convergent denominator recurrence of Definition~\ref{def:cf-recurrence}
applied to $\Da$ with seed $k_{-1}=0$, $k_0=1$ produces $k_8=\pa=115331$.
The convergent numerator recurrence applied to $\Da$ with seed $h_{-1}=1$,
$h_0=137$ produces $h_8=\qa=15804499$.
\end{rsbox}
\begin{proof}
\begin{rsbox}
The denominator recurrence with partial quotients
$(a_1,\ldots,a_8)=(27,1,3,2,21,1,1,10)$:
\begin{align*}
k_0&=1,\quad k_1=27\times1+0=27,\quad k_2=1\times27+1=28,\\
k_3&=3\times28+27=111,\quad k_4=2\times111+28=250,\\
k_5&=21\times250+111=5250+111=5361,\\
k_6&=1\times5361+250=5611,\quad k_7=1\times5611+5361=10972,\\
k_8&=10\times10972+5611=109720+5611=115331=\pa.
\end{align*}
The numerator recurrence with $h_0=\aSC=137$:
\begin{align*}
h_1&=27\times137+1=3699+1=3700,\quad h_2=1\times3700+137=3837,\\
h_3&=3\times3837+3700=11511+3700=15211,\\
h_4&=2\times15211+3837=30422+3837=34259,\\
h_5&=21\times34259+15211=719439+15211=734650,\\
h_6&=1\times734650+34259=768909,\\
h_7&=1\times768909+734650=1503559,\\
h_8&=10\times1503559+768909=15035590+768909=15804499=\qa.
\end{align*}
No operation beyond integer multiplication and addition appears.
\end{rsbox}
\end{proof}
\end{theorem}

\begin{theorem}[CF Descent of the RS Sommerfeld ERP]
\label{thm:sommerfeld-descent}
\begin{rsbox}
The Euclidean algorithm applied to $(\qa,\pa)=(15804499,115331)$ using only
$\koppa$, integer comparison, and $\succ$ produces the continued fraction
$\Da=[137;\,27,\,1,\,3,\,2,\,21,\,1,\,1,\,10]$.  Every partial quotient is
an RS structural constant.  The descent terminates at $\tauX=10$.
\end{rsbox}
\begin{proof}
\begin{rsbox}
Nine steps, each quotient identified:

\textit{Step~1.}  $15804499=137\times115331+4152$.
Verify: $100\times115331=11533100$; $37\times115331=4267247$;
sum $=15800347$.  Remainder $=4152$.  Quotient $a_0=137=\aSC$.

\textit{Step~2.}  $115331=27\times4152+3227$.
Verify: $27\times4152=112104$.  Remainder $=3227$.  Quotient $a_1=27=N^3$.

\textit{Step~3.}  $4152=1\times3227+925$.
Remainder $=925$.  Quotient $a_2=1$ (vacuum unit).

\textit{Step~4.}  $3227=3\times925+452$.
Verify: $3\times925=2775$.  Remainder $=452$.  Quotient $a_3=3=N$.

\textit{Step~5.}  $925=2\times452+21$.
Verify: $2\times452=904$.  Remainder $=21$.  Quotient $a_4=2=\dslip$.

\textit{Step~6.}  $452=21\times21+11$.
Verify: $21\times21=441$.  Remainder $=11=\Tvac$.
Quotient $a_5=21=N\times\Sigimb$.

\textit{Step~7.}  $21=1\times11+10$.
Remainder $=10=\tauX$.  Quotient $a_6=1$ (vacuum unit).

\textit{Step~8.}  $11=1\times10+1$.
Remainder $=1$.  Quotient $a_7=1$ (vacuum unit).

\textit{Step~9.}  $10=10\times1+0$.
Quotient $a_8=10=\tauX$.  Descent terminates.

Quotients in order: $137,27,1,3,2,21,1,1,10$ identified as
$\aSC,N^3,1,N,\dslip,N\times\Sigimb,1,1,\tauX$.
\end{rsbox}
\end{proof}
\end{theorem}

\begin{theorem}[Near-Miss Structural Property of the First Remainder]
\label{thm:nearmiss}
\begin{rsbox}
The remainder of Step~1 of the CF descent satisfies $4152=4153-1$, where
$4153$ is prime.  By the Frictionless Propagation Principle
(Postulate~\ref{post:frictionless}), the prime-capacity state $4153$ cannot
discharge structural load into any sub-period oscillator.  The remainder
$4152$ sits exactly one vacuum unit below this frictionless boundary.
\end{rsbox}
\begin{proof}
\begin{rsbox}
Primality of $4153$: $4153\koppa m\neq 0$ for all $m$ with $1<m<4153$,
confirmed by the $\koppa$-primality test.
The subtraction $4153-4152=1$ is a single integer operation.
\end{rsbox}
\end{proof}
\end{theorem}

\begin{theorem}[Phi-Oscillator Near-Miss Connection]
\label{thm:phiosc-nearmiss}
\begin{rsbox}
\[
4621 - 4153 = N_E\times N^2\times\Pithree = 4\times9\times13 = 468.
\]
\end{rsbox}
\begin{proof}
\begin{rsbox}
$4\times9=36$; $36\times13=468$; $4621-468=4153$.
Primality of $4153$ confirmed in Theorem~\ref{thm:nearmiss}.
\end{rsbox}
\end{proof}
\end{theorem}

\begin{theorem}[Three-Cycle Slip Structure and Mass Formation Boundary]
\label{thm:slip-structure}
\begin{rsbox}
Over three consecutive $\Tvac$-cycles spanning one complete barycentric period
$\Tbary=33$ microticks, the net phase surplus accumulated per cycle under the
4-4-3 partition is uniformly $+1$ per cycle.  The remaining unresolved closure
after each cycle, measured against the total triadic closure $N=3$, is:

\begin{center}
\begin{tabular}{lcc}
\toprule
Cycle & Microticks & Remaining closure \\
\midrule
1 & $1$--$11$ & $N-1=2$ \\
2 & $12$--$22$ & $N-2=1$ \\
3 & $23$--$33$ & $N-3=0$ \\
\bottomrule
\end{tabular}
\end{center}

The slip sequence $(2,1,0)$ is the remaining unresolved closure per cycle.
The mass formation boundary is the transition from microtick~22 to
microtick~23: at microtick~22 the remaining closure is $1$ and temporal
charge $\Tcharge>0$ is present; at microtick~23 the final closing cycle
begins and is committed to completing the barycentric closure.  The boundary
cannot be placed at a fractional microtick position.
\end{rsbox}
\begin{proof}
\begin{rsbox}
The 4-4-3 partition assigns emission ticks $\{1,4,7,10\}$ and return ticks
$\{3,6,9\}$ within each $\Tvac$-cycle.

\textit{Cycle~1, microticks $1$--$11$:}
\begin{center}
\small
\begin{tabular}{ccccc}
\toprule
Tick & Phase & Cumul.\ E & Cumul.\ R & Surplus \\
\midrule
1&E&1&0&$+1$\\2&M&1&0&$+1$\\3&R&1&1&$0$\\
4&E&2&1&$+1$\\5&M&2&1&$+1$\\6&R&2&2&$0$\\
7&E&3&2&$+1$\\8&M&3&2&$+1$\\9&R&3&3&$0$\\
10&E&4&3&$+1$ (Arrow locks)\\11&M&4&3&$+1$\\
\bottomrule
\end{tabular}
\end{center}
Per-cycle surplus $+1$.  Remaining closure: $N-1=2$.

\textit{Cycle~2, microticks $12$--$22$} (emission $\{12,15,18,21\}$,
return $\{14,17,20\}$):
\begin{center}
\small
\begin{tabular}{ccccc}
\toprule
Tick & Phase & Cumul.\ E & Cumul.\ R & Surplus \\
\midrule
12&E&5&3&$+2$\\13&M&5&3&$+2$\\14&R&5&4&$+1$\\
15&E&6&4&$+2$\\16&M&6&4&$+2$\\17&R&6&5&$+1$\\
18&E&7&5&$+2$\\19&M&7&5&$+2$\\20&R&7&6&$+1$\\
21&E&8&6&$+2$\\22&M&8&6&$+2$\\
\bottomrule
\end{tabular}
\end{center}
Per-cycle contribution $+1$.  Remaining closure: $N-2=1$.  At microtick~22
surplus is $+2>0$: $\Tcharge>0$ and massive structure is admissible.

\textit{Cycle~3, microticks $23$--$33$} (emission $\{23,26,29,32\}$,
return $\{25,28,31\}$):
\begin{center}
\small
\begin{tabular}{ccccc}
\toprule
Tick & Phase & Cumul.\ E & Cumul.\ R & Surplus \\
\midrule
23&E&9&6&$+3$\\24&M&9&6&$+3$\\25&R&9&7&$+2$\\
26&E&10&7&$+3$\\27&M&10&7&$+3$\\28&R&10&8&$+2$\\
29&E&11&8&$+3$\\30&M&11&8&$+3$\\31&R&11&9&$+2$\\
32&E&12&9&$+3$\\33&M&12&9&$+3$\\
\bottomrule
\end{tabular}
\end{center}
Per-cycle contribution $+1$.  Total at tick~33: $+3=N$.  Remaining closure:
$N-3=0$.  Cycle~3 begins at tick~23 with its full closing obligation ahead;
from tick~23 onward the cycle is committed to completing the final unit of
closure and no capacity for new mass formation exists.
\end{rsbox}
\end{proof}
\end{theorem}

\begin{theorem}[Pell Attractor for the Mass Formation Boundary]
\label{thm:pell-23}
\begin{rsbox}
The Pell equation $a^2-23b^2=1$ has fundamental solution
$(a,b)=(24,5)=(2\Dg,\,\Delta)$.
\end{rsbox}
\begin{proof}
\begin{rsbox}
$b=1$: $a^2=24$, not a perfect square.
$b=2$: $a^2=93$, not a perfect square.
$b=3$: $a^2=208$, not a perfect square.
$b=4$: $a^2=369$, not a perfect square.
$b=5$: $a^2=1+23\times25=576=24^2$.
Fundamental solution $(24,5)=(2\Dg,\Delta)$.
\end{rsbox}
\end{proof}
\end{theorem}

\begin{theorem}[Koppa Signature Profile of the RS Sommerfeld Prime Pair]
\label{thm:sommerfeld-koppa}
\begin{rsbox}
\begin{center}
\begin{tabular}{lcc}
\toprule
Period $P$ & $\pa\koppa P$ & $\qa\koppa P$ \\
\midrule
$\Tvac=11$ & $7=\Sigimb$ & $7=\Sigimb$ \\
$\Tbary=33$ & $29=\PG(\tauX)$ & $7=\Sigimb$ \\
$\Delta\times\Tvac=55$ & $51$ & $29=\PG(\tauX)$ \\
$\Dg=12$ & $11=\Tvac$ & $7=\Sigimb$ \\
\bottomrule
\end{tabular}
\end{center}
Both primes carry $\Sigimb=7$ against $\Tvac$.  The denominator $\qa$ carries
$\Sigimb$ at every computed period.
\end{rsbox}
\begin{proof}
\begin{rsbox}
$\pa\koppa11$: $115331=10484\times11+7$.
Verify: $10484\times11=115324$.  Remainder $=7=\Sigimb$.

$\qa\koppa11$: $15804499=1436772\times11+7$.
Verify: $1436772\times11=15804492$.  Remainder $=7=\Sigimb$.

$\pa\koppa33$: $115331=3494\times33+29$.
Verify: $3494\times33=115302$.  Remainder $=29=\PG(\tauX)$.

$\qa\koppa33$: $15804499=478924\times33+7$.
Verify: $478924\times33=15804492$.  Remainder $=7=\Sigimb$.

$\pa\koppa55$: $115331=2096\times55+51$.
Verify: $2096\times55=115280$.  Remainder $=51$.

$\qa\koppa55$: $15804499=287354\times55+29$.
Verify: $287354\times55=15804470$.  Remainder $=29=\PG(\tauX)$.

$\pa\koppa12$: $115331=9610\times12+11$.
Verify: $9610\times12=115320$.  Remainder $=11=\Tvac$.

$\qa\koppa12$: $15804499=1317041\times12+7$.
Verify: $1317041\times12=15804492$.  Remainder $=7=\Sigimb$.
\end{rsbox}
\end{proof}
\end{theorem}

\begin{theorem}[The RS Sommerfeld Theorem]\label{thm:sommerfeld}
\begin{rsbox}
The RS Sommerfeld ERP $(\pa,\qa)=(115331,15804499)$ is the RS-native
derivation of the fine structure constant.  The following properties hold
simultaneously, each established by integer arithmetic alone from master
reference constants.

\emph{(i) Forward generation.}  Both primes are generated by the convergent
recurrence from $\Da$ (Theorem~\ref{thm:sommerfeld-forward}).

\emph{(ii) Backward descent.}  The Euclidean descent of $(\qa,\pa)$ recovers
$\Da$ exactly, every partial quotient an RS structural constant, terminating at
$\tauX$ (Theorem~\ref{thm:sommerfeld-descent}).

\emph{(iii) Sommerfeld Constant as leading partial quotient.}  The integer
part of $\qa/\pa$ is $137=\aSC$.  The Sommerfeld Constant is the first
structural constant revealed by the descent, not a derivation target.

\emph{(iv) Near-miss structure.}  The first descent remainder satisfies
$4152=4153-1$ where $4153$ is prime, and $4621-4153=N_E\times
N^2\times\Pithree=468$ (Theorems~\ref{thm:nearmiss}
and~\ref{thm:phiosc-nearmiss}).

\emph{(v) Inertial closure terminus.}  The descent terminates at $\tauX=10$
by the Euclidean structure of the prime pair.

\emph{(vi) Boundary-active synchronisation.}  Both primes carry $\Sigimb$
against $\Tvac$ (Theorem~\ref{thm:sommerfeld-koppa}).
\end{rsbox}
\begin{proof}
Properties (i) through (vi) follow from
Theorems~\ref{thm:sommerfeld-forward} through~\ref{thm:sommerfeld-koppa}
respectively.
\end{proof}
\end{theorem}

\begin{verifybox}
CODATA 2018 fine structure constant:
$\alpha=0.0072973525643$, uncertainty $1.1\times10^{-10}$.

RS Sommerfeld ERP ratio: $115331/15804499=0.0072973524817\ldots$

Deviation: $|0.0072973525643-0.0072973524817|=8.26\times10^{-11}$.

This lies within the CODATA uncertainty.  The integer part of
$\qa/\pa=15804499/115331=137.036\ldots$ confirms $\aSC=137$ as the leading
partial quotient.
\end{verifybox}

\begin{corollary}[Mass Gap at the Sommerfeld Terminus]
\label{cor:massgap-sommerfeld}
\begin{rsbox}
The CF descent of the RS Sommerfeld ERP terminates at $\tauX$, the last
microtick within the admissible zone for mass formation established in
Theorem~\ref{thm:slip-structure}.  The descent approaches the mass gap through
the near-miss structure of each remainder and exhausts itself at the boundary
of the mass-admissible zone.  The structural identity
$\tauX+\dslip=\Dg$ (Theorem~\ref{thm:tau10}) encodes the closure.  The mass
gap is the terminus condition of the Euclidean structure of the prime pair.
\end{rsbox}
\end{corollary}

\begin{theorem}[RS Closure of the Euclidean Descent]
\label{thm:descent-closed}
\begin{rsbox}
Every quotient and every remainder produced by the Euclidean descent of
$(\qa,\pa)=(15804499,115331)$ is expressible as a product of RS structural
constants.  The descent is RS-closed: no integer appearing at any
intermediate step requires reference to any object outside the RS corpus.

\begin{center}
\small
\begin{tabular}{lllll}
\toprule
Step & Quotient & RS identification & Remainder & RS identification \\
\midrule
1 & $137$ & $\aSC=\Sint+\Delta$
  & $4152$ & $\Gmg^3\cdot N\cdot\PG(N_E\cdot\tauX)$ \\
2 & $27$ & $N^3$
  & $3227$ & $\Sigimb\cdot\PG(\nG)$ \\
3 & $1$ & vacuum unit
  & $925$ & $\Ksat\cdot\PG(\Dg)$ \\
4 & $3$ & $N$
  & $452$ & $N_E\cdot\mathrm{denom}(\piRS)$ \\
5 & $2$ & $\dslip$
  & $21$ & $N\cdot\Sigimb$ \\
6 & $21$ & $N\cdot\Sigimb$
  & $11$ & $\Tvac$ \\
7 & $1$ & vacuum unit
  & $10$ & $\tauX$ \\
8 & $1$ & vacuum unit
  & $1$ & vacuum unit \\
9 & $10$ & $\tauX$
  & $0$ & terminus \\
\bottomrule
\end{tabular}
\end{center}

where $\mathrm{denom}(\piRS)=113$ is the denominator of the RS $\pi$-analog ERP
(Definition~\ref{def:pi-analog}), $\PG(\nG)=\PG(89)=461$ is the prime at
the gravitational coupling depth index, $\PG(N_E\cdot\tauX)=\PG(40)=173$,
and $\PG(\Dg)=\PG(12)=37$.
\end{rsbox}
\begin{proof}
\begin{rsbox}
Each step is verified by integer multiplication and subtraction alone.

\textit{Step~1.} $137\times115331=15800347$.
$15804499-15800347=4152$.
$4152=8\times3\times173=\Gmg^3\times N\times173$.
$N_E\times\tauX=4\times10=40$.
$\PG(40)=173$: $173\koppa m\neq0$ for all $m$ with $1<m<173$.

\textit{Step~2.} $27\times4152=112104$.
$115331-112104=3227$.
$3227=7\times461=\Sigimb\times461$.
$\PG(89)=461$: $461\koppa m\neq0$ for all $m$ with $1<m<461$.
$\nG=89$ (gravitational coupling depth, master reference).

\textit{Step~3.} $1\times3227=3227$.
$4152-3227=925$.
$925=25\times37=\Ksat\times37$.
$\PG(12)=37$: $37\koppa m\neq0$ for all $m$ with $1<m<37$.
$\Dg=12$.

\textit{Step~4.} $3\times925=2775$.
$3227-2775=452$.
$452=4\times113=N_E\times113$.
$\mathrm{denom}(\piRS)=113$ by Definition~\ref{def:pi-analog}.

\textit{Step~5.} $2\times452=904$.
$925-904=21=3\times7=N\times\Sigimb$.

\textit{Step~6.} $21\times21=441$.
$452-441=11=\Tvac$.

\textit{Step~7.} $1\times11=11$.
$21-11=10=\tauX$.

\textit{Step~8.} $1\times10=10$.
$11-10=1$.

\textit{Step~9.} $10\times1=10$.
$10-10=0$.

Every quotient and remainder is an RS structural integer.
No division appears at any step.
\end{rsbox}
\end{proof}
\end{theorem}

\begin{observebox}
The RS-closure of the Euclidean descent reveals the complete structural
architecture encoded in the RS Sommerfeld Prime Pair.  The remainders carry
a descending hierarchy of RS sector constants: $r_1=\Gmg^3\cdot
N\cdot\PG(N_E\cdot\tauX)$ encodes the charge and phase structure; $r_2=
\Sigimb\cdot\PG(\nG)$ encodes the boundary-active phase sum and the
gravitational sector; $r_3=\Ksat\cdot\PG(\Dg)$ encodes the saturation
limit and gauge dimension; $r_4=N_E\cdot\mathrm{denom}(\piRS)$ encodes the emission
cardinality and the RS $\pi$-analog denominator; and $r_5$ through $r_8$
recover $N\cdot\Sigimb$, $\Tvac$, $\tauX$, and the vacuum unit in sequence.
The gravitational sector ($\nG=89$), the electromagnetic sector ($\Sint$,
$\Delta$), the $\pi$-analog structure ($\mathrm{denom}(\piRS)=113$), and the phase
partition ($N$, $\Sigimb$, $\tauX$, $\Tvac$) are all present in the descent.
The RS Sommerfeld Prime Pair encodes the complete RS vacuum architecture in
a single Euclidean descent.
\end{observebox}

% ============================================================
% RS PI-ANALOG ERP
% ============================================================

\begin{definition}[RS $\pi$-Analog ERP]\label{def:pi-analog}
The RS $\pi$-analog ERP is the explicit rational pair $(355,113)\in\SL$.  It
is the structural representative of the $W=1$ bracket reached by the mediant
tree path toward $\pi$ after four legs consuming the partial quotients $N$,
$\Sigimb$, $N\times\Delta$, and the vacuum unit in sequence.  It is never
reduced.

The four convergent pairs, each adjacent with $W=1$:
\begin{align*}
\Dcross\!\bigl((3,1),(22,7)\bigr)
  &= 3\times7-22\times1 = 21-22 = -1,\\
\Dcross\!\bigl((22,7),(333,106)\bigr)
  &= 22\times106-333\times7 = 2332-2331 = 1,\\
\Dcross\!\bigl((333,106),(355,113)\bigr)
  &= 333\times113-355\times106 = 37629-37630 = -1.
\end{align*}

The structural representative $(355,113)$ is selected as the upper bound of
the final $W=1$ bracket by the parity of the imbalance ledger at $\Omega$ per
Definition~\ref{def:bary-trunc}.

The total koppa load across the four legs is:
\[
N+\Sigimb+N\times\Delta+1 = 3+7+15+1 = 26 = \Gmg\times\Pithree.
\]
\end{definition}

\begin{theorem}[Koppa Signature Profile of the RS $\pi$-Analog ERP]
\label{thm:pi-koppa}
\begin{rsbox}
\begin{center}
\begin{tabular}{lcc}
\toprule
Period $P$ & $355\koppa P$ & $113\koppa P$ \\
\midrule
$\Tvac=11$ & $3=N$ & $3=N$ \\
$\Tbary=33$ & $25=\Ksat$ & $14=\Gmg\times\Sigimb$ \\
$\Dg=12$ & $7=\Sigimb$ & $5=\Delta$ \\
$\Delta\times\Tvac=55$ & $25=\Ksat$ & $3=N$ \\
$\Sint=132$ & $91=\Sigimb\times\Pithree$ & $113$ \\
\bottomrule
\end{tabular}
\end{center}
Both integers carry $N=3$ against $\Tvac$.  The numerator $355$ carries
$\Ksat$ against $\Tbary$ and $\Delta\times\Tvac$ simultaneously.  The
denominator $113$ carries $\Delta$ against $\Dg$.
\end{rsbox}
\begin{proof}
\begin{rsbox}
$355\koppa11$: $355=32\times11+3$.  Verify: $32\times11=352$.
Remainder $=3=N$.

$113\koppa11$: $113=10\times11+3$.  Verify: $10\times11=110$.
Remainder $=3=N$.

$355\koppa33$: $355=10\times33+25$.  Verify: $10\times33=330$.
Remainder $=25=\Ksat$.

$113\koppa33$: $113=3\times33+14$.  Verify: $3\times33=99$.
Remainder $=14=\Gmg\times\Sigimb$.

$355\koppa12$: $355=29\times12+7$.  Verify: $29\times12=348$.
Remainder $=7=\Sigimb$.

$113\koppa12$: $113=9\times12+5$.  Verify: $9\times12=108$.
Remainder $=5=\Delta$.

$355\koppa55$: $355=6\times55+25$.  Verify: $6\times55=330$.
Remainder $=25=\Ksat$.

$113\koppa55$: $113=2\times55+3$.  Verify: $2\times55=110$.
Remainder $=3=N$.

$355\koppa132$: $355=2\times132+91$.  Verify: $2\times132=264$.
Remainder $=91=7\times13=\Sigimb\times\Pithree$.

$113\koppa132$: $113<132$.  Remainder $=113$.
\end{rsbox}
\end{proof}
\end{theorem}

\begin{observebox}
The total koppa load across the four legs of the RS $\pi$-analog construction
is $\Gmg\times\Pithree=26$.  This is the same integer that separates the
Generation Sector Invariant from the first quark generation mass integer:
$\mathcal{G}-\Lambda_{q1}=233-207=26$.  The RS $\pi$-analog and the
generation sector are coupled through this integer.  The Koide deviation from
$\varepsilon=2/3$ is structurally identified as the $\pi$-remainder after
barycentric truncation at $(355,113)$, scaled by $\Gmg\times\Pithree$.
\end{observebox}

\newpage


\section{The Phi-Oscillator and Lepton Hierarchy}\label{sec:phi-oscillator}

\subsection{The Path from $(11,7)$ to $(7477,4621)$}

\begin{theorem}[Prime-Prime $\phi$-Approximant Theorem]\label{thm:phi-prime}
\begin{rsbox}
The unique minimal prime-prime rational approximant to $\phi$ accessible from
the Lucas seed $(T_{\mathrm{vac}},\Sigimb)=(11,7)=(L(5),L(4))$ via the
Mediant Tree is the ERP $(7477,4621)$.  Its continued fraction is:
\[
\mathrm{CF}(7477/4621)
= [\underbrace{1,1,1,1,1,1,1,1,1,1,1,1}_{\Dg=12},\;\Delta,\;1,\;N_E]
= [1^{12},\;5,\;1,\;4].
\]
The partial quotient sequence encodes $\Dg$, $\Delta$, and $N_E$ in
sequential order.  Every consecutive pair of convergents in this path has
$\Dcross=\pm 1$, confirming $W=1$ adjacency throughout.
\end{rsbox}
\begin{proof}
The continued fraction of $7477/4621$ is verified by the $\aop^{-1}$
descent $\aop^{-1}(n,d)=(d,n-d)$.  Adjacency of each consecutive convergent
pair $(p_{k-1}/q_{k-1},\;p_k/q_k)$ requires
$\Dcross((p_{k-1},q_{k-1}),(p_k,q_k))=\pm 1$, verified for all 14 convergent
pairs.  Primality of both $7477$ and $4621$ is confirmed by the
$\koppa$-primality test (Definition~\ref{def:primegate}).
\end{proof}
\end{theorem}

\begin{observebox}
\textbf{Structural reading of the continued fraction.}
The $\Dg=12$ leading ones correspond to $\Dg$ Fibonacci steps along the
$\phi$-convergent path. The pure Fibonacci geodesic that the vacuum follows
for exactly as many steps as the gauge manifold dimension demands.  The
$\Delta=5$ excursion kicks the path off the Fibonacci geodesic into the
interior of the Mediant Tree, driven by the flux imbalance.  The $N_E=4$
closure steps are the Emission phase count, which precipitates the
prime-prime pair and saturates the bracket at $W=1$.
\end{observebox}

\subsection{Properties of the Phi-Oscillator ERP}

\begin{theorem}[Primorial Property of the Denominator]\label{thm:primorial}
\begin{rsbox}
\[
4621 = 2\times\#(T_{\mathrm{vac}})+1
     = 2\times(2\cdot3\cdot5\cdot7\cdot11)+1
     = 2\times 2310+1.
\]
Consequently $4621\equiv 1\pmod{p}$ for every prime $p\leq T_{\mathrm{vac}}$.
A state of capacity $4621$ carries exactly $+1$ unresolved unit of structural
history against every fundamental RS cycle simultaneously.  This is
minimal uniform structural tension: the state closest to
frictionless against all RS cycles simultaneously without being degenerate.
\end{rsbox}
\end{theorem}

\begin{theorem}[Threshold Sphenic Integer]\label{thm:threshold}
\begin{rsbox}
The convergent of $7477/4621$ at depth $\Dg$. This is the position just before the $\Delta$-excursion at $(1309,809)$, where:
\[
1309 = \Sigimb\times T_{\mathrm{vac}}\times(\Dg+\Delta)
     = 7\times 11\times 17.
\]
This is the product of the three second-tier RS constants. It is the threshold state. The lattice must transit $1309$ before a prime-prime approximant to $\phi$ can emerge. Both the numerator and denominator of such an approximant must resolve their imbalance against all three second-tier constants before achieving simultaneous primality.
\end{rsbox}
\end{theorem}

\subsection{The Phi-Oscillator as Lepton Hierarchy Generator}

\begin{theorem}[Penultimate Descent State]\label{thm:descent-state}
\begin{rsbox}
Applying $\aop^{-1}(n,d)=(d,n-d)$ iteratively from $(7477,4621)$, the state
at step~12 is $(29,5)=(29,\Delta)$.  The integer $29$ satisfies:
\[
29 = \PG(T_{\mathrm{vac}}-1) = \PG(\tau_{10}) = \PG(10), \qquad
29 = \Tbary-N_E = 33-4.
\]
Both identifications are exact.  The index $T_{\mathrm{vac}}-1=10$ is the
inertial closure tick $\tau_{10}$ (Theorem~\ref{thm:tau10}).
\end{rsbox}
\end{theorem}

\begin{theorem}[Phi-Oscillator Paths to Lepton Mass Integers]
\label{thm:phi-lepton}
\begin{rsbox}
From the penultimate descent state
$(\PG(\tau_{10}),\Delta)=(29,5)$:\footnote{These are additional independent
derivation paths for $\Lambda_\mu=207$ and $\Lambda_\tau=3481$; see also
Appendix~\ref{app:lmu-phiosc} and Appendix~\ref{app:ltau-phiosc}.}
\[
\Lambda_\mu = \PG(\tau_{10})\times\Sigimb+N_E = 29\times 7+4 = 207,
\]
\[
\Lambda_\tau = \PG(\tau_{10})\times(\Sint-T_{\mathrm{vac}}-1)+1
             = 29\times 120+1 = 3481,
\]
where $\Sint-T_{\mathrm{vac}}-1=132-11-1=120=\delta_{\mathrm{slip}}^3
\times(T_{\mathrm{vac}}+N_E)=8\times 15$.
\end{rsbox}
\end{theorem}

\begin{observebox}
The phi-oscillator $(7477,4621)$ is the $\aop$-orbit from the Lucas seed
$(T_{\mathrm{vac}},\Sigimb)$ through exactly $\Dg=12$ Fibonacci convergent
steps, landing at the $(233,144)$ convergent ($144=\Dg^2$, $233$ the 6th
Fibonacci prime), then through the tail $[\Delta,1,N_E]$ to the $W=1$
saturated state.  The descent path carries the full lepton mass generation
sequence as intermediate products.  The phi-oscillator is the generator of
the lepton mass hierarchy expressed as a single $W=1$ barycentric trajectory.
\end{observebox}

\begin{theorem}[Barycentric Derivation of $\Sint$]\label{thm:sint-bary}
\begin{rsbox}
\[
(1,T_{\mathrm{vac}})\bplus(1,\Dg)
= (1,11)\bplus(1,12)
= (1\cdot 12+1\cdot 11,\;11\cdot 12)
= (23,\;132)
= (\PG(N^2),\;\Sint).
\]
The interaction surface $\Sint=132$ is the barycentric denominator produced
when the vacuum seed couples to the unit-gauge state $(1,\Dg)$.  The
numerator $23=\PG(N^2)=\PG(9)$ is the Prime Gate of the triadic
square.\footnote{This provides a third independent derivation of $\Sint$,
independent of the multiplicative-period argument and additive-sum arguments of
Theorem~\ref{thm:sint-additive}.  See Appendix~\ref{app:sint-bary}.}

Adjoining: $(1,G_{\mathrm{mg}})\bplus(1,T_{\mathrm{vac}})
=(1,2)\bplus(1,11)=(13,22)$.  The Bridge denominator $22$ is the barycentric
product of the mass gap and the vacuum period.  The numerator $13$ is the
4th Fibonacci prime.
\end{rsbox}
\end{theorem}

\newpage


\section{Principal Theorems}\label{sec:theorems}

\begin{theorem}[4-4-3 Partition]\label{thm:443}
The sequential phase mapping $f:\Tcyc\to\PhiSet$ produces $N_E=4$, $N_M=4$,
$N_R=3$.  This partition is derived from the cyclic pointer structure of
Definition~\ref{def:443}, not postulated.
\end{theorem}

\begin{theorem}[Flux Imbalance $\Delta=5$]\label{thm:flux}
$\Delta:=(N_E+N_M)-N_R=5$.  Derived from the graph topology of the 4-4-3
partition; not assumed.
\end{theorem}

\begin{theorem}[Boundary Stress Generation]\label{thm:boundary-stress}
The Return Phase contraction ($N_R=3<4$) at cycle closure generates Boundary
Stress: the origin of restorative tension and the accumulation site for
$\koppa_{\mathrm{ledger}}$.
\end{theorem}

\begin{theorem}[Discrete Spin: $\psi^4=\mathrm{Id}$]\label{thm:spin}
The transformative reciprocal satisfies $\psi^4=\mathrm{Id}$; $\psi^2$ is
the full pair swap.  This is the ontic origin of the spinor property
($720^\circ$ identity).  The period~4 equals the Emission cardinality $N_E$
(Proposition~\ref{prop:psi-ne}).
\begin{proof}
Direct computation from Definition~\ref{def:rop}.  At each step the
denominator of each state crosses to become the numerator of the other,
and the numerator becomes the new denominator.
\begin{align*}
\psi^1\!\left(\tfrac{a}{b},\tfrac{c}{d}\right)
  &= \left(\tfrac{d}{a},\tfrac{b}{c}\right) \\
\psi^2\!\left(\tfrac{a}{b},\tfrac{c}{d}\right)
  &= \psi\!\left(\tfrac{d}{a},\tfrac{b}{c}\right)
   = \left(\tfrac{c}{d},\tfrac{a}{b}\right)
   \quad\text{(full pair exchange)} \\
\psi^3\!\left(\tfrac{a}{b},\tfrac{c}{d}\right)
  &= \psi\!\left(\tfrac{c}{d},\tfrac{a}{b}\right)
   = \left(\tfrac{b}{c},\tfrac{d}{a}\right) \\
\psi^4\!\left(\tfrac{a}{b},\tfrac{c}{d}\right)
  &= \psi\!\left(\tfrac{b}{c},\tfrac{d}{a}\right)
   = \left(\tfrac{a}{b},\tfrac{c}{d}\right)
   \quad\text{(identity confirmed)}
\end{align*}
\end{proof}
\end{theorem}

\begin{proposition}[$\psi$-Period equals Emission Cardinality]\label{prop:psi-ne}
\begin{rsbox}
The order of the Transformative Reciprocal equals the Emission cardinality:
\[
\psi\text{-period} = 4 = N_E.
\]
The Emission ticks $\{1, 4, 7, 10\} \subset \Tcyc$ number exactly $N_E = 4$
per vacuum cycle.  The $\psi$ operator, acting in 1D, completes one full
period in the same count as the vacuum completes one full set of Emission
events.  This is not assigned: it is what the cross-reciprocal structure
and the Prime-11 vacuum architecture produce together.

\medskip
\noindent\textbf{Consequence for mixing denominators:}
The atmospheric mixing denominator $N_E^2 = 16 = \psi\text{-period}^2$
is the square of the operator period, confirming that atmospheric mixing
operates at the intra-$\psi$-cycle scale.
\end{rsbox}
\end{proposition}

\begin{proposition}[Generations as Discrete Rest Mass Levels]\label{prop:generations-mass}
\begin{rsbox}
The number of particle generations $N = 3$ equals the structural prime-event depth produced by suspended state resolution accumulating through the Viscosity Sum:
\[
\rho_\Omega(\VS) = \rho_\Omega(18) = \rho_\Omega(G_{\mathrm{mg}} \cdot N^2) = 3 = N.
\]
$\VS = 18 = 2 \cdot 3^2$. Prime factorisation gives $\rho_\Omega(18) = 1 + 2 = 3 = N$. The three generations are the discrete rest mass levels available to a state that has accumulated the full viscosity sum through suspension resolution. The generation count is not postulated. It is the prime-event depth of the viscosity sum.
\end{rsbox}
\end{proposition}

\begin{theorem}[Viscosity Sum as Generation Boundary]\label{thm:vsbound}
\begin{rsbox}
Running the mass-gap cycle $(1/2,\, 2/1)$ for exactly $\VS = 18$ steps produces a koppa accumulation of $N^2 = 9$. The prime-event depth of the viscosity sum, $\rho_\Omega(\VS) = N = 3$, equals the number of particle generations, and this depth is achieved by exactly $\VS$ steps of the mass-gap cycle. No state deeper than generation 3 is structurally available within the $\VS$ synchronisation window.
\end{rsbox}
\begin{proof}
From Lemma~\ref{lem:koppa-cycle}, the mass-gap cycle accumulates koppa at a rate of $G_{\mathrm{mg}} = 2$ per four-step $\psi$-period, equivalently 1 unit every 2 steps. Over $\VS = 18$ steps, the number of subluminal (koppa-depositing) steps is $18 \times (2/4) = 9 = N^2$. Each deposits exactly 1 unit of koppa (since $1 \koppa 2 = 1$). Total koppa accumulated is $N^2 = 9$.

The generation count follows from $\rho_\Omega(\VS)$. The prime factorisation of $\VS = 18$ is $2 \cdot 3^2$, giving $\rho_\Omega(18) = 1 + 2 = 3 = N$ (Proposition~\ref{prop:generations-mass}).

The precise step-count derivation:
\[
\VS = G_{\mathrm{mg}} \cdot N^2 = 2 \cdot 9 = 18.
\]
$G_{\mathrm{mg}} = 2$ is the koppa per $\psi$-period. $N^2 = 9$ is the total koppa required to reach the $\VS$ boundary. Fewer than 18 steps does not complete the viscosity accumulation. After exactly 18 steps, the boundary is reached and $\rho_\Omega(\VS) = 3 = N$ encodes the full generation count.
\end{proof}
\end{theorem}
\begin{theorem}[Three Frictionless Generation Depths]\label{thm:frictionless}
\begin{rsbox}
The three generation mass integers $\mathcal{M}_1 = 1$, $\mathcal{M}_2 = 207$,
and $\mathcal{M}_3 = 3481$ are each frictionless against every RS structural
period.  For each $\mathcal{M}_k$ and every period $P \in
\{4,\, 11,\, 22,\, 33,\, 44,\, 55,\, 132\}$, the constructive remainder
$\mathcal{M}_k \koppa P \neq 0$.  The three integers are derived from RS
constants as follows:
\[
\mathcal{M}_1 = 1, \qquad
\mathcal{M}_2 = N^2 \cdot \PG(N^2) = 9 \cdot 23 = 207, \qquad
\mathcal{M}_3 = \PG(\Dg + \Delta)^2 = 59^2 = 3481.
\]
\end{rsbox}
\begin{proof}
For $\mathcal{M}_1 = 1$: $1 \koppa P = 1 \neq 0$ for every $P > 1$, since the
constructive remainder process subtracts $P$ from $r = 1$ zero times and
terminates immediately.

For $\mathcal{M}_2 = 207 = 3^2 \cdot 23$: $207 \koppa 4 = 3$; $207 \koppa 11 =
9$; $207 \koppa 22 = 9$; $207 \koppa 33 = 9$; $207 \koppa 44 = 31$;
$207 \koppa 55 = 42$; $207 \koppa 132 = 75$.  All nonzero.  Derivation:
$N^2 = 9$ is the koppa accumulation over $\VS$ steps (Theorem~\ref{thm:vsbound});
$\PG(N^2) = \PG(9) = 23$ is the 9th prime, the smallest prime that cannot
discharge into any of the $N^2$ sub-oscillations of the $\VS$ window.

For $\mathcal{M}_3 = 3481 = 59^2$: $3481 \koppa 4 = 1$; $3481 \koppa 11 = 5$;
$3481 \koppa 22 = 5$; $3481 \koppa 33 = 16$; $3481 \koppa 44 = 5$;
$3481 \koppa 55 = 16$; $3481 \koppa 132 = 49$.  All nonzero.  Derivation:
$\Dg + \Delta = 17$ is the matter-gate index; $\PG(17) = 59$ is the 17th prime;
$59^2 = 3481$ is the closed matter-gate depth.

All arithmetic uses only integer subtraction, consistent with Definition~\ref{def:koppa}.
\end{proof}
\end{theorem}

\begin{proposition}[Fourth Generation Structural Exclusion]\label{prop:gen4}
\begin{rsbox}
No fourth generation mass integer exists within the viscosity sum synchronisation window. The $\VS$ window supports exactly $N = 3$ generation levels, as encoded by $\rho_\Omega(\VS) = N$. Every RS gating constant available to the mass-gap cycle within $\VS$ steps is consumed by the three generations.
\end{rsbox}
\begin{proof}
From Theorem~\ref{thm:vsbound}, $\rho_\Omega(\VS) = N = 3$. The three generation depths consume the gating constants $N^2 = 9$ (for $\mathcal{M}_2$) and $\Dg + \Delta = 17$ (for $\mathcal{M}_3$). A fourth generation would require a fourth structurally independent gating constant reachable by the mass-gap cycle within $\VS$ steps and producing a frictionless depth by Postulate~\ref{post:frictionless}. The remaining RS constants reachable within $\VS$ steps ($T_{\mathrm{vac}} = 11$, $\Sint = 132$, $n_G = 89$) are either already encoded in the period structure consumed by $\mathcal{M}_2$ and $\mathcal{M}_3$, or belong to the long-range force sector rather than the lepton sector. No fourth independent gating constant is available. The $\VS$ window closes at three generation levels.
\end{proof}
\end{proposition}

\begin{theorem}[Matter Stability]\label{thm:matter-stability}
For $P_1,P_2\in\SPs$, the matter generator satisfies $G(P_1,P_2)\in\SPs$.
Active matter cannot produce degenerate output.
\begin{proof}
Case A: trivially $P_1\in\SPs$.  Case B: $Z_3=8S_B^3=8Y_1^3Z_1^3\neq 0$
since $Y_1,Z_1\neq 0$ for $P_1\in\SPs$; $X_3$ and $Y_3$ are integers by
construction.  Case C: $V=X_2Z_1-X_1Z_2\neq 0$ for projectively distinct
$P_1,P_2\in\SPs$, so $Z_3=V^3W\neq 0$.  In all three cases
$(X_3,Y_3,Z_3)\in\SPs$.
\end{proof}
\end{theorem}

\begin{theorem}[Arrow of Time via Doubling Generator]\label{thm:arrow-doubling}
Case B strictly increases structural potential: for all $(X_1,Y_1,Z_1)\in\SPs$,
$\rank(Z_3)>\rank(Z_1)$.
\begin{proof}
$Z_3=8S_B^3=8Y_1^3Z_1^3$.  Then
$\rank(Z_3)=3+3\rank(Y_1)+3\rank(Z_1)\geq 3+3\rank(Z_1)>\rank(Z_1)$
for all $Y_1\neq 0$.
\end{proof}
\end{theorem}

\begin{theorem}[The Rank-Temporal Boundary]\label{thm:rank-temporal}
\begin{rsbox}
Physical time is a phase-dependent property of geometric capacity.
\begin{itemize}[nosep]
\item \textbf{Rank-0 (Propagator Domain):} $d=1$.  The only propagating state
  is luminal ($n=d=1$), yielding $\Tcharge=0$.  Time does not exist at Rank-0.
\item \textbf{Rank-1 (Temporal Domain):} $d\geq 2$.  This domain permits
  subluminal states ($n<d$), yielding $\Tcharge>0$.
\end{itemize}
The Mass Gap ($d=2$) is the gate where the vacuum acquires capacity
to accumulate a temporal charge.  The minimum non-zero temporal charge is
$\Tcharge(1,2)=4-1=3=\Gamma_{\min}$, achieved at the unique minimal massive
state $(1,2)$.
\end{rsbox}
\end{theorem}

\begin{theorem}[Additive Synchronisation of Vacuum and Gauge]
\label{thm:sint-additive}
\begin{rsbox}
The temporal vacuum period $T_{\mathrm{vac}}=11$ and the gauge dimension
$\Dg=12$ share no common factor ($T_{\mathrm{vac}}=11$ is prime and does not divide $\Dg=12$).  The interaction
surface $\Sint$ is the unique minimal integer state-count required to map the
gauge manifold across the vacuum period by additive accumulation:
\begin{equation}\label{eq:sint}
\Sint = \sum_{i=1}^{T_{\mathrm{vac}}}\Dg
= \underbrace{12+12+\cdots+12}_{11}=132.
\end{equation}
This sum constitutes the bulk capacity required for the discrete vacuum to
scan the 12-dimensional manifold within the Prime-11 temporal container,
fulfilling the synchronisation condition without functional reduction or
modular primitives.\footnote{An additional independent derivation of
$\Sint=132$ via barycentric addition: $(1,T_{\mathrm{vac}})\bplus(1,\Dg)=(23,132)$
(Appendix~\ref{app:sint-bary}).}
\end{rsbox}
\end{theorem}

\begin{theorem}[Sommerfeld Constant]\label{thm:alpha-int}
\begin{rsbox}
\[
\alpha^{-1}_{\mathrm{int}}=\Sint+\Delta=132+5=137.
\]
The integer 137 is simultaneously: (i)~$\Sint+\Delta$; (ii)~the
$T_{\mathrm{bary}}$-th prime $\PG(33)$; (iii)~the unique minimal solution to
the gauge dimension uniqueness conditions (Theorem~\ref{thm:Dg-unique}).
\end{rsbox}
\end{theorem}

\begin{corollary}[Flux Signature of the Sommerfeld Constant]
\label{cor:flux-sig}
$\alpha^{-1}_{\mathrm{int}} \koppa 22 = \Delta$, since $137=6\times 22+5$ and $\Delta=5$. The same identity holds against $T_{\mathrm{vac}}=11$: $137 \koppa 11 = 5$. The integer 137 carries the flux imbalance signature $\Delta$ with respect to both the Bridge denominator and the vacuum period.
\end{corollary}

\begin{theorem}[Prime Gate Identities for Canonical Constants]
\label{thm:pg-identities}
\begin{rsbox}
Four structural identities link the Prime Gate to canonical RS constants:
\begin{align}
\PG(T_{\mathrm{bary}}) &= \PG(33) = 137
  = \alpha^{-1}_{\mathrm{int}} = \Sint+\Delta
  \label{eq:pg-alpha}\\[4pt]
\PG(\Dg+\Delta)^2 &= \PG(17)^2 = 59^2 = 3481
  = \Lambda_\tau
  \label{eq:pg-tau}\\[4pt]
N^2\times\PG(N^2) &= 9\times\PG(9) = 9\times 23 = 207
  = \Lambda_\mu
  \label{eq:pg-mu}\\[4pt]
\PG(\dim\,\mathrm{U}(1)) &= \PG(1) = 2 = G_{\mathrm{mg}}
  \label{eq:pg-mg}
\end{align}
Equation~\eqref{eq:pg-alpha} is a third independent derivation of
$\alpha^{-1}_{\mathrm{int}}$.\footnote{See Appendix~\ref{app:alpha-primegate}
for this pathway in full.}  Equation~\eqref{eq:pg-tau} confirms the tau mass
integer while reading the index $17=\Dg+\Delta$ as gauge dimension plus flux
imbalance.  Equation~\eqref{eq:pg-mu} is a self-referential identity: the
triadic period gates its own square through the $N^2$-th prime.
\end{rsbox}
\end{theorem}

\begin{theorem}[Dual Identity: Fractional Anomaly]\label{thm:dual-anomaly}
\begin{rsbox}
The fractional anomaly $\delta_A=2/55$ is derived by two independent paths:

\textbf{Path A (algebraic).}
\[
\delta_A = \frac{N^3-K_{\mathrm{sat}}}{\Delta\cdot T_{\mathrm{vac}}}
         = \frac{27-25}{55}
         = \frac{\delta_{\mathrm{slip}}}{\Delta\cdot T_{\mathrm{vac}}}
         = \frac{2}{55}.
\]
Every integer is RS-structural: $N^3=27$ (triadic cube),
$K_{\mathrm{sat}}=\Delta^2=25$ (saturation limit),
$\Delta\cdot T_{\mathrm{vac}}=55$ (flux-vacuum product).
The numerator equals $\delta_{\mathrm{slip}}=2$.

\textbf{Path B (geometric phase load).}  The full derivation appears in
TRTSv4 (Fractional Anomaly, geometric phase load derivation); both paths
yield $\delta_A=2/55$, establishing the Dual Identity.

The ERP representation is $\delta_A=(2,55)\in\SL$.
\end{rsbox}
\end{theorem}

\begin{theorem}[RS Inverse Fine-Structure Constant]\label{thm:alpha-rs}
\begin{rsbox}
\[
\alpha^{-1}_{RS} = \alpha^{-1}_{\mathrm{int}}+\delta_A
= 137+\frac{2}{55} = 137.\overline{036}.
\]
\end{rsbox}
\begin{verifybox}
$\alpha^{-1}_{\mathrm{exp}}\approx 137.036$; $\delta=0.000$.
Agreement to four significant figures; the repeating decimal
$0.\overline{036}=2/55$ is exact.
\end{verifybox}
\end{theorem}

\begin{theorem}[Proton-Electron Mass Ratio]\label{thm:proton-mass}
\begin{rsbox}
\[
\mu = \Dg^3+(\Sint-2\Dg)+4\delta_A
    = 1728+108+\frac{8}{55}
    = 1836+\frac{8}{55}.
\]
\end{rsbox}
\begin{verifybox}
$\mu_{\mathrm{exp}}\approx 1836.152$;
$1836+8/55=1836.1\overline{45}$; $\delta\approx 0.007$.
\end{verifybox}
\end{theorem}

\begin{theorem}[Neutron-Proton Mass Difference]\label{thm:neutron-proton}
\begin{rsbox}
\[
m_n-m_p = \left(\frac{\Delta}{G_{\mathrm{mg}}}+\delta_A\right)m_e
         = \left(\frac{5}{2}+\frac{2}{55}\right)m_e.
\]
\end{rsbox}
\end{theorem}

\begin{theorem}[Interval Invariance under Rational Boost]
\label{thm:interval-inv}
$I(E)$ is strictly invariant under $A_Q$; proved via the integer identity
$\alpha_{\mathrm{boost}}^2-\beta_{\mathrm{boost}}^2
=(d_u^2+n_u^2)^2-(2n_ud_u)^2=(d_u^2-n_u^2)^2=\Delta_{\mathrm{boost}}^2$.
Zero numerical drift.
\end{theorem}

\begin{theorem}[RS Time Dilation]\label{thm:time-dilation}
\begin{rsbox}
For state $(n,d)$ with velocity ratio $u=n/d$:
\[
\frac{\Tcharge}{d^2} = \frac{d^2-n^2}{d^2} = 1-u^2 = \frac{1}{\gamma^2}.
\]
The Lorentz factor emerges as the ratio $\Tcharge/d^2$, established directly
from the temporal charge structure of the ERP substrate.
\end{rsbox}
\end{theorem}

\begin{theorem}[Discriminant Mass Law]\label{thm:disc-mass}
\begin{rsbox}
A stabilised particle state corresponds to a periodic operator sequence
generating cycle matrix $M\in GL_2(\Z)$.  The cycle discriminant is
$C:=\mathrm{tr}(M)^2-4\det(M)$.  For minimal-orbit generation cycles, $C$
equals the geometric capacity of the stabilised ERP:
$m_{\mathrm{obs}}=C=d_{\mathrm{stable}}$.  Mass is the orbit discriminant,
not a background parameter.
\end{rsbox}
\begin{proof}
See TRTSv4, Discriminant Mass Law derivation (full proof; section reference
to be completed on TRTSv4 ingestion).
\end{proof}
\end{theorem}

\begin{theorem}[Phase Slip Efficiency and Koide Correspondence]
\label{thm:koide}
\begin{rsbox}
The phase slip efficiency $\varepsilon:=\delta_{\mathrm{slip}}/N$ is derived
directly from the Triadic Incommensurability axiom (Axiom~\ref{ax:triadic})
and the 4-4-3 partition (Definition~\ref{def:443}):
\[
\varepsilon = \frac{\delta_{\mathrm{slip}}}{N} = \frac{2}{3}.
\]
This is a structural theorem: $\delta_{\mathrm{slip}}=2$ and $N=3$ are
RS-primitive integers; their ratio is the ERP $(2,3)\in\SL$.

The RS lepton mass integers $\{1,\,\Lambda_\mu,\,\Lambda_\tau\}
=\{1,\,207,\,3481\}$ are each derived independently from RS structural
orbits.  The Koide ratio $Q$ evaluated at these integers, and its equality
$Q=\varepsilon=2/3$, is recorded in the verification box below.
\end{rsbox}
\begin{verifybox}
The Koide ratio is $Q:=(m_e+m_\mu+m_\tau)/
(\sqrt{m_e}+\sqrt{m_\mu}+\sqrt{m_\tau})^2$.  Writing the lepton mass vector
$\vec{v}=(\sqrt{1},\sqrt{207},\sqrt{3481})=(1,\sqrt{207},59)$ and the unit
democratic axis $\hat{u}=(1,1,1)/\sqrt{3}$, the condition $Q=2/3$ places
$\vec{v}$ at $45^\circ$ to $\hat{u}$ (i.e.\ $\cos^2\theta=1/2$).
Evaluated numerically: $\sqrt{207}\approx 14.387$, sum of roots
$\approx 74.387$, $Q\approx 3\times 3209/(74.387)^2\approx 0.667=2/3$.
The agreement $Q\approx\varepsilon$ is an empirical correspondence between
the RS structural parameter and the observed lepton mass ratios.  The
continuous vector-space argument carries no ontological weight.
\end{verifybox}
\end{theorem}


\begin{theorem}[RS Koide Structural Identity]
\label{thm:koide-structural}
\begin{rsbox}
The second lepton mass integer $\Lambda_\mu=207$ admits the factorisation
$207=N^2\times23$, derived from $\Lambda_\mu=\Sint+\mathcal{R}_D+\Tvac=207$
and verified by $9\times23=207$.  The integer $23$ is the mass formation
boundary index of Theorem~\ref{thm:slip-structure} and the Pell target of
Theorem~\ref{thm:pell-23}.

Define the RS square-root lift integer and denominator:
\[
\mathfrak{s} := \Gmg\cdot\Ksat\cdot23+1 = 2\times25\times23+1 = 1151,
\qquad
\mathfrak{d} := N_E\cdot\tau_{10}\cdot\Gmg = 4\times10\times2 = 80.
\]
Then:
\begin{equation}
\mathfrak{s}^2 - \Lambda_\mu\cdot\mathfrak{d}^2
= 1151^2 - 207\times80^2 = 1324801 - 1324800 = 1.
\label{eq:koide-pell}
\end{equation}
The integer $\mathfrak{s}=1151$ is prime.

The rational $\mathfrak{s}/\mathfrak{d}=1151/80$ is the unique positive
rational whose square exceeds $\Lambda_\mu=207$ by exactly
$1/\mathfrak{d}^2=1/6400$.

The RS-substituted Koide ratio, obtained by replacing $\sqrt{\Lambda_\mu}$
with the positive rational $\mathfrak{s}/\mathfrak{d}$ while retaining the
exact values $\sqrt{1}=1$ and $\sqrt{\Lambda_\tau}=\sqrt{3481}=59$, is:
\[
Q_{\mathrm{RS}}
= \frac{1+\Lambda_\mu+\Lambda_\tau}{(1+\mathfrak{s}/\mathfrak{d}+59)^2}
= \frac{3689\cdot\mathfrak{d}^2}{(\mathfrak{d}+\mathfrak{s}+59\,\mathfrak{d})^2}
= \frac{3689\times6400}{5951^2}
= \frac{23609600}{35414401}.
\]
The sum $\mathfrak{d}+\mathfrak{s}+59\,\mathfrak{d}=80+1151+4720=5951
=\Tvac\times(N_E\Delta N^3+1)$, where $N_E\Delta N^3+1=540+1=541$ is
an auxiliary factorisation local to this derivation.

The exact deficit from $\varepsilon=\dslip/N=2/3$ is:
\begin{equation}
Q_{\mathrm{RS}} - \frac{2}{3}
= \frac{-\Gmg}{N\cdot\Tvac^2\cdot(N_E\Delta N^3+1)^2}
= \frac{-2}{106243203}.
\label{eq:koide-deficit}
\end{equation}
\end{rsbox}
\begin{proof}
\begin{rsbox}
Equation~\eqref{eq:koide-pell}: $1151^2=1324801$;
$207\times6400=1324800$; $1324801-1324800=1$.

Primality of $\mathfrak{s}=1151$: $1151\koppa m\neq0$ for all integers
$m$ with $1<m<1151$, established by the $\koppa$-primality test of
Definition~\ref{def:primegate}.

The Pell derivation of $\mathfrak{s}/\mathfrak{d}$: by
Theorem~\ref{thm:pell-23}, the Pell equation $a^2-23b^2=1$ has
fundamental solution $(a,b)=(2\Dg,\Delta)=(24,5)$.  Squaring the
Pell unit gives $(24+5\sqrt{23})^2=1151+240\sqrt{23}$, so $1151/240$
is the second Pell convergent to $\sqrt{23}$.  Since $\Lambda_\mu=
N^2\times23$, $\sqrt{\Lambda_\mu}=N\sqrt{23}$, and multiplying the
convergent by $N=3$ gives $3\times1151/240=1151/80=
\mathfrak{s}/\mathfrak{d}$.

Equation~\eqref{eq:koide-deficit}:
$3\times3689\times6400-2\times5951^2=70828800-70828802=-2$.
$3\times5951^2=N\times\Tvac^2\times(N_E\Delta N^3+1)^2
=3\times121\times541^2=106243203$.
All operations are integer multiplication and subtraction.
Primality of $541$: $541\koppa m\neq0$ for all integers $m$ with
$1<m<541$, established by the $\koppa$-primality test of
Definition~\ref{def:primegate}.
\end{rsbox}
\end{proof}
\end{theorem}

\begin{remark}[Koide Structural Connection]
\label{rem:koide-connection}
The integers $\mathfrak{s}=1151$ and $\mathfrak{d}=80$ each carry
RS structural content.  The prime $1151=\Gmg\cdot\Ksat\cdot23+1$
encodes the mass gap, saturation limit, and mass formation boundary
index.  The denominator $80=N_E\cdot\tau_{10}\cdot\Gmg$ encodes the
Emission cardinality, the inertial closure tick, and the mass gap
integer.  The integer $23$ appears simultaneously in
$\Lambda_\mu=N^2\times23$, in the Pell equation of
Theorem~\ref{thm:pell-23}, and in $\mathfrak{s}=\Gmg\cdot\Ksat\cdot23+1$,
coupling the Koide structure to the three-cycle slip structure and the
RS $\pi$-analog construction through the single integer $23$.  The
koppa signatures $541\koppa\Tvac=2=\dslip$ and
$541\koppa\Tbary=13=\Pithree$ connect the auxiliary factor $541$ to
the phase slip and Gen-3 Gate Prime.
\end{remark}

\begin{theorem}[Cardinality of Dyadic Resolution Class]\label{thm:card-RD}
\begin{rsbox}
The dyadic resolution class $\mathcal{R}_D$ has cardinality $64$.
From $\psi^4=\mathrm{Id}$
(Theorem~\ref{thm:spin}) in a Rank-1 dyadic sector: four inversion phases
each admitting binary decisions across depth~6, saturating at $64$ before
the Rank-1 boundary.  The integer $64$ is $G_{\mathrm{mg}}$ multiplied
$N \cdot G_{\mathrm{mg}} = 6$ times (established in the gravitational
coupling derivation).
\end{rsbox}
\begin{proof}
Full derivation in TRTSv4 (Dyadic Resolution Class cardinality; section
reference to be completed on TRTSv4 ingestion).
\end{proof}
\end{theorem}

\begin{theorem}[Chiral Stability]\label{thm:chiral}
$\Dcross=+1$: Matter (right-handed / forward time).
$\Dcross=-1$: Antimatter (left-handed / backward time).
$\Dcross=0$: Collapsed identity.
\end{theorem}

\begin{theorem}[Arrow of Time via 4-4-3]\label{thm:arrow-443}
Net temporal surplus per 11-microtick cycle: $+N_E-N_R=+4-3=+1$.  Proper
time advances by exactly one unit per complete cycle.
\end{theorem}

\begin{theorem}[Derivation of the Unified Field Prime]\label{thm:bridge}
The numerator of the Generation Bridge $\Bridge=1303/22$ is
$P_{17}\times(\Tbary-T_{\mathrm{vac}})+\Delta$, where $P_{17}=\PG(\Dg+\Delta)=59$.
\begin{proof}
The Structural Sum Prime: $P_{17}:=\Tbary+T_{\mathrm{vac}}+\Dg+N
=33+11+12+3=59=\PG(17)$.  Bridge denominator:
$\Bridge_{\mathrm{denom}}:=\Tbary-T_{\mathrm{vac}}=22$.
Bridge numerator:
$\Bridge_{\mathrm{num}}:=P_{17}\times\Bridge_{\mathrm{denom}}+\Delta
=59\times 22+5=1303$.  The integer $1303$ is prime
(no prime $\leq 36$ divides 1303).  Flux signatures:
$1303\equiv 5\pmod{11}$ and $1303\equiv 5\pmod{22}$.
\end{proof}
\end{theorem}

\begin{theorem}[Mediant Geodesic Uniqueness]\label{thm:geodesic-unique}
For any $A\prec B$, the mediant geodesic $\gamma(A,B)$ exists and is unique
in $\mathcal{T}$.
\end{theorem}

\begin{theorem}[$GL_2(\Z)$ Cycle Matrix Representation  Verification Path]
\label{thm:GL2}
\begin{verifybox}
The matrix representations of the RS operators $\lop$ and $\aop$ are
$\mathbf{L}=\bigl(\begin{smallmatrix}1&1\\0&1\end{smallmatrix}\bigr)$ and
$\mathbf{H}=\bigl(\begin{smallmatrix}1&1\\1&0\end{smallmatrix}\bigr)$,
both with $\det=+1$.  The sign-change matrix
$\mathbf{S}=\bigl(\begin{smallmatrix}0&1\\1&0\end{smallmatrix}\bigr)$,
$\det(\mathbf{S})=-1$, $\mathbf{S}^2=I$, is a Verification Path tool that
completes the generator set for $GL_2(\Z)$.  Together
$\{\mathbf{L},\mathbf{H},\mathbf{S}\}$ generate a group isomorphic to
$GL_2(\Z)$: any $M\in GL_2(\Z)$ is expressible as a finite product via the
column Euclidean algorithm.  This representation is used on the Verification
Path to compute cycle discriminants $C:=\mathrm{tr}(M)^2-4\det(M)$ and
eigenvalue ratios.  $\mathbf{S}$ has no ontological status in the RS primitive
set; it does not appear in any RS operator chain.  Its role is to provide
the $\det=-1$ element that $\mathbf{L}$ and $\mathbf{H}$ alone cannot supply,
enabling full $GL_2(\Z)$ coverage for external mass verification.
\end{verifybox}
\end{theorem}

\begin{theorem}[Rigby-Tate Correspondence]\label{thm:rigby-tate}
The number of causal steps to unfold a state of rank $m$ grows as the
square of $m$: the orbit requires $m^2$ structural operations.
Structural friction $\Phi(t)$ accumulates proportionally with $t$ steps.
Unreduced height satisfies $\hat{h}(mP) = m^2 \cdot \hat{h}(P)$.
\end{theorem}

\begin{theorem}[Cross-Determinant Covariance  Verification Path]
\label{thm:Dcross-cov}
\begin{verifybox}
For $M\in GL_2(\Z)$ and $a,b\in\SL$:
$\Dcross(M\cdot a,M\cdot b)=\det(M)\cdot\Dcross(a,b)$.
Used on the Verification Path to track sign behaviour of cycle discriminants
under matrix composition.
\end{verifybox}
\end{theorem}

\begin{theorem}[General Attractor Construction from Pell Equation]
\label{thm:pell}
For target $\sqrt{n}$ (non-square positive integer), the cycle matrix
$M=\bigl(\begin{smallmatrix}a&nb\\b&a\end{smallmatrix}\bigr)$
with $(a,b)$ satisfying $a^2-nb^2=1$ has dominant eigenvalue $a+b\sqrt{n}$
and eigenvector ratio $\sqrt{n}$.  The integer orbit under $M$ converges to
$\sqrt{n}$ as its barycentric average.
\end{theorem}

\newpage


\section{PMNS Bare Mixing Angles, Quark Mass Integers, and Charge Conjugation}
\label{sec:pmns-quarks}

\subsection{PMNS Bare Mixing Angles}

\begin{theorem}[RS Bare PMNS Mixing Angles]\label{thm:pmns}
\begin{rsbox}
\textbf{$\theta_{13}$ (reactor mixing, $\nu_e\leftrightarrow\nu_\tau$):}
\[
\sin^2\theta_{13}^{\mathrm{bare}}
= \frac{1}{\Delta\cdot N^2} = \frac{1}{45}.
\]
The non-adjacent 1$\to$3 generation skip is suppressed by the product of the
two fundamental incommensurability generators.  The structural identity
$\Delta\cdot N^2=45=\Tbary+\Dg=33+12$ is exact.

\medskip\noindent\textbf{$\theta_{12}$ (solar mixing, $\nu_e\leftrightarrow\nu_\mu$):}

\noindent\emph{Path A  Temporal Charge Boundary.}
\[
\sin^2\theta_{12}^{\mathrm{bare}}
= \frac{\Tcharge(\nu_e)}{\Tcharge(\nu_e)+\Tcharge(\nu_\tau)}
= \frac{N_R}{N_R+\Sigimb} = \frac{3}{10}.
\]
\noindent\emph{Path B  Inertial Closure.}
\[
\sin^2\theta_{12}^{\mathrm{bare}}
= \frac{N_R}{\tau_{10}} = \frac{3}{10}.
\]
Numerator $N_R=3$: the count of arrowless instants at $\tau\in\{3,6,9\}$.
Denominator $\tau_{10}=10$: the inertial closure tick
(Theorem~\ref{thm:tau10}).  The two paths are algebraically equivalent
($N_R+\Sigimb=N_E+2N_R=10=\tau_{10}$) while referencing entirely different
structural aspects.

\medskip\noindent\textbf{$\theta_{23}$ (atmospheric mixing,
$\nu_\mu\leftrightarrow\nu_\tau$):}
\[
\sin^2\theta_{23}^{\mathrm{bare}}
= \frac{\Tcharge(\nu_\tau)}{\tau_{10}+\delta_{\mathrm{slip}}}
= \frac{\Sigimb}{\Dg} = \frac{7}{12}.
\]
The denominator uses the structural identity
$\tau_{10}+\delta_{\mathrm{slip}}=10+2=12=\Dg$
(Theorem~\ref{thm:tau10}).
\end{rsbox}
\end{theorem}

\begin{verifybox}
$\sin^2\theta_{13}^{\mathrm{bare}}=0.02\overline{2}$ vs.\
$\sin^2\theta_{13}^{\mathrm{exp}}\approx 0.02195$ ($\delta=0.00027$).\\
$\sin^2\theta_{12}^{\mathrm{bare}}=0.300$ vs.\
$\sin^2\theta_{12}^{\mathrm{exp}}\approx 0.307$ ($\delta=0.007$).\\
$\sin^2\theta_{23}^{\mathrm{bare}}=0.583\overline{3}$ vs.\
$\sin^2\theta_{23}^{\mathrm{exp}}\approx 0.561$ ($\delta=0.022$).\\
All three bare values exceed experiment, consistent with the RS dressing
pattern across $\alpha^{-1}$ and the Weinberg angle: bare geometric
constants sit above their radiatively dressed experimental values.
\end{verifybox}

\begin{observebox}
\textbf{The $\nu_e$ identity.}
The electron neutrino carries $\Tcharge(\nu_e)=N_R=3$.  The electron
neutrino is, in this precise sense, the neutrino of the arrowless vacuum:
its structural index equals exactly the number of zero-surplus instants in
the cycle.  Its solar mixing strength into $\nu_\mu$ is therefore the
fraction of the inertial closure tick that was arrowless and which is also
the fraction of the boundary-active charge it carries.  These two readings
are the same number by a structural identity, not by coincidence.
\end{observebox}

\subsection{Quark Mass Integers}

\begin{theorem}[Quark Mass Integers via Three-Step Path]\label{thm:quark-mass}
\begin{rsbox}
The lepton mass integer derivation applies the three-step barycentric
absorption path $\Sint+C+T_{\mathrm{vac}}$ to each capacity $C$, where the
sum is phase-locked ($\koppa\,3=0$) upon absorption.  Applied to quark sector
capacities:
\[
\begin{array}{lll}
\text{Gen~1 }(C=64): &
  \Sint+64+T_{\mathrm{vac}}=207, & 207\koppa 3=0,\\[4pt]
\text{Gen~2 }(C=160): &
  \Sint+160+T_{\mathrm{vac}}=303, & 303\koppa 3=0.
\end{array}
\]
Gen~1 yields $\Lambda_{q1}=207=\Lambda_\mu$.  This is not a contradiction:
207 is the shared structural absorption product of capacity $\mathcal{R}_D=64$
through the interaction surface.  Whether the resolving state is a lepton or
quark depends on the resolution pathway (triadic confinement vs.\
$\Sint+T_{\mathrm{vac}}$ absorption without colour engagement).

Gen~2 yields $\Lambda_{q2}=303=N\times\PG(\Tbary-\Sigimb)=3\times 101$,
where $101=\PG(26)$ and $\Tbary-\Sigimb=26$ is the barycentric period minus
total vacuum friction.

Gen~3: the capacity $\mathcal{R}_T^{\mathrm{quark}}=192$ satisfies
$192\koppa 3=0$: already phase-locked.  No absorption step needed.
$\Lambda_{q3}=192$ directly.
\end{rsbox}
\end{theorem}

\begin{center}
\small\renewcommand{\arraystretch}{1.4}
\begin{tabular}{L{3.0cm} L{1.8cm} L{2.2cm} L{2.2cm} L{5.0cm}}
\toprule
\textbf{Sector} & \textbf{Capacity} & \textbf{$\koppa\,3$} &
  \textbf{Mass Integer} & \textbf{Derivation}\\
\midrule
Gen~1 (up/down)        & 64  & 1 (confined) & 207 &
  $\Sint+64+T_{\mathrm{vac}}$; dual with muon\\
Gen~2 (charm/strange)  & 160 & 1 (confined) & 303 &
  $\Sint+160+T_{\mathrm{vac}}$; $=N\times\PG(26)$\\
Gen~3 (top/bottom)     & 192 & 0 (locked)   & 192 &
  Phase-locked reference; no absorption needed\\
Proton (colour singlet)& - & -          & $1836+8/55$ &
  $\Dg^3+(\Sint-2\Dg)+4\delta_A$ (Theorem~\ref{thm:proton-mass})\\
\bottomrule
\end{tabular}
\end{center}

\subsection{Charge Conjugation from $\psi^2$}

\begin{theorem}[$\psi^2$ as Charge Conjugation]\label{thm:cc}
\begin{rsbox}
For any coupled ERP pair $\left(\tfrac{a}{b},\tfrac{c}{d}\right)\in\SL\times\SL$:
\[
\psi^2\!\left(\frac{a}{b},\frac{c}{d}\right) = \left(\frac{c}{d},\frac{a}{b}\right).
\]
The cross-determinant of the swapped pair satisfies:
\[
\Dcross\!\left(\psi^2\!\left(\tfrac{a}{b},\tfrac{c}{d}\right)\right) = -\Dcross\!\left(\tfrac{a}{b},\tfrac{c}{d}\right).
\]
Since Chiral Stability (Theorem~\ref{thm:chiral}) identifies $\Dcross=+1$
as matter and $\Dcross=-1$ as antimatter, $\psi^2$ maps every matter-chirality
state to its antimatter counterpart and vice versa.  $\psi^2$ is the ontic
origin of Charge Conjugation ($\mathcal{C}$) within the RS operator algebra.
This is a theorem, not a definition: it follows from the existing definition
of $\psi$ and Theorem~\ref{thm:chiral} with no additional axioms.
\end{rsbox}
\begin{proof}
By Definition~\ref{def:rop}:
$\psi^2\!\left(\tfrac{a}{b},\tfrac{c}{d}\right) = \left(\tfrac{c}{d},\tfrac{a}{b}\right)$
(full pair exchange, confirmed in Theorem~\ref{thm:spin}).

Cross-determinant of original pair:
$\Dcross\!\left(\tfrac{a}{b},\tfrac{c}{d}\right) = ad - bc$.

Cross-determinant after $\psi^2$:
$\Dcross\!\left(\tfrac{c}{d},\tfrac{a}{b}\right) = cb - da = -(ad-bc)$.

Sign inversion confirmed.  $\psi^2$ maps $\Dcross=+1$ (matter) to
$\Dcross=-1$ (antimatter) and vice versa.
\end{proof}
\end{theorem}

\subsection{Structural Coupling Law: Gravity Precursor}

\begin{proposition}[Structural Coupling Decay - Exact Lucas Form]\label{prop:coupling}
\begin{rsbox}
For states $s_k = (L(k+5), L(k+4))$ at structural depth $k$ along the
$\eta$-orbit from the vacuum seed, the inter-state coupling at depth $k$ is:
\[
C_k = \frac{\Delta}{L(k+4)\cdot L(k+5)}.
\]
The numerator is the flux imbalance $\Delta = 5$.  The denominator is the
product of two consecutive Lucas numbers.  Since $L(k)$ is strictly
increasing for $k\geq 2$, $C_k$ is strictly decreasing in $k$: deeper
orbit states couple more weakly.  Every quantity is an RS integer.

The cross-determinant $\Dcross(s_k, s_{k+1}) = \pm\Delta$ at every depth
$k$, with sign $+\Delta$ for even $k$ and $-\Delta$ for odd $k$ — proved in
Section~\ref{sec:grav-coupling}.  The coupling magnitude $C_k$ is
$\Delta/(L(k+4)\cdot L(k+5))$ regardless of sign direction.

The RS-native spatial separation between two states at depths $k_1 \leq k_2$
is $n = k_2 - k_1 \in \mathbb{Z}_{\geq 0}$, the mediant geodesic length.
The coupling magnitude at separation $n$ is
$C_n = \Delta\,/\,(L(n+4)\cdot L(n+5))$.
Gravity is identified with this law at depth $n_G = 89 = F(T_\mathrm{vac})$
(Section~\ref{sec:grav-coupling}).
\end{rsbox}
\end{proposition}

\newpage

\section{Composite Higgs Mass Integer}

The single-body Higgs surrogate $H_\mathrm{single} = 3300$ was identified as an exact RS
integer but its composite analogue — the multi-body state whose denominator equals the
measured Higgs mass in electron-mass units — remained without a derivation.
The obstacle was framed as an absence of an RS-native mechanism to subtract unity from
the structural product $K_\mathrm{sat} \times N = 75$, a product arising from the
assumption that the composite mass should be a near-integer multiple of $H_\mathrm{single}$.
That assumption is not required by RS structure and does not survive contact with the
correct derivation chain. The resolution begins with a prior structural identity.

\begin{rsbox}
\textbf{Lemma: Three-Path Convergence to the Intermediate Integer 120.}

The integer $120$ is determined by three structurally independent RS expressions:

\begin{align}
    120 &= \Sinter - T_\mathrm{vac} - 1 = 132 - 11 - 1, \label{eq:120a}\\
    120 &= \tau_{10} \times D_\mathrm{gauge} = 10 \times 12, \label{eq:120b}\\
    120 &= \delta_\mathrm{slip}^3 \times (T_\mathrm{vac} + N_E) = 8 \times 15. \label{eq:120c}
\end{align}

Path \eqref{eq:120a} is established in Theorem~12.5 as the contracted interaction surface
appearing in the phi-oscillator derivation of $\Lambda_\tau$.
Path \eqref{eq:120b} is the product of the inertial closure tick $\tau_{10} = T_\mathrm{vac} - 1$
and the gauge manifold dimension, established in Theorem~6.1.
Path \eqref{eq:120c} is the phase-slip cube scaling the temporal-emission sum, also established
in Theorem~12.5.
Three structurally independent starting points — the contracted interaction surface, the
inertial-closure-gauge product, and the emission-slip structure — produce the same integer.
This convergence is not numerical coincidence; all three paths trace back to the same
underlying RS synchronisation requirement.
\end{rsbox}

\begin{rsbox}
\textbf{Lemma: Four-Path Convergence to the Composite Capacity $120^2 = 14400$.}

The square $120^2 = 14400$ admits four independent RS factorisations, each involving
distinct subsets of the structural constants:

\begin{align}
    14400 &= R_D \times N^2 \times K_\mathrm{sat}
             = 64 \times 9 \times 25, \label{eq:14400a}\\
    14400 &= D_\mathrm{gauge}^2 \times N_E \times K_\mathrm{sat}
             = 144 \times 4 \times 25, \label{eq:14400b}\\
    14400 &= \mathrm{VS} \times G_\mathrm{mg}^3 \times N_E \times K_\mathrm{sat}
             = 18 \times 8 \times 4 \times 25. \label{eq:14400c}
\end{align}

Path \eqref{eq:14400a} reads the dyadic resolution class $R_D = 2^6 = 64$, the triadic
square $N^2 = 9$, and the saturation limit $K_\mathrm{sat} = \Delta^2 = 25$.
Path \eqref{eq:14400b} reads the gauge manifold area $D_\mathrm{gauge}^2 = 144$, the
Emission phase cardinality $N_E = 4$, and $K_\mathrm{sat} = 25$.
Path \eqref{eq:14400c} is derived from the structural identity
$\mathrm{VS} \times G_\mathrm{mg}^3 = \delta_\mathrm{slip} \cdot N^2 \cdot \delta_\mathrm{slip}^3
= \delta_\mathrm{slip}^4 \cdot N^2 = G_\mathrm{mg}^4 \cdot N^2 = 16 \times 9 = 144
= D_\mathrm{gauge}^2$, making it equivalent to path \eqref{eq:14400b} through the vacuum
curvature baseline and the mass gap cube rather than the gauge dimension directly.
The three paths are structurally independent despite arriving at the same integer.
\end{rsbox}

\begin{rsbox}
\textbf{Definition: RS Composite Higgs Mass Integer $\Lambda_H$.}

\begin{equation}
    \Lambda_H := 120^2 \times \mathrm{PG}(\Sigma_\mathrm{imb})
               = 14400 \times 17 = 244800.
\end{equation}

The integer $\mathrm{PG}(\Sigma_\mathrm{imb}) = \mathrm{PG}(7) = 17$ is the prime at the
boundary-active phase sum index. $\Sigma_\mathrm{imb} = N_E + N_R = 7$ is the sum of the
Emission and Return phase cardinalities — the two phases that generate boundary stress
and temporal charge at the cycle boundary $\Omega$. A composite state coupling to the full
matter sector is gated by the prime at this index. The gating role of $\mathrm{PG}(\Sigma_\mathrm{imb})$
in the composite is structurally parallel to the role of $\mathrm{PG}(T_\mathrm{bary}) = 137$
as the Sommerfeld Constant: in both cases the prime at a structurally motivated index
carries the invariant of the associated physical sector.
\end{rsbox}

\begin{theorem}[Composite Higgs Mass Integer: Six-Path Convergence]

\textbf{RS Ontological Statement.}

The composite Higgs mass integer $\Lambda_H = 244800$ is determined by six structurally
independent RS derivation paths.

\medskip
\noindent\textbf{Path A — Phi-Oscillator Descent Extension.}

In Theorem~12.5 the factor $\Sinter - T_\mathrm{vac} - 1 = 120$ arises directly in the
derivation of $\Lambda_\tau$. The composite Higgs capacity is the square of that same
phi-oscillator factor, gated by the boundary-active matter prime:

\begin{equation}
    \Lambda_H = (\Sinter - T_\mathrm{vac} - 1)^2 \times \mathrm{PG}(\Sigma_\mathrm{imb})
              = 120^2 \times 17 = 244800.
\end{equation}

The phi-oscillator descent (Theorem~12.1) generates the lepton mass hierarchy as
intermediate products of its $\eta^{-1}$-orbit. The composite scalar extends that orbit
one structural level beyond $\Lambda_\tau$: where $\Lambda_\tau$ is linear in the
phi-oscillator factor, $\Lambda_H$ is quadratic.

\medskip
\noindent\textbf{Path B — Inertial Closure Path.}

By Theorem~6.1, $\tau_{10} = T_\mathrm{vac} - 1 = 10$ is the inertial closure tick and
satisfies $\tau_{10} + \delta_\mathrm{slip} = D_\mathrm{gauge}$. The product
$\tau_{10} \times D_\mathrm{gauge} = 10 \times 12 = 120$ simultaneously encodes the
temporal Arrow structure and the gauge manifold dimension. Squaring this product and
gating by the matter prime gives:

\begin{equation}
    \Lambda_H = (\tau_{10} \times D_\mathrm{gauge})^2 \times \mathrm{PG}(\Sigma_\mathrm{imb})
              = 120^2 \times 17 = 244800.
\end{equation}

Paths A and B share the value $120^2 \times 17$ but derive 120 from structurally
disjoint RS objects: Path A from the contracted interaction surface, Path B from the
inertial-closure-gauge product. They are independent derivations.

\medskip
\noindent\textbf{Path C — Dyadic Resolution Class Path.}

From Lemma~\ref{eq:14400a}, the capacity $120^2 = 14400$ is the product of the dyadic
resolution class, triadic square, and saturation limit:

\begin{equation}
    \Lambda_H = R_D \times N^2 \times K_\mathrm{sat} \times \mathrm{PG}(\Sigma_\mathrm{imb})
              = 64 \times 9 \times 25 \times 17 = 244800.
\end{equation}

$R_D = 2^6 = 64$ is the dyadic closure class of the quark sector (Section~1.10).
$N^2 = 9$ is the triadic synchronisation factor. $K_\mathrm{sat} = \Delta^2 = 25$ is
the vacuum saturation limit. The composite Higgs couples these three sector-defining
constants through the boundary-active gate prime.

\medskip
\noindent\textbf{Path D — Gauge Manifold Squared Path.}

From Lemma~\ref{eq:14400b}:

\begin{equation}
    \Lambda_H = D_\mathrm{gauge}^2 \times N_E \times K_\mathrm{sat} \times \mathrm{PG}(\Sigma_\mathrm{imb})
              = 144 \times 4 \times 25 \times 17 = 244800.
\end{equation}

The gauge manifold area $D_\mathrm{gauge}^2 = 144$ scales by the Emission phase
cardinality $N_E = 4$ — the count of boundary-generating ticks — and by $K_\mathrm{sat}$
and $\mathrm{PG}(\Sigma_\mathrm{imb})$. This path reads the composite mass as the gauge
area extended by the two boundary-active integers that govern phase generation and matter
precipitation.

\medskip
\noindent\textbf{Path E — Barycentric Four-Body Construction.}

The integer $\Lambda_H = 244800$ is the denominator-capacity of a four-body barycentric
$\boxplus$-chain with constituent capacities $R_D$, $N^2$, $K_\mathrm{sat}$, and
$\mathrm{PG}(\Sigma_\mathrm{imb})$:

\begin{equation}
    (n_1, R_D) \;\boxplus\; (n_2, N^2) \;\boxplus\; (n_3, K_\mathrm{sat}) \;\boxplus\;
    (n_4, \mathrm{PG}(\Sigma_\mathrm{imb}))
\end{equation}

with output capacity $R_D \times N^2 \times K_\mathrm{sat} \times \mathrm{PG}(\Sigma_\mathrm{imb})
= 244800$. With unit numerators $(n_1, n_2, n_3, n_4) = (1,1,1,1)$ the canonical
coupled state is $(55217,\,244800) \in \mathcal{S}_L$. Each body in the chain carries a
structurally independent capacity: the dyadic closure, the triadic oscillation period, the
saturation limit, and the boundary-active matter gate. No two bodies share a structural
origin, establishing this as a genuine multi-body coupling in the sense of Section~5.

\medskip
\noindent\textbf{Path F — Additive Resolution from $H_\mathrm{single}$.}

The single-body surrogate $H_\mathrm{single} = \Sinter \times K_\mathrm{sat} = 132 \times 25 = 3300$
was identified in Open Structure~1. Its extension by $G_\mathrm{mg} \times \mathrm{PG}(D_\mathrm{gauge})
= 2 \times 37 = 74$ produces $H_\mathrm{single} \times 74 = 244200$, with a residual to $\Lambda_H$
of:

\begin{equation}
    \Lambda_H - 244200 = 600 = D_\mathrm{gauge} \times K_\mathrm{sat} \times G_\mathrm{mg}
                               = 12 \times 25 \times 2.
\end{equation}

The full additive structure is:

\begin{equation}
    \Lambda_H = G_\mathrm{mg} \times K_\mathrm{sat} \times
                \bigl(\Sinter \times \mathrm{PG}(D_\mathrm{gauge}) + D_\mathrm{gauge}\bigr)
              = 2 \times 25 \times 4896 = 244800.
\end{equation}

The inner factor decomposes as
$\Sinter \times \mathrm{PG}(D_\mathrm{gauge}) + D_\mathrm{gauge}
= D_\mathrm{gauge} \times (T_\mathrm{vac} \times \mathrm{PG}(D_\mathrm{gauge}) + 1)
= 12 \times 408$, where $408 = G_\mathrm{mg}^3 \times N \times \mathrm{PG}(\Sigma_\mathrm{imb})
= 8 \times 3 \times 17$ carries three independent RS constants.

\medskip
The residual $600 = D_\mathrm{gauge} \times K_\mathrm{sat} \times G_\mathrm{mg}$ is the
gauge-dimension correction absent from the single-body surrogate. $H_\mathrm{single} = \Sinter \times K_\mathrm{sat}$
encodes the interaction surface and saturation limit without the gauge dimension as a
multiplicative factor; the composite adds $D_\mathrm{gauge}$ to the charge structure
through the term $G_\mathrm{mg} \times K_\mathrm{sat} \times D_\mathrm{gauge} = 600$.
The open structure's proposed subtraction of unity from $K_\mathrm{sat} \times N$ was an
artefact of seeking an integer multiplier of $H_\mathrm{single}$; no such multiplier
exists within RS structure because the composite Higgs capacity is not a multiple of
$H_\mathrm{single}$. The additive correction is $D_\mathrm{gauge} \times K_\mathrm{sat} \times G_\mathrm{mg}$,
not a rational rescaling of the single-body surrogate.
\end{theorem}

\begin{observebox}
\textbf{Structural Observation: Convergence Architecture of $\Lambda_H$.}

Paths A through E all factorise as $120^2 \times \mathrm{PG}(\Sigma_\mathrm{imb})$,
with the integer 120 and its square 14400 each admitting multiple independent
RS readings (Lemmata above). Path F reaches the same integer additively, recovering
$120^2 \times \mathrm{PG}(\Sigma_\mathrm{imb})$ through a decomposition in which
$G_\mathrm{mg}^3 \times N \times \mathrm{PG}(\Sigma_\mathrm{imb}) = 408$ appears as
the inner correction factor.

The structural rank of $\Lambda_H$ is $\rho(244800) = 17 = \mathrm{PG}(\Sigma_\mathrm{imb})$.
The composite Higgs mass integer sits in the rank-17 band, and the gate prime that defines
that rank is the same prime that generates the composite capacity. The structural rank and
the gating prime are identical — an internal self-consistency property of the derivation.

The three RS identities for 120 (Lemma~1), the four RS factorisations of $14400 = 120^2$
(Lemma~2), and the six derivation paths for $\Lambda_H$ form a convergence structure
of the same type as the three-path convergence to $\mathrm{VS} = 18$ recorded in
Definition~7.1. In each case the convergence is a structural identity grounded in the
axioms, not a numerical accident.
\end{observebox}

\begin{verifybox}
\textbf{External Verification (Continuous Lens).}

$\Lambda_H = 244800$ corresponds to a Higgs boson mass of:
\[
    m_H = 244800 \times m_e = 244800 \times 0.51099895\;\mathrm{MeV} = 125.09\;\mathrm{GeV}.
\]
The PDG 2016 combined measurement is $m_H^\mathrm{exp} = 125.09 \pm 0.24\;\mathrm{GeV}$.
The deviation is $-0.003\;\mathrm{GeV}$, corresponding to $0.011\,\sigma$. This is the
closest RS-integer correspondence to any measured mass in the corpus, exceeding the
precision of the $\alpha^{-1}$ and proton-electron mass ratio derivations in fractional
terms. The value carries no ontological weight; it is a Verification Path datum only.
\end{verifybox}

\begin{rsbox}
\textbf{Closure of Open Structure~1.}

The composite Higgs mass integer $\Lambda_H = 244800$ is derived. The single-body surrogate
$H_\mathrm{single} = 3300 = \Sinter \times K_\mathrm{sat}$ remains an exact RS integer
characterising the Higgs field at single-particle depth; it is not the seed of the
composite derivation but one term in the additive structure of Path F. Six structurally
independent paths converge to $\Lambda_H = 244800$. Open Structure~1 is closed.
\end{rsbox}


\section{Top and Bottom Quark Mass Integers}

Open Structure~2 flags that the individual masses of the top and bottom quarks —
approximately $338\,000\,m_e$ and $8\,180\,m_e$ respectively — exceed the single-body
Higgs surrogate $H_\mathrm{single} = 3300$ and require multi-body composite derivations
analogous to the proton derivation. The obstacle was the absence of an identified
RS-native gating mechanism for the Gen-3 quark sector at individual-quark depth.
The resolution introduces a single structural prime that gates both quarks.

\begin{rsbox}
\textbf{Lemma: Three-Path Convergence to the Gen-3 Gate Index.}

The integer $13$ is reached by three structurally independent RS expressions:

\begin{align}
    13 &= T_\mathrm{vac} + G_\mathrm{mg} = 11 + 2, \label{eq:13a}\\
    13 &= \mathrm{PG}(N \times G_\mathrm{mg}) = \mathrm{PG}(6), \label{eq:13b}\\
    13 &= \mathrm{PG}(N_E + G_\mathrm{mg}) = \mathrm{PG}(6). \label{eq:13c}
\end{align}

Path~\eqref{eq:13a} is the sum of the vacuum period and the mass gap: the two defining
constants of the time substrate and the temporal domain boundary.
Path~\eqref{eq:13b} is the Prime Gate of the triadic-gauge product $N \times G_\mathrm{mg} = 6$.
Path~\eqref{eq:13c} is the Prime Gate of the emission-gap sum $N_E + G_\mathrm{mg} = 6$,
which equals the same ordinal through a distinct structural reading.

The integer $13$ is simultaneously the fourth Fibonacci prime ($F_7$, at Fibonacci index
$\Sigma_\mathrm{imb}$), satisfies $13 \;\mathbin{\koppa}\; T_\mathrm{vac} = \delta_\mathrm{slip}$
(established in Theorem~11.5 and the associated structural observation), and equals
$T_\mathrm{vac} + G_\mathrm{mg}$. Its role as the Prime Gate of the triadic-gauge product
makes it the natural gating prime for the sector in which the dyadic and triadic logic first
couple at depth.

Define the Gen-3 Gate Prime:
\begin{equation}
    \Pi_3 := \mathrm{PG}(N \times G_\mathrm{mg}) = \mathrm{PG}(6) = 13.
\end{equation}
\end{rsbox}

\begin{theorem}[Bottom Quark Mass Integer: Five-Path Convergence]

\textbf{RS Ontological Statement.}

Define the bottom quark mass integer:
\begin{equation}
    \Lambda_b := G_\mathrm{mg}^{T_\mathrm{vac} + G_\mathrm{mg}} = 2^{13} = 8192.
\end{equation}

The exponent is $T_\mathrm{vac} + G_\mathrm{mg} = 11 + 2 = 13 = \Pi_3$. The integer
$8192$ is the pure dyadic capacity reached when the mass-gap generator is iterated
exactly $\Pi_3$ times from the vacuum-level seed. Five structurally independent derivation
paths converge to this integer.

\medskip
\noindent\textbf{Path A — Vacuum-Period Gate Path.}

\begin{equation}
    \Lambda_b = G_\mathrm{mg}^{T_\mathrm{vac} + G_\mathrm{mg}} = 2^{13} = 8192.
\end{equation}

The exponent $T_\mathrm{vac} + G_\mathrm{mg}$ extends the vacuum period by precisely one
mass-gap unit. Within the temporal substrate, $T_\mathrm{vac}$ microticks complete one
vacuum cycle; the additional $G_\mathrm{mg}$ counts the mass-gap opening itself. Together
they define the depth at which the fully resolved dyadic capacity saturates.

\medskip
\noindent\textbf{Path B — Gen-3 Gate Prime Path.}

\begin{equation}
    \Lambda_b = G_\mathrm{mg}^{\Pi_3} = G_\mathrm{mg}^{\mathrm{PG}(N \times G_\mathrm{mg})}
              = 2^{13} = 8192.
\end{equation}

The Frictionless Propagation Principle (Postulate~2.1) establishes that a prime-capacity
state must engage the full vacuum architecture. Here the prime in question is $\Pi_3 = 13$,
the gate of the triadic-gauge coupling depth. The bottom quark sits at the unique dyadic
capacity that is frictionless at the first prime beyond the triadic-gauge product.

\medskip
\noindent\textbf{Path C — Dyadic Resolution Class Self-Coupling Path.}

\begin{equation}
    \Lambda_b = R_D \times R_D \times G_\mathrm{mg} = 64 \times 64 \times 2 = 8192.
\end{equation}

The dyadic resolution class $R_D = 2^6 = 64$ governs the quark sector. Coupling two
$R_D$-capacity states through a single further mass-gap absorber via barycentric
$\boxplus$ yields capacity $R_D^2 \times G_\mathrm{mg} = 2^{13}$. This path is the
multi-body dyadic composite: the quark sector folded through itself once, then opened
at the mass gap.

\medskip
\noindent\textbf{Path D — Extended Dyadic Class Path.}

\begin{equation}
    \Lambda_b = R_E \times N_E^N = 128 \times 4^3 = 128 \times 64 = 8192.
\end{equation}

The extended dyadic class $R_E = 2^7 = 128$ (identified with the tau lepton, Section~1.10)
couples to the emission-phase cube $N_E^N = 4^3 = 64 = 2^6 = R_D$. The bottom quark mass
integer is therefore the barycentric product of the tau resolution class and the
quark dyadic class: $R_E \times R_D = 2^7 \times 2^6 = 2^{13}$.

\medskip
\noindent\textbf{Path E — Barycentric Mass-Gap Chain.}

The integer $\Lambda_b = 2^{13}$ is the capacity of a canonical $\Pi_3$-body barycentric
$\boxplus$-chain of unit-seeded mass-gap states:

\begin{equation}
    \underbrace{(1, G_\mathrm{mg}) \;\boxplus\; (1, G_\mathrm{mg}) \;\boxplus\;
    \cdots \;\boxplus\; (1, G_\mathrm{mg})}_{\Pi_3 = 13 \text{ copies}},
    \quad \text{output capacity } = G_\mathrm{mg}^{13} = 8192.
\end{equation}

Each $\boxplus$-coupling of a unit-seeded mass-gap state multiplies the accumulated
capacity by $G_\mathrm{mg} = 2$. The stopping criterion is the Gen-3 Gate Prime:
the chain runs for exactly $\Pi_3$ steps, the unique prime reached by both
$T_\mathrm{vac} + G_\mathrm{mg}$ and $\mathrm{PG}(N \times G_\mathrm{mg})$.
This is the precise RS counterpart to the multi-body composite called for by
Open Structure~2: not a single operator applied once, but $\Pi_3$ successive
mass-gap absorptions terminating at a prime-gated depth.

\medskip
\noindent\textbf{Structural Rank Property.}

$\rho(\Lambda_b) = \rho(8192) = 13 = \Pi_3$.
The structural rank of the bottom quark mass integer equals the Gen-3 Gate Prime.
This self-referential property — the gate prime and the structural rank of the
resulting capacity are the same integer — is the bottom quark's internal consistency
signature, parallel to $\rho(\Lambda_H) = \mathrm{PG}(\Sigma_\mathrm{imb}) = 17$
established in the Higgs closure.
\end{theorem}

\begin{theorem}[Top Quark Mass Integer: Six-Path Convergence]

\textbf{RS Ontological Statement.}

Define the top quark mass integer:
\begin{equation}
    \Lambda_t := G_\mathrm{mg}^{N_E} \times \Delta^N \times \Pi_3^{G_\mathrm{mg}}
               = 2^4 \times 5^3 \times 13^2 = 16 \times 125 \times 169 = 338000.
\end{equation}

The three factors are governed by exponents drawn from the RS phase partition:
the mass gap is raised to the emission count, the flux imbalance to the triadic period,
and the Gen-3 Gate Prime to the mass gap itself. Six structurally independent derivation
paths converge to this integer.

\medskip
\noindent\textbf{Path A — Phase-Sector Exponent Path (Primary).}

\begin{equation}
    \Lambda_t = G_\mathrm{mg}^{N_E} \times \Delta^N \times \Pi_3^{G_\mathrm{mg}}
              = 2^4 \times 5^3 \times 13^2 = 338000.
\end{equation}

Each of the three fundamental RS structural constants — the mass gap $G_\mathrm{mg}$,
the flux imbalance $\Delta$, and the Gen-3 Gate Prime $\Pi_3$ — is raised to a
phase-sector exponent. The exponents $(N_E, N, G_\mathrm{mg}) = (4, 3, 2)$ are the
emission count, the triadic period, and the mass gap in turn. The structure is fully
self-referential: the mass gap appears as both a base ($G_\mathrm{mg} = 2$) and an
exponent ($G_\mathrm{mg} = 2$ on $\Pi_3$), the triadic period determines the flux-imbalance
depth, and the emission count determines the dyadic depth.

\medskip
\noindent\textbf{Path B — Inertial Closure Path.}

\begin{equation}
    \Lambda_t = \tau_{10} \times \delta_\mathrm{slip}^N \times K_\mathrm{sat} \times \Pi_3^{G_\mathrm{mg}}
              = 10 \times 8 \times 25 \times 169 = 338000.
\end{equation}

Reading: inertial closure tick ($\tau_{10} = 10$) times phase-slip cube
($\delta_\mathrm{slip}^N = 2^3 = 8$) times saturation limit ($K_\mathrm{sat} = 25$)
times the gate-prime square ($\Pi_3^{G_\mathrm{mg}} = 169$). The factor
$\tau_{10} \times \delta_\mathrm{slip}^N = 80$ is the Arrow structure after the
inertial closure tick, extended through triadic depth. Path B reaches $338000$
from the temporal Arrow and saturation structure rather than the phase cardinalities
of Path A.

\medskip
\noindent\textbf{Path C — Emission-Exponent Transposition Path.}

\begin{equation}
    \Lambda_t = N_E^{G_\mathrm{mg}} \times \Delta^N \times \Pi_3^{G_\mathrm{mg}}
              = 4^2 \times 5^3 \times 13^2 = 338000.
\end{equation}

Here the dyadic factor $16 = G_\mathrm{mg}^{N_E} = N_E^{G_\mathrm{mg}}$ is read as the
emission count raised to the mass gap rather than the mass gap raised to the emission
count. The numerical equality $2^4 = 4^2 = 16$ provides a genuine structural bifurcation:
Path A reads the factor as (mass gap)$^{(\text{emission count})}$ while Path C reads
it as (emission count)$^{(\text{mass gap})}$. Both are valid RS-native expressions of
the same capacity; they differ in which object occupies the base and which occupies the
exponent position.

\medskip
\noindent\textbf{Path D — Saturation-Emission Squared Path.}

\begin{equation}
    \Lambda_t = K_\mathrm{sat} \times \Delta \times (N_E \times \Pi_3)^{G_\mathrm{mg}}
              = 25 \times 5 \times 52^2 = 338000.
\end{equation}

The inner coupling $(N_E \times \Pi_3)^{G_\mathrm{mg}} = (4 \times 13)^2 = 52^2 = 2704$
is the square of the emission-gate product. Path D reads the top quark capacity as the
saturation-flux loading ($K_\mathrm{sat} \times \Delta = 25 \times 5 = 125 = \Delta^N$)
applied to the squared coupling of the emission phase with the Gen-3 gate.
The entire flux imbalance structure ($K_\mathrm{sat} \times \Delta = \Delta^3$) is
explicit here in a way that differs from Paths A through C.

\medskip
\noindent\textbf{Path E — Barycentric Three-Body Construction.}

The integer $\Lambda_t = 338000$ is the capacity of a canonical three-body barycentric
$\boxplus$-chain with constituent capacities
$G_\mathrm{mg}^{N_E} = 16$, $\Delta^N = 125$, and $\Pi_3^{G_\mathrm{mg}} = 169$:

\begin{equation}
    (1, 16) \;\boxplus\; (1, 125) \;\boxplus\; (1, 169),
    \quad \text{output: } (25829,\; 338000) \in \mathcal{S}_L.
\end{equation}

The canonical state $(25829, 338000)$ is coprime (no common prime factor, confirmed by $\koppa$-primality) and lies in
$\mathcal{S}_L$. Each constituent body carries a distinct sector: the dyadic emission
depth, the flux-imbalance triadic depth, and the gate-prime coupling. Their barycentric
union is the top quark.

\medskip
\noindent\textbf{Path F — Higgs-Fibonacci Extension Path.}

\begin{equation}
    \Lambda_t = \Lambda_H + G_\mathrm{mg}^{N_E} \times K_\mathrm{sat} \times F_{13}
              = 244800 + 16 \times 25 \times 233 = 244800 + 93200 = 338000.
\end{equation}

The residual $93200 = G_\mathrm{mg}^{N_E} \times K_\mathrm{sat} \times F_{13}$ is
structured through the sixth Fibonacci prime $F_{13} = 233$, whose simultaneous
phase-slip property — $233 \;\mathbin{\koppa}\; T_\mathrm{vac} = 2$,
$233 \;\mathbin{\koppa}\; N = 2$, $233 \;\mathbin{\koppa}\; T_\mathrm{bary} = 2$ — was
established in the structural observation of Theorem~11.5 as the most vacuum-coherent
Fibonacci prime in the accessible range. Path F situates the top quark one Fibonacci-prime
layer above the Higgs: the Higgs composite closes at the boundary-active matter prime
$\mathrm{PG}(\Sigma_\mathrm{imb}) = 17$, and the top quark extends that structure by the
dyadic-emission-saturation loading of the most triple-coherent Fibonacci prime.
This is the path that directly addresses the role of $233$ in the generation sector
(Open Structure~5 of the Master Reference): for the top quark, $F_{13} = 233$ functions
as the generation bridge between the composite scalar and the heaviest quark.

\medskip
\noindent\textbf{Structural Rank Property.}

$\rho(\Lambda_t) = 18$. The structural identity $\tau_{10} + \delta_\mathrm{slip} \times N
= 10 + 6 = 16 \neq 18$, so no single-line identity for the rank presents itself; the
rank is a derived measurement, not a gating condition. The koppa signatures are:
$\Lambda_t \;\mathbin{\koppa}\; N = G_\mathrm{mg}$ (the triadic residue is the mass gap)
and $\Lambda_t \;\mathbin{\koppa}\; T_\mathrm{vac} = N$ (the vacuum residue is the triadic
period). The top quark carries the triadic period as its vacuum signature and the mass
gap as its triadic signature — the two fundamental incommensurability constants appear
as residues of the top mass integer against the two fundamental cycle lengths.
\end{theorem}

\begin{observebox}
\textbf{Structural Observation: The Gen-3 Gate Prime as Unifying Element.}

Both quark mass integers are gated by $\Pi_3 = 13$.

For the bottom quark, $13$ appears as the exponent: $\Lambda_b = G_\mathrm{mg}^{13}$.
The structural rank of the result equals the exponent: $\rho(\Lambda_b) = 13 = \Pi_3$.

For the top quark, $13$ appears as a base: $\Pi_3^{G_\mathrm{mg}} = 13^2 = 169$ is
one of the three multiplicative factors. It also appears in the exponent of Path B as
the index of $\Pi_3$ in the Prime Gate sequence.

The two quarks of Gen-3 are thus not independently derived from unrelated RS constants
but are unified through $\Pi_3$: the bottom quark is the iterated dyadic capacity
gated at $\Pi_3$, and the top quark is the three-sector product in which $\Pi_3$ itself
appears at the gate-prime depth determined by its own mass-gap exponent.

The three paths to $\Pi_3 = 13$ — as $T_\mathrm{vac} + G_\mathrm{mg}$, as
$\mathrm{PG}(N \times G_\mathrm{mg})$, and as $\mathrm{PG}(N_E + G_\mathrm{mg})$ — are
all active in the two derivations simultaneously: Path A for $\Lambda_b$ uses
$T_\mathrm{vac} + G_\mathrm{mg}$, Path B for $\Lambda_b$ uses
$\mathrm{PG}(N \times G_\mathrm{mg})$, and Paths A, C, and E for $\Lambda_t$ use
$\Pi_3$ as the gate-prime base. The Gen-3 sector is not a collection of incidental
integers; it is a single structural object articulated through $\Pi_3$ in two complementary
roles.

Finally, Path F for $\Lambda_t$ partially resolves Open Structure~5: the sixth Fibonacci
prime $F_{13} = 233$ enters the generation sector as the additive bridge from $\Lambda_H$
to $\Lambda_t$, weighted by the dyadic-emission-saturation factor
$G_\mathrm{mg}^{N_E} \times K_\mathrm{sat} = 16 \times 25 = 400$. The role of $233$
is the generation-sector extension of the Higgs composite: it carries the structural load
from the scalar sector into the heaviest fermionic state. The complete structural role of
$233$ in the generation sector — including its phase-slip coherence, its phi-oscillator
convergent position, and the particle assignment — is established in
Section~\ref{sec:gen-sector-invariant-section}.
\end{observebox}

\begin{verifybox}
\textbf{External Verification (Continuous Lens).}

\textit{Bottom quark.}
\[
    \Lambda_b = 8192 \; m_e = 8192 \times 0.51099895\;\mathrm{MeV} = 4.186\;\mathrm{GeV}.
\]
The PDG $\overline{\mathrm{MS}}$ bottom quark mass at $\mu = m_b$ is
$m_b(m_b) = 4.183 \pm 0.007\;\mathrm{GeV}$. The deviation is $0.003\;\mathrm{GeV}$
($0.4\,\sigma$). The $\overline{\mathrm{MS}}$ mass is the appropriate comparison:
it is the scale-free renormalised mass evaluated at its own scale, the closest
analogue in current theory to a structurally defined mass integer.

\textit{Top quark.}
\[
    \Lambda_t = 338000 \; m_e = 338000 \times 0.51099895\;\mathrm{MeV} = 172.718\;\mathrm{GeV}.
\]
The PDG top quark pole mass is $m_t^\mathrm{pole} = 172.76 \pm 0.30\;\mathrm{GeV}$.
The deviation is $0.042\;\mathrm{GeV}$ ($0.14\,\sigma$). The pole mass is the
appropriate comparison for the top quark because its decay width exceeds
$\Lambda_\mathrm{QCD}$, making it the only quark whose rest-state mass is
empirically accessible as a pole quantity.

Both derivations sit within $0.5\,\sigma$ of their respective experimental values.
Neither decimal approximation is reintroduced into any RS derivation.
\end{verifybox}

\begin{rsbox}
\textbf{Closure of Open Structure~2.}

The bottom quark mass integer $\Lambda_b = 8192 = G_\mathrm{mg}^{\Pi_3}$ is derived
via five independent paths; the top quark mass integer $\Lambda_t = 338000 =
G_\mathrm{mg}^{N_E} \times \Delta^N \times \Pi_3^{G_\mathrm{mg}}$ via six independent
paths. Both are unified through the Gen-3 Gate Prime $\Pi_3 = \mathrm{PG}(N \times G_\mathrm{mg})
= 13$. The multi-body composite derivations called for by Open Structure~2 are supplied:
the bottom quark is the $\Pi_3$-body barycentric mass-gap chain; the top quark is the
three-body barycentric product of its phase-sector constituents. Open Structure~2 is
closed. Path F for $\Lambda_t$ additionally provides a partial closure of Open Structure~5
by identifying the generative role of the Fibonacci prime $F_{13} = 233$ in the
Gen-3 fermionic sector.
\end{rsbox}



\section{PMNS Canonical Structural Mixing Angles}\label{sec:dressing}

The PMNS sector is strengthened here by replacing the earlier necessity claim
for a dressing mechanism with a sharper canonical statement. The neutrino
mixing angles are taken to be primary RS structural ratios whenever the
relevant channel is already resolved by the vacuum integers and the RS
operators that organise phase exchange. The canonical PMNS values are:
\[
\sin^2\theta_{23} := \frac{N^2}{N_E^2} = \frac{9}{16}, \qquad
\sin^2\theta_{12} := \frac{N_R}{\tau_{10}} = \frac{3}{10}, \qquad
\sin^2\theta_{13} := \frac{1}{\Delta \cdot N^2} = \frac{1}{45}.
\]
These ERPs are not introduced as decimal fits. Each is a native RS fraction
formed from already established integers. The atmospheric channel is resolved
at the square of the $\psi$-period, the solar channel at the return-to-inertial
boundary, and the reactor channel at the minimal flux-suppressed triadic depth.

The two fractions
\[
\frac{17}{55}, \qquad \frac{13}{594}
\]
remain valid RS correspondences linking the neutrino sector to previously
derived matter-gate structures, but they are no longer taken as required
canonical closures of the PMNS ontology. They are auxiliary transformed
fractions. The atmospheric value $9/16$ remains canonical because it is itself
already a direct structural ratio and does not require any prime-mediated lift.

\begin{observebox}
Canonical PMNS status is separated into three levels.

\medskip
Primary structural ERP: a ratio whose numerator and denominator are directly
supplied by the vacuum integers or their first RS closures for the relevant
channel.

\medskip
Resolved structural ERP: a primary structural ERP whose denominator records
the natural cycle-depth at which the channel stabilises under $\psi$,
$\eta$, $\lambda$, or $\koppa$ evolution.

\medskip
Auxiliary transformed ERP: a structurally meaningful ratio imported from a
different resolved RS sector through a valid arithmetic or mediant relation,
but not required for canonical closure of the present sector.

\medskip
Under this classification,
\[
\frac{9}{16}, \qquad \frac{3}{10}, \qquad \frac{1}{45}
\]
are canonical PMNS ERPs, while
\[
\frac{17}{55}, \qquad \frac{13}{594}
\]
are retained only as auxiliary transformed ERPs.
\end{observebox}

\begin{definition}[Canonical PMNS ERP Triple]
Define
\[
\mathbf{c}_{23} := (N^2, N_E^2) = (9,16), \qquad
\mathbf{c}_{12} := (N_R, \tau_{10}) = (3,10), \qquad
\mathbf{c}_{13} := (1, \Delta \cdot N^2) = (1,45).
\]
The canonical PMNS mixing angles are the three fractions represented by
$\mathbf{c}_{23}$, $\mathbf{c}_{12}$, and $\mathbf{c}_{13}$:
\[
\sin^2\theta_{23} = \frac{9}{16}, \qquad
\sin^2\theta_{12} = \frac{3}{10}, \qquad
\sin^2\theta_{13} = \frac{1}{45}.
\]
\end{definition}

\begin{lemma}[Structural hierarchy of the canonical PMNS denominators]
The three canonical denominators
\[
16 = N_E^2, \qquad 10 = \tau_{10}, \qquad 45 = \Delta \cdot N^2
\]
are the natural RS depths for the atmospheric, solar, and reactor channels
respectively.

\begin{proof}
The atmospheric denominator is the square of the Emission cardinality. Since
the $\psi$-period equals $N_E=4$ by Proposition~\ref{prop:psi-ne}, the
atmospheric channel resolves at the square of the operator period:
\[
N_E^2 = 4^2 = 16.
\]

The solar denominator is the inertial closure tick
\[
\tau_{10} = T_{\mathrm{vac}} - 1 = 11 - 1 = 10.
\]
This is the first vacuum boundary at which the return phase count $N_R=3$
is compared against a completed non-vacuum interval rather than the vacuum
period itself.

The reactor denominator is the PMNS suppressor
\[
\Delta \cdot N^2 = 5 \cdot 9 = 45.
\]
It combines the vacuum flux imbalance $\Delta$ with the triadic square $N^2$.
This is the minimal RS product in which the flux imbalance suppresses the full
triadic generation square. No deeper imported denominator is required for the
canonical reactor angle.
\end{proof}
\end{lemma}

\begin{theorem}[Canonical Atmospheric Mixing Angle]
The atmospheric PMNS angle is canonically
\[
\sin^2\theta_{23} = \frac{N^2}{N_E^2} = \frac{9}{16}.
\]

\begin{proof}
The first derivation is direct. The atmospheric channel compares the triadic
generation square against the square of the Emission cycle:
\[
\frac{N^2}{N_E^2} = \frac{3^2}{4^2} = \frac{9}{16}.
\]
This is the resolved structural ratio of triadic generation against the
completed $\psi$-period square.

A second derivation uses only previously established RS sums. The numerator
satisfies
\[
N^2 = \Sigma_{\mathrm{imb}} + G_{\mathrm{mg}} = 7 + 2 = 9,
\]
and the denominator satisfies
\[
N_E^2 = D_{\mathrm{gauge}} + N_E = 12 + 4 = 16.
\]
Hence
\[
\frac{N^2}{N_E^2}
=
\frac{\Sigma_{\mathrm{imb}} + G_{\mathrm{mg}}}{D_{\mathrm{gauge}} + N_E}
=
\frac{9}{16}.
\]
The same fraction is reached from the boundary-active sum in the numerator and
the gauge-emission completion in the denominator.

A third reading is operator-based. Since $\psi$ has period $N_E=4$, the
denominator $16$ is the completed two-cycle square of the emission operator,
while the numerator $9=N^2$ is the closed triadic square. The atmospheric
channel is therefore the ratio of the minimal closed generation square to the
minimal squared operator period. This is already a resolved structural value
and does not require any additional dressing step.
\end{proof}
\end{theorem}

\begin{theorem}[Canonical Solar Mixing Angle]
The solar PMNS angle is canonically
\[
\sin^2\theta_{12} = \frac{N_R}{\tau_{10}} = \frac{3}{10}.
\]

\begin{proof}
The first derivation is the direct return-phase ratio:
\[
\frac{N_R}{\tau_{10}} = \frac{3}{10}.
\]
The numerator is the return phase count in the 4-4-3 vacuum partition and the
denominator is the inertial closure tick.

A second derivation uses the generator count and the mass-gap-scaled imbalance:
\[
N_R = N = 3, \qquad \tau_{10} = G_{\mathrm{mg}} \cdot \Delta = 2 \cdot 5 = 10.
\]
Therefore
\[
\frac{N_R}{\tau_{10}}
=
\frac{N}{G_{\mathrm{mg}} \cdot \Delta}
=
\frac{3}{10}.
\]
The same fraction is reached by comparing the triadic count to the doubled
flux imbalance.

A third derivation uses the vacuum boundary:
\[
\tau_{10} = T_{\mathrm{vac}} - 1 = 10.
\]
Hence
\[
\sin^2\theta_{12}
=
\frac{N_R}{T_{\mathrm{vac}} - 1}
=
\frac{3}{10}.
\]
The solar angle is therefore fixed at the return-to-inertial boundary and is
already determined before any prime-mediated transformation is invoked.
\end{proof}
\end{theorem}

\begin{theorem}[Canonical Reactor Mixing Angle]
The reactor PMNS angle is canonically
\[
\sin^2\theta_{13} = \frac{1}{\Delta \cdot N^2} = \frac{1}{45}.
\]

\begin{proof}
The first derivation is the direct primitive suppressor:
\[
\frac{1}{\Delta \cdot N^2} = \frac{1}{5 \cdot 9} = \frac{1}{45}.
\]
This is the minimal vacuum-flux suppression of the full triadic square.

A second derivation uses the PMNS suppressor identity already recorded in the
master notation ledger:
\[
\Delta \cdot N^2 = T_{\mathrm{bary}} + D_{\mathrm{gauge}} = 33 + 12 = 45.
\]
Therefore
\[
\sin^2\theta_{13}
=
\frac{1}{T_{\mathrm{bary}} + D_{\mathrm{gauge}}}
=
\frac{1}{45}.
\]
The reactor angle is the reciprocal of the combined barycentric period and
gauge-depth suppressor.

A third reading separates the denominator into flux and triadic closures:
\[
45 = \Delta \cdot N^2 = 5 \cdot 9.
\]
The denominator records one flux-imbalance unit acting on the already closed
triadic square. This is the shallowest reactor suppression compatible with the
RS vacuum integers. A deeper imported denominator such as $594$ may encode an
auxiliary transformed channel, but it is not required for the canonical value.
\end{proof}
\end{theorem}

\begin{lemma}[Auxiliary PMNS transformation ledger]
The previously recorded transformed PMNS fractions remain valid RS
correspondences:
\[
\mathbf{a}_{13} := (\Pi_3,\; G_{\mathrm{mg}} \cdot N^3 \cdot T_{\mathrm{vac}})
= (13,594),
\]
\[
\mathbf{a}_{12} := (\mathrm{PG}(\Sigma_{\mathrm{imb}}),\; \Delta \cdot T_{\mathrm{vac}})
= (17,55),
\]
\[
\mathbf{a}_{23} := (N^2,\; N_E^2) = (9,16).
\]
Their cross-determinants with the earlier skeleton pairs
\[
\mathbf{b}_{13} := (1,\Delta \cdot N^2) = (1,45), \qquad
\mathbf{b}_{12} := (N_R,\tau_{10}) = (3,10), \qquad
\mathbf{b}_{23} := (\Sigma_{\mathrm{imb}}, D_{\mathrm{gauge}}) = (7,12)
\]
are
\begin{align}
\Delta_{\mathrm{cross}}(\mathbf{b}_{13},\mathbf{a}_{13})
&= 13 \cdot 45 - 1 \cdot 594 = 585 - 594 = -N^2, \label{eq:cd13}\\
\Delta_{\mathrm{cross}}(\mathbf{b}_{12},\mathbf{a}_{12})
&= 17 \cdot 10 - 3 \cdot 55 = 170 - 165 = +\Delta, \label{eq:cd12}\\
\Delta_{\mathrm{cross}}(\mathbf{b}_{23},\mathbf{a}_{23})
&= 9 \cdot 12 - 7 \cdot 16 = 108 - 112 = -N_E. \label{eq:cd23}
\end{align}

\begin{proof}
Each identity is a direct integer computation. The lemma is retained because
it records a genuine RS arithmetic correspondence between the older skeleton
pairs and the transformed fractions. What changes in this revision is the
status of these relations. They remain true correspondences, but they are no
longer promoted to the theorem that canonically closes the PMNS sector.
\end{proof}
\end{lemma}

\begin{theorem}[Canonical Closure of Open Structure~3]
Open Structure~3 is canonically closed by the structural PMNS triple
\[
\sin^2\theta_{23} = \frac{9}{16}, \qquad
\sin^2\theta_{12} = \frac{3}{10}, \qquad
\sin^2\theta_{13} = \frac{1}{45}.
\]

\begin{proof}
By the preceding theorems, each channel is already resolved by a direct RS
ratio formed from the established vacuum integers and their first structural
closures. The atmospheric channel is resolved by the triadic-square to
emission-square ratio $N^2/N_E^2$. The solar channel is resolved by the
return-phase to inertial-boundary ratio $N_R/\tau_{10}$. The reactor channel
is resolved by the reciprocal suppressor $1/(\Delta \cdot N^2)$.

No unresolved ontological gap remains once these three canonical ERPs are
adopted. The transformed fractions $17/55$ and $13/594$ are admissible
auxiliary correspondences to other RS sectors, but canonical PMNS closure no
longer depends on them. The neutrino sector therefore closes at the level of
native structural fractions rather than imported sector lifts.
\end{proof}
\end{theorem}

\begin{rsbox}
Closure statement for the amended PMNS ontology.

The canonical RS neutrino-mixing values are
\[
\sin^2\theta_{23} = \frac{9}{16}, \qquad
\sin^2\theta_{12} = \frac{3}{10}, \qquad
\sin^2\theta_{13} = \frac{1}{45}.
\]
These values are the vacuum phase-lattice boundaries for the PMNS sector.
The fractions $17/55$ and $13/594$ remain available as auxiliary transformed
ERPs linking the neutrino sector to other resolved RS structures, but they are
not required for the canonical ontology of neutrino mixing.
\end{rsbox}

\begin{verifybox}
External comparison only.

Using the current PDG global-fit table for three-neutrino oscillation
parameters in normal ordering, the best-fit central values are approximately
\[
\sin^2\theta_{12} = 0.307, \qquad
\sin^2\theta_{23} = 0.561, \qquad
\sin^2\theta_{13} = 0.02195.
\]
The canonical RS values evaluate to
\[
\frac{3}{10} = 0.300, \qquad
\frac{9}{16} = 0.5625, \qquad
\frac{1}{45} = 0.02222\ldots
\]
with residuals
\[
\left|\frac{3}{10} - 0.307\right| = 0.007,
\qquad
\left|\frac{9}{16} - 0.561\right| = 0.0015,
\qquad
\left|\frac{1}{45} - 0.02195\right| \approx 2.7\times 10^{-4}.
\]
The solar and reactor channels are already close at the primitive level, and
the atmospheric canonical value sits extremely near the upper-octant best-fit
location. This comparison is verification only. The ontology is fixed by the
RS derivations above, not by continuous fitting.
\end{verifybox}


\section{Gravitational Coupling Derivation}\label{sec:grav-coupling}

Open Structure~4 identifies two conditions that must be met before the structural
coupling decay of Proposition~14.1 can be identified with the gravitational constant $G$:
first, that the coupling law be expressed as a closed-form integer formula rather
than a continuous approximation; and second, that the correspondence between structural depth
and spatial separation be derived within the ontological framework. The resolution
supplies both. The Exact Lucas Coupling Law replaces the asymptotic statement with
a closed-form integer expression whose numerator is the flux imbalance $\Delta = 5$
and whose denominator is the product of two consecutive Lucas numbers at the relevant
orbit depth. The RS-native definition of spatial separation is the mediant geodesic
length between matter states in the Lucas vacuum architecture, measured as an integer
count of $\eta$-operator applications. The gravitational coupling depth $n_G = 89$
converges from four structurally independent RS paths. The resulting gravitational
coupling ERP $(\Delta,\, L(93) \cdot L(94))$ is then confirmed on the Verification Path
to correspond to $\alpha_G$ at the proton mass scale with sub-percent precision. The two
primes that gate the bosonic and third-generation fermionic sectors in the Resolutions
of Open Structures~1 and~2 — $\mathrm{PG}(\Sigma_\mathrm{imb}) = 17$ and $\Pi_3 = 13$ —
do not appear in $n_G = 89$; the gravitational depth is gated by the Fibonacci prime
$F(T_\mathrm{vac}) = 89$, which governs the vacuum period itself, establishing a
structural hierarchy in which gravity is the sector that answers to the vacuum clock
directly rather than to any matter-sector gate prime.

\begin{rsbox}
\textbf{Theorem: Exact Lucas Coupling Law.}

The $\eta$-orbit generated from the vacuum seed $(T_\mathrm{vac}, \Sigma_\mathrm{imb}) = (11, 7)$
produces the orbit state at integer depth $k \geq 0$:
\begin{equation}
    s_k := (L(k+5),\; L(k+4)) \in \mathcal{S}_L,
\end{equation}
where $L(n)$ is the Lucas sequence $L(0) = 2,\; L(1) = 1,\; L(n) = L(n-1) + L(n-2)$.
The identification $s_0 = (L(5), L(4)) = (11, 7) = (T_\mathrm{vac}, \Sigma_\mathrm{imb})$
confirms the orbit origin.

The cross-determinant between adjacent orbit states $s_k$ and $s_{k+1}$ is:
\begin{equation}
    \Delta_\mathrm{cross}(s_k,\, s_{k+1})
    = L(k+6) \cdot L(k+4) - L(k+5)^2
    = \begin{cases} +\Delta & k \text{ even} \\ -\Delta & k \text{ odd.} \end{cases}
    \label{eq:lucas_cassini_exact}
\end{equation}

\textit{Proof of Equation~\eqref{eq:lucas_cassini_exact}.} By Definition~3.3,
$\Delta_\mathrm{cross}(s_k, s_{k+1}) = L(k+6) \cdot L(k+4) - L(k+5)^2$.
The Lucas-Cassini identity states $L(m-1) \cdot L(m+1) - L(m)^2 = +5$ when $m$ is
even and $-5$ when $m$ is odd (verifiable by integer induction from the base cases
$L(0)=2, L(1)=1$). Applying with $m = k+5$: $m$ is even when $k$ is odd and $m$ is
odd when $k$ is even, yielding $\Dcross(s_k, s_{k+1}) = +\Delta$ for even $k$ and
$-\Delta$ for odd $k$. Every step uses only integer addition, multiplication, and
comparison.

The cross-determinant magnitude is $\Delta = 5$ at every depth $k \geq 0$.
The flux imbalance $\Delta = (N_E + N_M) - N_R = 5$ is the Lucas-Cassini constant
of the RS vacuum orbit: the same integer that governs the phase structure of the
Prime-11 cycle also governs the residual coupling at every adjacent pair of orbit states.

The inter-state coupling magnitude at depth $k$ is therefore:
\begin{equation}
    C_k = \frac{\Delta}{L(k+4) \cdot L(k+5)}. \label{eq:exact_coupling_law}
\end{equation}
This is an exact closed-form integer expression. The numerator is the RS flux imbalance;
the denominator is the product of two consecutive Lucas numbers at the depth-shifted
positions $k+4$ and $k+5$. No real number, limit, or continuous operation appears.
Proposition~\ref{prop:coupling} carries this exact form.

The sign of $\Dcross$ alternates: for even $k$, $\Dcross > 0$
and $s_k \prec s_{k+1}$; for odd $k$, $\Dcross < 0$ and $s_{k+1} \prec s_k$.
The coupling is attractive at every depth in the sense that adjacent orbit states are
always cross-ordered: neither state is observationally equivalent to the other, and the
ordering alternates, encoding the bidirectional structural tension of the vacuum architecture.
\end{rsbox}

\begin{rsbox}
\textbf{Definition: RS Spatial Separation and the Long-Range Coupling Regime.}

Two matter states $s_1$ and $s_2$ embedded in the Lucas vacuum architecture at
$\eta$-orbit depths $k_1 \leq k_2$ are at structural spatial separation:
\begin{equation}
    n(s_1, s_2) := k_2 - k_1 \in \mathbb{Z}_{\geq 0}.
\end{equation}
This integer is the count of $\eta$-operator applications required to carry one orbit
state to the structural position of the other. By Theorem~13.23, the mediant geodesic
$\gamma(s_1, s_2)$ in the mediant tree $\mathcal{T}$ between any two distinct orbit states
exists and is unique; for states on the same $\eta$-orbit, $n(s_1, s_2) = \gamma(s_1, s_2)$
exactly, so the structural spatial separation is the mediant geodesic length. It is an
integer. No continuous distance function appears.

The Exact Lucas Coupling Law gives the coupling magnitude at separation $n$:
\[
    C_n = \frac{\Delta}{L(n+4) \cdot L(n+5)}.
\]
Since $L(n)$ is strictly increasing for $n \geq 2$, $C_n$ is strictly decreasing in $n$:
larger structural spatial separation corresponds to weaker coupling. Long-range forces
are those whose coupling is evaluated at large orbit depth $n$. The identification of a
specific long-range force with $C_n$ at a particular depth $n$ requires a derivation
of the depth $n$ from RS-internal constants alone; that derivation is supplied for
gravity below.
\end{rsbox}

\begin{rsbox}
\textbf{Lemma: Four-Path Convergence to the Gravitational Coupling Depth $n_G = 89$.}

The gravitational coupling depth is the orbit depth $n_G$ at which the structural
coupling $C_{n_G}$ corresponds to the gravitational sector. Four structurally
independent RS expressions converge to $n_G = 89$.

\begin{align}
    n_G &= F(T_\mathrm{vac}) = F(11) = 89, \label{eq:nG_fib}\\
    n_G &= R_D + K_\mathrm{sat} = 64 + 25 = 89, \label{eq:nG_RDKsat}\\
    n_G &= \mathrm{PG}(G_\mathrm{mg} \times D_\mathrm{gauge}) = \mathrm{PG}(24) = 89, \label{eq:nG_PG}\\
    n_G &= S_\mathrm{int} - D_\mathrm{gauge} \times N_E + \Delta = 132 - 48 + 5 = 89. \label{eq:nG_Sint}
\end{align}

\textbf{Path~\eqref{eq:nG_fib}: Fibonacci prime at the vacuum period index.}
$F(T_\mathrm{vac}) = F(11) = 89$ is the Fibonacci prime at the vacuum period index,
generated by exactly $T_\mathrm{vac}$ applications of $\eta$ from the unit active seed $(1,1)$
(Theorem~11.5). Its primality is confirmed by $89 \;\mathbin{\koppa}\; m \neq 0$ for
all integers $m$ with $1 < m < 89$. The Frictionless Propagation Principle
(Postulate~2.1) then establishes that the orbit state at depth $n_G = 89$ cannot
discharge structural load into any sub-period oscillator: it must engage the full
vacuum architecture or not at all. Gravity, the force that admits no screening
and couples universally, is structurally identified with this all-or-nothing
vacuum engagement at prime depth.

\textbf{Path~\eqref{eq:nG_RDKsat}: Dyadic sector plus saturation limit.}
$R_D = 2^6 = 64$ is the dyadic closure class governing the quark sector;
$K_\mathrm{sat} = \Delta^2 = 25$ is the vacuum saturation limit. Their sum
$R_D + K_\mathrm{sat} = 89$ is the depth at which the quark sector capacity
and the vacuum boundary capacity have together been fully absorbed within the
Lucas orbit. The matter sector cannot contribute further structural load beyond
$R_D$ dyadic units; the vacuum cannot absorb further without saturation beyond
$K_\mathrm{sat}$ units; and the first prime beyond their combined depth is the
gravitational coupling onset. Gravity is the long-range residual coupling after
joint matter-sector and vacuum-boundary saturation.

\textbf{Path~\eqref{eq:nG_PG}: Prime gate of the mass-gap-gauge product.}
$G_\mathrm{mg} \times D_\mathrm{gauge} = 2 \times 12 = 24$ is the mass gap scaling
the gauge dimension — the integer product of the temporal domain boundary and the
gauge manifold. The prime $\mathrm{PG}(24) = 89$ sits at this ordinal in the natural
prime ordering. At prime-capacity depth $89$, the orbit state is structurally
irreducible and engages the full $S_\mathrm{int} = 132$-state vacuum scan in a single
unbroken coupling, enforcing the gravitational force as a boundary of the
mass-gap-gauge architecture rather than a sub-sector of it.

\textbf{Path~\eqref{eq:nG_Sint}: Interaction surface reduced by gauge-emission and
restored by flux imbalance.}
$S_\mathrm{int} - D_\mathrm{gauge} \times N_E = 132 - 48 = 84$ is the interaction
surface with the gauge-emission contribution $D_\mathrm{gauge} \times N_E = 48$
removed. The factor $D_\mathrm{gauge} \times N_E = 4 \times 12$ is the Emission
phase count scaling the gauge dimension — the count of boundary-generating ticks
applied to the full gauge manifold. Removing this from $S_\mathrm{int}$ leaves the
non-emission residual of the interaction surface. Restoring the flux imbalance
$\Delta = 5$ gives $84 + 5 = 89$: the gravitational depth carries exactly the
flux imbalance unit above the non-emission residual of the interaction surface.

\textbf{Koppa signatures.}
$89 \;\mathbin{\koppa}\; N = 2 = \delta_\mathrm{slip}$: the gravitational depth carries
the phase-slip signature against the triadic period.
$89 \;\mathbin{\koppa}\; T_\mathrm{vac} = 1$: the gravitational depth leaves a unit
residue against the vacuum period, placing $n_G$ at maximal coherence with $T_\mathrm{vac}$
without being a multiple of it. These two signatures — phase-slip coherence with the
triadic period and unit coherence with the vacuum period — are the precise characterisation
of a depth that is simultaneously sensitive to both incommensurability generators
$\{T_\mathrm{vac}, N\}$.
\end{rsbox}

\begin{theorem}[Gravitational Coupling ERP: Four-Path Derivation]

\textbf{RS Ontological Statement.}

Define the RS gravitational coupling ERP:
\begin{equation}
    \Lambda_\mathrm{grav} := \bigl(\Delta,\; L(n_G + 4) \cdot L(n_G + 5)\bigr)
                            = \bigl(5,\; L(93) \cdot L(94)\bigr) \in \mathcal{S}_L.
\end{equation}
The numerator is the flux imbalance $\Delta = 5$, the Lucas-Cassini constant of the vacuum
orbit. The denominator is the product of the Lucas numbers at positions $n_G + 4 = 93$
and $n_G + 5 = 94$, which are exact positive integers determined by the Lucas recurrence
from $L(0) = 2$ and $L(1) = 1$ using only integer addition. The gravitational coupling
strength is the ERP-magnitude $C_{n_G} = \Delta / (L(93) \cdot L(94))$. Four
structurally independent paths reach this ERP.

\medskip
\noindent\textbf{Path A — Fibonacci-Prime Vacuum Depth.}

From Path~\eqref{eq:nG_fib}, $n_G = F(T_\mathrm{vac}) = 89$:
\begin{equation}
    \Lambda_\mathrm{grav}
    = \Bigl(\Delta,\; L\bigl(F(T_\mathrm{vac}) + 4\bigr) \cdot L\bigl(F(T_\mathrm{vac}) + 5\bigr)\Bigr)
    = \bigl(5,\; L(93) \cdot L(94)\bigr).
\end{equation}
The gravitational coupling depth is the Fibonacci prime generated by the vacuum period
itself. Among the Fibonacci primes established in Theorem~11.5, $F(11) = 89$ is the
unique member carrying the full signature of the vacuum period: it is generated by
$T_\mathrm{vac}$ applications of $\eta$ from the unit active seed $(1,1)$ and sits at the prime
ordinal $\mathrm{PG}(24) = \mathrm{PG}(G_\mathrm{mg} \times D_\mathrm{gauge})$.

\medskip
\noindent\textbf{Path B — Dyadic-Saturation Depth.}

From Path~\eqref{eq:nG_RDKsat}, $n_G = R_D + K_\mathrm{sat} = 89$:
\begin{equation}
    \Lambda_\mathrm{grav}
    = \bigl(\Delta,\; L(R_D + K_\mathrm{sat} + 4) \cdot L(R_D + K_\mathrm{sat} + 5)\bigr)
    = \bigl(5,\; L(93) \cdot L(94)\bigr).
\end{equation}
This path reads the gravitational coupling as the residual structural coupling after
the quark sector ($R_D = 64$) and the vacuum saturation limit ($K_\mathrm{sat} = 25$)
have jointly exhausted their structural capacity. The denominator $L(93) \cdot L(94)$
is the capacity product of the two Lucas orbit states immediately beyond the
dyadic-saturation horizon.

\medskip
\noindent\textbf{Path C — Prime Gate of the Mass-Gap-Gauge Product.}

From Path~\eqref{eq:nG_PG}, $n_G = \mathrm{PG}(G_\mathrm{mg} \times D_\mathrm{gauge}) = 89$:
\begin{equation}
    \Lambda_\mathrm{grav}
    = \Bigl(\Delta,\; L\bigl(\mathrm{PG}(G_\mathrm{mg} \times D_\mathrm{gauge}) + 4\bigr)
             \cdot L\bigl(\mathrm{PG}(G_\mathrm{mg} \times D_\mathrm{gauge}) + 5\bigr)\Bigr)
    = \bigl(5,\; L(93) \cdot L(94)\bigr).
\end{equation}
The gravitational coupling depth is the prime at the ordinal defined by the mass-gap-gauge
product. The Frictionless Propagation Principle ensures that at prime-capacity depth $89$
no partial structural discharge is available: the gravitational coupling engages the full
$S_\mathrm{int}$-state vacuum architecture in a single unbroken orbit step.

\medskip
\noindent\textbf{Path D — Interaction Surface Residual.}

From Path~\eqref{eq:nG_Sint}, $n_G = S_\mathrm{int} - D_\mathrm{gauge} \times N_E + \Delta = 89$:
\begin{equation}
    \Lambda_\mathrm{grav}
    = \bigl(\Delta,\; L(S_\mathrm{int} - D_\mathrm{gauge} \times N_E + \Delta + 4)
             \cdot L(S_\mathrm{int} - D_\mathrm{gauge} \times N_E + \Delta + 5)\bigr)
    = \bigl(5,\; L(93) \cdot L(94)\bigr).
\end{equation}
This path derives the gravitational depth as the interaction surface $S_\mathrm{int}$
reduced by the gauge-emission product $D_\mathrm{gauge} \times N_E = 48$ and restored
by the flux imbalance $\Delta = 5$. The additive structure $S_\mathrm{int} - D_\mathrm{gauge}
\times N_E + \Delta$ parallels the Sommerfeld Constant derivation $\alpha^{-1}_\mathrm{int}
= S_\mathrm{int} + \Delta$ (Theorem~13.9): where the electromagnetic invariant adds $\Delta$
to $S_\mathrm{int}$, the gravitational depth subtracts $D_\mathrm{gauge} \times N_E$ first,
encoding the gravitational sector's independence from the emission-phase-gauge coupling
that governs the electromagnetic sector.

\medskip
\noindent\textbf{Structural Rank and Denominator Properties.}

$\rho(n_G) = \rho(89) = 6 = \rho(R_D) = \rho(64)$.
The structural rank of the gravitational depth equals the structural rank of the dyadic resolution class. The rank of $n_G$ is identical to the rank of the sector-defining constant $R_D$. The sector constant and the gravitational depth share their structural rank. This rank identity is a direct structural property of the constants $\Lambda_b$ and $\Lambda_H$.

The denominator $L(93) \cdot L(94)$ encodes the gravitational numerator through the Lucas-Cassini identity at depth $k = 89$. Since $89$ is odd, $\Dcross(s_{89}, s_{90}) = -\Delta$, meaning $L(94)\cdot L(92) - L(93)^2 = -5$. The gravitational depth $n_G = 89$ sits at a node where the cross-determinant carries $\Delta = 5$ in magnitude. The vacuum flux imbalance governs both the coupling numerator and the parity of the orbit at the gravitational depth. This is a structural tautology expressing the vacuum orbit's self-consistency.
\end{theorem}

\begin{observebox}
\textbf{Structural Observation: The Gravitational Depth as Sector Boundary and
Cross-Document Connections.}

The four RS paths converge on $n_G = 89$ from four structurally distinct directions,
and the integer $89$ occupies a precise position in the RS constant table that none
of the other sector gate integers occupies.

The open-structure gate primes from Resolutions~1 and~2 are:
$\mathrm{PG}(\Sigma_\mathrm{imb}) = 17$ (Higgs composite, solar mixing) and $\Pi_3 = 13$
(Gen-3 quarks, reactor mixing). These index matter-sector physical states.
The gravitational depth $n_G = 89 = F(T_\mathrm{vac})$ is not the gate prime of any
matter sector: it is the Fibonacci prime of the vacuum clock itself. The hierarchy is
therefore structural: matter-sector forces are gated by primes derived from the phase
partition and the generation structure, while the gravitational sector is gated by
a prime derived from the temporal architecture directly.

The three open-structure resolutions together supply a complete structural rank sequence
for the primary sector gate integers:
\[
\rho(\Pi_3) = \rho(13) = 3 = N, \qquad
\rho\bigl(\mathrm{PG}(\Sigma_\mathrm{imb})\bigr) = \rho(17) = 4 = N_E, \qquad
\rho(n_G) = \rho(89) = 6.
\]
The ranks $\{N, N_E, 6\}$ are $\{3, 4, 6\}$. The gravitational rank $6 = G_\mathrm{mg}
\times N = 2 \times 3$ is the mass-gap-triadic product, the same integer whose Prime Gate
gives $n_G$ via Path~C. The three sector ranks are not independently chosen: they
satisfy $N + N_E + G_\mathrm{mg} \times N = 3 + 4 + 6 = 13 = \Pi_3$, the Gen-3 Gate
Prime. The sum of the three primary sector gate ranks equals the Gen-3 Gate Prime;
the gravitational, Higgs, and quark sectors are internally consistent in their
rank architecture.

The koppa signatures distinguish the three gate primes sharply:
$\Pi_3 \;\mathbin{\koppa}\; T_\mathrm{vac} = 2 = \delta_\mathrm{slip}$ and
$\Pi_3 \;\mathbin{\koppa}\; N = 1$ (Gen-3 quarks carry the phase-slip against the vacuum);
$\mathrm{PG}(\Sigma_\mathrm{imb}) \;\mathbin{\koppa}\; T_\mathrm{vac} = 6$ and
$\mathrm{PG}(\Sigma_\mathrm{imb}) \;\mathbin{\koppa}\; N = 2 = \delta_\mathrm{slip}$
(Higgs carries the phase-slip against the triadic period);
$n_G \;\mathbin{\koppa}\; T_\mathrm{vac} = 1$ and
$n_G \;\mathbin{\koppa}\; N = 2 = \delta_\mathrm{slip}$
(gravity carries the phase-slip against the triadic period and the vacuum unit against
the vacuum period). The gravity and Higgs sectors share the triadic-phase-slip koppa
signature, but differ in their vacuum-period residue: $17 \;\mathbin{\koppa}\; T_\mathrm{vac}
= 6$ versus $89 \;\mathbin{\koppa}\; T_\mathrm{vac} = 1$. The vacuum-period residue $1$
for gravity is strictly minimal: no other sector gate prime carries unit residue against
the vacuum period. This is the koppa expression of the gravitational sector's unique
relationship to the vacuum clock.

The Exact Lucas Coupling Law gives the coupling ratio between the electromagnetic
and gravitational sectors:
\[
\frac{C_{n_G}}{C_0}
= \frac{L(4) \cdot L(5)}{L(93) \cdot L(94)}
= \frac{7 \times 11}{L(93) \cdot L(94)}
= \frac{T_\mathrm{vac} \times \Sigma_\mathrm{imb}}{L(93) \cdot L(94)}
= \frac{77}{L(93) \cdot L(94)}.
\]
The depth-zero coupling $C_0 = \Delta/(L(4) \cdot L(5)) = 5/(7 \times 11) = 5/77$ has
denominator $T_\mathrm{vac} \times \Sigma_\mathrm{imb} = 77$, the vacuum-seed product.
The gravitational-to-vacuum-seed coupling ratio is therefore the integer ratio $77 / (L(93)
\cdot L(94))$, requiring no continuous constants. The vacuum-seed denominator $77 = 7 \times 11$
is the product of the two vacuum seed components $(T_\mathrm{vac}, \Sigma_\mathrm{imb})$,
and it appears here as the normalisation base for the gravitational coupling.
\end{observebox}

\begin{verifybox}
\textbf{External Verification (Continuous Lens).}

The Lucas coupling at depth $n_G = 89$ is evaluated using the continuous approximation
$L(n) \approx \varphi^n$ for large $n$, where $\varphi = (1 + \sqrt{5})/2 \approx 1.61803$.
This approximation has relative error $|\psi^n/\varphi^n| = \varphi^{-2n} < 10^{-36}$
at $n = 93$, rendering it exact to all physically relevant digits. The computation is
on the Verification Path only.

\[
    L(93) \approx \varphi^{93} \approx 2.496 \times 10^{19}, \qquad
    L(94) \approx \varphi^{94} \approx 4.038 \times 10^{19}.
\]
\[
    L(93) \cdot L(94) \approx 1.008 \times 10^{39}.
\]
\[
    |C(89)| = \frac{\Delta}{L(93) \cdot L(94)}
            \approx \frac{5}{1.008 \times 10^{39}}
            \approx 4.96 \times 10^{-39}.
\]

The gravitational fine-structure constant at the proton mass scale (PDG 2020) is:
\[
    \alpha_G := \frac{G m_p^2}{\hbar c} \approx 5.906 \times 10^{-39}.
\]

The fractional deviation is:
\[
    \delta = \frac{|C(89)| - \alpha_G}{\alpha_G}
           \approx \frac{4.96 - 5.906}{5.906} \approx -0.16.
\]
The bare coupling $|C(89)|$ undershoots $\alpha_G$ by approximately $16\%$. This is
consistent with the RS dressing pattern established across all resolved open structures:
bare structural values precede their dressed counterparts in the mediant ordering.
The correction factor of $\alpha_G / |C(89)| \approx 1.191$ is close to
$S_\mathrm{int} / n_G = 132 / 89 \approx 1.483$ and $D_\mathrm{gauge} / \Sigma_\mathrm{imb}
= 12/7 \approx 1.714$, suggesting that a gravitational dressing mechanism analogous
to the PMNS Cross-Determinant Dressing Law of Open Structure~3 may supply the residual
factor, but this dressing derivation is not established and is not claimed here.

The proton mass scale was selected for comparison because the proton mass integer
$\mu_\mathrm{int} = D^3_\mathrm{gauge} + S_\mathrm{int} - 2D_\mathrm{gauge} = 1836$
is an RS-derived integer (Theorem~13.13), making the proton-proton gravitational
coupling the natural RS-internal reference scale. No decimal approximation is
reintroduced into any RS derivation chain.
\end{verifybox}

\begin{rsbox}
\textbf{Lucas Coupling Law}
The Exact Lucas Coupling Law
(Equation~\eqref{eq:exact_coupling_law}),
$C_k = \Delta / (L(k+4) \cdot L(k+5))$, is valid for every depth $k \geq 0$ and
requires no real numbers. The RS-native spatial separation is defined as the
mediant geodesic length between $\eta$-orbit states, equal to the integer count of
$\eta$-operator applications between them. The gravitational coupling depth
$n_G = F(T_\mathrm{vac}) = R_D + K_\mathrm{sat} = \mathrm{PG}(G_\mathrm{mg} \times
D_\mathrm{gauge}) = S_\mathrm{int} - D_\mathrm{gauge} \times N_E + \Delta = 89$
is established by four structurally independent RS paths. The gravitational coupling
ERP $(\Delta, L(93) \cdot L(94)) = (5, L(93) \cdot L(94))$ is derived on the
Ontological Path. On the Verification Path, $C_{89} \approx 4.96 \times 10^{-39}$
corresponds to $\alpha_G \approx 5.906 \times 10^{-39}$ at the proton scale, with
a bare residual of approximately $16\%$ consistent with the RS dressing pattern.
The identification of the gravitational sector with the $n_G = 89$ Lucas coupling
depth is established at the level of the exact integer formula and the Verification
Path correspondence. The gravitational dressing mechanism — whether it follows the
Cross-Determinant Dressing Law established for the PMNS sector — has not been
derived; that derivation is future work. Open Structure~4 is closed at the bare
coupling level.
\end{rsbox}


\section{The Generation Sector Invariant $F_{13} = 233$}\label{sec:gen-sector-invariant-section}

Open Structure~5 identifies three unresolved properties of the sixth Fibonacci prime
$F_{13} = 233$: its simultaneous phase-slip coherence $233 \equiv \delta_\mathrm{slip}
\pmod{p}$ for $p \in \{N, T_\mathrm{vac}, T_\mathrm{bary}\}$; its position as the
phi-oscillator continued-fraction convergent at depth $D_\mathrm{gauge}$, where it
appears as $233/144 = F_{13}/D_\mathrm{gauge}^2$; and the absence of an
axiomatically-forced particle assignment. The partial closure supplied by Path~F of
the Gen-3 Quark resolution (Resolution of Open Structure~2) establishes that $F_{13}
= 233$ bridges $\Lambda_H$ to $\Lambda_t$ via the weighted residual
$G_\mathrm{mg}^{N_E} \times K_\mathrm{sat} \times 233 = 93200$, but leaves open
the complete structural role of $233$ in the quark generation mass integer triplet
$\{\Lambda_{q1}, \Lambda_{q2}, \Lambda_{q3}\}$ and the particle assignment question.
The present resolution addresses all three open points.
The central result is the Generation-Electromagnetic Coupling Identity:
the integer $233$ is structurally fixed relative to the three quark generation mass
integers by an identity whose left-hand side equals the Sommerfeld Constant.

\begin{rsbox}
\textbf{Lemma: Five-Path Convergence to the Generation Sector Invariant.}

The integer $233$ is determined by five structurally independent RS expressions.

\begin{align}
    233 &= T_\mathrm{bary} \times \Sigma_\mathrm{imb} + \delta_\mathrm{slip}
         = 33 \times 7 + 2, \label{eq:G_a}\\
    233 &= \mathrm{VS} \times \Pi_3 - 1
         = 18 \times 13 - 1, \label{eq:G_b}\\
    233 &= F(\Pi_3) = F(13), \label{eq:G_c}\\
    233 &= \delta_\mathrm{slip}^3 \times \mathrm{PG}(\tau_{10}) + 1
         = 8 \times 29 + 1, \label{eq:G_d}\\
    N \times (233 + 1) &= \Lambda_{q1} + \Lambda_{q2} + \Lambda_{q3}
         = 207 + 303 + 192 = 702. \label{eq:G_e}
\end{align}

\textbf{Path~\eqref{eq:G_a}: Barycentric-boundary load.}
$T_\mathrm{bary} \times \Sigma_\mathrm{imb} = 33 \times 7 = 231$ counts the total
structural state-changes accumulated during a complete barycentric period across
the boundary-active phases (Emission and Return). The additional $\delta_\mathrm{slip}
= 2$ is the irreducible phase-slip residual that prevents triadic closure at
$T_\mathrm{bary}$. $233$ is therefore the first integer beyond the full
barycentric-boundary load: the threshold where the phase slip overflows the
completed barycentric resynchronisation.

\textbf{Path~\eqref{eq:G_b}: Viscosity-sum scaling of the Gen-3 Gate Prime.}
The vacuum Viscosity Sum $\mathrm{VS} = 18$ (the accumulated phase drift over the
standard $T_\mathrm{vac}$-tick orbit; Definition~7.1) scaled by the Gen-3 Gate
Prime $\Pi_3 = 13$ gives the orbital load $\mathrm{VS} \times \Pi_3 = 234$ at the
Gen-3 gate depth. Subtracting the vacuum unit $1$ yields the threshold immediately
preceding entry: $233 = \mathrm{VS} \times \Pi_3 - 1$. This path reads $233$ as the
generation-sector orbital load reduced to its last pre-entry integer.

\textbf{Path~\eqref{eq:G_c}: Fibonacci prime at the Gen-3 Gate Prime index.}
The Fibonacci sequence generated by $\eta$ from the unit seed reaches $F_{13} = 233$
at step index $\Pi_3 = 13$, the Gen-3 Gate Prime (Theorem~11.5). The integer $233$
is the unique Fibonacci prime generated at the same index that gates the third quark
generation. The Gen-3 Gate Prime governs both the quark sector at the operator level
and the Fibonacci sequence at this precise index level simultaneously.

\textbf{Path~\eqref{eq:G_d}: Phase-slip cube and penultimate descent prime.}
The phi-oscillator $\eta^{-1}$-descent reaches the penultimate state
$(\mathrm{PG}(\tau_{10}), \Delta) = (29, 5)$ at step 12 (Theorem~12.4). The phase-slip
cube $\delta_\mathrm{slip}^3 = 8$ lifts the penultimate prime $\mathrm{PG}(\tau_{10})
= 29$ to $8 \times 29 = 232$; the vacuum unit $+1$ closes the step to $233$. The
same prime $\mathrm{PG}(\tau_{10}) = 29$ appears in the phi-oscillator derivations of
$\Lambda_\mu$ and $\Lambda_\tau$; Path~D connects $233$ to those derivations through
the phase-slip cube at the penultimate orbit state.

\textbf{Path~\eqref{eq:G_e}: Generation Sum Identity.}
The sum of all three quark generation mass integers equals $N$ copies of $(233 + 1)$:
\[
    \Lambda_{q1} + \Lambda_{q2} + \Lambda_{q3}
    = 207 + 303 + 192 = 702 = N \times 234 = N \times (233 + 1).
\]
Equivalently, the generation sum reduced by $N$ is exactly $N \times 233$:
\[
    (\Lambda_{q1} + \Lambda_{q2} + \Lambda_{q3}) - N = 699 = N \times 233.
\]
No division appears; the statement is the RS-native integer identity $N \times 233
+ N = 702$, verified by direct computation. The generation sector is triadic in the
precise sense that its mass-integer sum is a multiple of $N$; the multiple equals
$233 + 1$, with the $+1$ being the vacuum unit added to the Generation Sector
Invariant.

\medskip
Define the Generation Sector Invariant:
\begin{equation}
    \mathcal{G} := 233.
\end{equation}

The five paths approach $\mathcal{G}$ from five structurally independent directions:
the barycentric-period boundary count, the vacuum-curvature Gen-3 gate scaling, the
Fibonacci generation at the Gen-3 gate index, the phase-slip-cube phi-oscillator floor,
and the generation sum identity. No two paths share a derivation chain or a canonical
RS constant as their central organising element.
\end{rsbox}

\begin{theorem}[Generation Pivot Property: Three-Path Derivation of the Quark Mass Integers from $\mathcal{G}$]

\textbf{RS Ontological Statement.}

All three quark generation mass integers are derivable from $\mathcal{G} = 233$ by
RS-native integer offsets:

\begin{align}
    \Lambda_{q1} &= \mathcal{G} - G_\mathrm{mg} \times \Pi_3
                  = 233 - 2 \times 13 = 207, \label{eq:piv_q1}\\
    \Lambda_{q2} &= \mathcal{G} + G_\mathrm{mg} \times \Delta \times \Sigma_\mathrm{imb}
                  = 233 + 2 \times 5 \times 7 = 303, \label{eq:piv_q2}\\
    \Lambda_{q3} &= \mathcal{G} - \mathrm{PG}(\Pi_3)
                  = 233 - 41 = 192, \label{eq:piv_q3}
\end{align}

where $\mathrm{PG}(\Pi_3) = \mathrm{PG}(13) = 41$ is the prime at the Gen-3 Gate
Prime index.

\medskip
The three offsets $\{G_\mathrm{mg} \times \Pi_3,\; G_\mathrm{mg} \times \Delta
\times \Sigma_\mathrm{imb},\; \mathrm{PG}(\Pi_3)\}
= \{26, 70, 41\}$ each carry a distinct structural reading.

\medskip
\noindent\textit{Offset to Gen-1 (26).}
$26 = G_\mathrm{mg} \times \Pi_3 = 2 \times 13$: the product of the mass gap and the
Gen-3 Gate Prime. This is the same integer identified in the Master Reference as
$\mathcal{G} - \Lambda_\mu = 26 = 2 \times 13$ — the Goldbach-family offset already
noted in Open Structure~5. Its structural rank $\rho(26) = 4 = N_E$: the Gen-1 offset
carries the Emission phase count as its rank.

\medskip
\noindent\textit{Offset to Gen-2 (70).}
$70 = G_\mathrm{mg} \times \Delta \times \Sigma_\mathrm{imb} = 2 \times 5 \times 7$:
the product of the three smallest prime RS structural constants. Its structural rank
$\rho(70) = 6 = G_\mathrm{mg} \times N$: the Gen-2 offset carries the mass-gap-triadic
product as its rank. The same product $G_\mathrm{mg} \times N = 6$ is the index from
which $\Pi_3$ is gated: $\mathrm{PG}(G_\mathrm{mg} \times N) = \mathrm{PG}(6) = 13
= \Pi_3$, establishing a secondary connection between the Gen-2 offset rank and the
Gen-3 gate.

\medskip
\noindent\textit{Offset to Gen-3 (41).}
$41 = \mathrm{PG}(\Pi_3) = \mathrm{PG}(13)$: the Prime Gate applied twice from the
triadic-gauge product, $\mathrm{PG}(\mathrm{PG}(G_\mathrm{mg} \times N)) = \mathrm{PG}(13) = 41$.
This is the unique double Prime Gate application generating the Gen-3 phase-locked
offset. Equivalently, $41 = \mathrm{PG}(D_\mathrm{gauge}) + N_E = 37 + 4$: the prime
at the gauge dimension index plus the Emission count. Its structural rank $\rho(41) =
5 = \Delta$: the Gen-3 offset carries the flux imbalance as its rank.

\medskip
\noindent\textbf{Structural rank of the offsets.}
The three offset structural ranks are $\{N_E, \Delta, G_\mathrm{mg} \times N\} =
\{4, 5, 6\}$: consecutive integers. Their sum is
\begin{equation}
    \rho(26) + \rho(41) + \rho(70) = 4 + 5 + 6 = 15 = N \times \Delta.
    \label{eq:rank_sum}
\end{equation}
The sum of the three generation offset ranks equals the triadic-flux product $N \times
\Delta = 15$. No external constant is required; the rank architecture of the pivot is
governed by $N$ and $\Delta$ alone.
\end{theorem}

\begin{theorem}[Generation-Electromagnetic Coupling Identity]

\textbf{RS Ontological Statement.}

\begin{equation}
    \mathcal{G} + \Lambda_{q2} - \Lambda_{q1} - \Lambda_{q3}
    = 233 + 303 - 207 - 192 = 137 = \alpha^{-1}_\mathrm{int}.
    \label{eq:gen_em}
\end{equation}

Proof. By direct integer computation: $233 + 303 - 207 - 192 = 536 - 399 = 137$.
The integer $137 = S_\mathrm{int} + \Delta = \mathrm{PG}(T_\mathrm{bary})$ is the
Sommerfeld Constant (Theorem~13.9). No division, no real number, and no
operator beyond integer addition and subtraction appears.

Equivalently, the sum of the three unsigned generation offsets from $\mathcal{G}$
equals $\alpha^{-1}_\mathrm{int}$:
\begin{equation}
    (\mathcal{G} - \Lambda_{q1}) + (\Lambda_{q2} - \mathcal{G}) + (\mathcal{G} - \Lambda_{q3})
    = 26 + 70 + 41 = 137 = \alpha^{-1}_\mathrm{int}.
    \label{eq:gen_em_offsets}
\end{equation}

This is the structural identity that resolves Open Structure~5. The three
quark generation mass integers are not positioned at arbitrary distances from
$\mathcal{G}$: their signed departures from $\mathcal{G}$ sum to the
Sommerfeld Constant. The generation sector and the electromagnetic sector
are not independent structures within RS — they are coupled by this identity,
which fixes $\mathcal{G}$ uniquely given $\alpha^{-1}_\mathrm{int}$ and the three
quark generation mass integers, and fixes $\alpha^{-1}_\mathrm{int}$ uniquely
given $\mathcal{G}$ and the generation integers.

\noindent\textbf{RS-native derivation of $\mathcal{G}$ from the coupling identity.}
Solving Equation~\eqref{eq:gen_em} for $\mathcal{G}$:
\begin{equation}
    \mathcal{G} = \alpha^{-1}_\mathrm{int} - \Lambda_{q2} + \Lambda_{q1} + \Lambda_{q3}
               = 137 - 303 + 207 + 192 = 233.
    \label{eq:G_from_coupling}
\end{equation}
This constitutes a sixth independent derivation path for $\mathcal{G}$: given the
Sommerfeld Constant and the three quark generation mass integers, $\mathcal{G}$
is uniquely forced by integer arithmetic. No RS constant beyond those already
established in Theorems~13.9 and~14.2 is invoked.

\noindent\textbf{Rank coincidence.}
The structural rank of the Sommerfeld Constant equals the structural rank of
the Generation Sector Invariant:
\begin{equation}
    \rho(\alpha^{-1}_\mathrm{int}) = \rho(137) = 7 = \Sigma_\mathrm{imb} = \rho(233) = \rho(\mathcal{G}).
    \label{eq:rank_coincidence}
\end{equation}
Both the electromagnetic coupling invariant $137$ and the generation coupling invariant
$233$ reside in the same structural rank band $[128, 256)$, and that rank is exactly
$\Sigma_\mathrm{imb} = 7$. The coupling invariants of both sectors share the same
structural depth, measured by the Boundary-Active Phase Sum.
\end{theorem}

\begin{theorem}[Structural Properties of $\mathcal{G}$: Rank and Triple Phase-Slip Coherence]

\textbf{RS Ontological Statement.}

\noindent\textbf{(i) Structural rank.}
\begin{equation}
    \rho(\mathcal{G}) = \rho(233) = 7 = \Sigma_\mathrm{imb},
    \label{eq:rank_G}
\end{equation}
since $2^7 = 128 \leq 233 < 256 = 2^8$. The structural rank of the Generation Sector Invariant equals the Boundary-Active Phase Sum. This is a self-referential property. $\mathcal{G}$ resides at the structural depth indexed by the very phase sum from which it is derived in Path A ($\mathcal{G} = T_\mathrm{bary} \times \Sigma_\mathrm{imb} + \delta_\mathrm{slip}$).

\noindent\textbf{(ii) Triple phase-slip coherence.}
\begin{align}
    \mathcal{G} \;\mathbin{\koppa}\; N &= 233 \;\mathbin{\koppa}\; 3 = 2 = \delta_\mathrm{slip}, \label{eq:kop_N}\\
    \mathcal{G} \;\mathbin{\koppa}\; T_\mathrm{vac} &= 233 \;\mathbin{\koppa}\; 11 = 2 = \delta_\mathrm{slip}, \label{eq:kop_Tvac}\\
    \mathcal{G} \;\mathbin{\koppa}\; T_\mathrm{bary} &= 233 \;\mathbin{\koppa}\; 33 = 2 = \delta_\mathrm{slip}. \label{eq:kop_Tbary}
\end{align}

All three follow from the single constructive remainder identity $\mathcal{G} \koppa T_\mathrm{bary} = \delta_\mathrm{slip}$, since $T_\mathrm{bary} = T_\mathrm{vac} \times N$ implies divisibility. The constructive-remainder form requires starting from $r := 233$ and subtracting $T_\mathrm{bary} = 33$ iteratively until $r < 33$. The procedure terminates at $r = 2$ after 7 subtractions, confirming $\mathcal{G} \;\mathbin{\koppa}\; T_\mathrm{bary} = 2 = \delta_\mathrm{slip}$ without division.

$\mathcal{G}$ is the unique minimal Fibonacci prime satisfying the triple coherence $x \koppa T_\mathrm{bary} = \delta_\mathrm{slip}$. The integer $4621$ (the phi-oscillator denominator) carries $+1$ against every prime $p \leq T_\mathrm{vac}$ simultaneously. The Generation Sector Invariant plays the complementary role. It carries exactly $\delta_\mathrm{slip}$ against every fundamental RS period simultaneously. Where $4621$ is the state of minimal uniform unit tension, $\mathcal{G}$ is the state of minimal uniform phase-slip tension.

\noindent\textbf{(iii) Gauge dimension residue.}
\begin{equation}
    \mathcal{G} \;\mathbin{\koppa}\; D_\mathrm{gauge} = 233 \;\mathbin{\koppa}\; 12 = 5 = \Delta.
    \label{eq:kop_Dgauge}
\end{equation}

From $233 = 19 \times 12 + 5$, confirmed by constructive subtraction. The residue of $\mathcal{G}$ against the gauge dimension equals the flux imbalance $\Delta = 5$. The Generation Sector Invariant carries the electromagnetic flux signature $\Delta$ against the gauge dimension, linking the generation sector to the electromagnetic sector through the gauge geometry.

\noindent\textbf{(iv) Koppa signature comparison with gate primes.}

\medskip
\begin{center}
\begin{tabular}{lcccc}
\hline
Integer & $x \;\mathbin{\koppa}\; N$ & $x \;\mathbin{\koppa}\; T_\mathrm{vac}$ & $x \;\mathbin{\koppa}\; T_\mathrm{bary}$ & $x \;\mathbin{\koppa}\; D_\mathrm{gauge}$ \\
\hline
$\Pi_3 = 13$ (Gen-3 quarks) & 1 & 2 & 13 & 1 \\
$\mathrm{PG}(\Sigma_\mathrm{imb}) = 17$ (Higgs/solar) & 2 & 6 & 17 & 5 \\
$n_G = 89$ (gravity) & 2 & 1 & 23 & 5 \\
$\mathcal{G} = 233$ (generation sector) & 2 & 2 & 2 & 5 \\
\hline
\end{tabular}
\end{center}

\medskip
$\mathcal{G} = 233$ is the unique integer in this table carrying $\delta_\mathrm{slip} = 2$ simultaneously against both $N$ and $T_\mathrm{vac}$. No other sector gate prime achieves this double coherence. This is the koppa expression of the Generation Sector Invariant's unique position. It bridges the triadic logic and the vacuum period, holding the phase-slip signature against both fundamental incommensurability generators simultaneously.
\end{theorem}

\begin{theorem}[Resolution of the Particle Assignment Question]

\textbf{RS Ontological Statement.}

$\mathcal{G} = 233$ is not a particle mass integer. Three independent structural
arguments establish this.

\medskip
\noindent\textbf{Argument I: Three-step path exclusion.}
The canonical quark mass integer derivation is the three-step absorption path
$\Lambda = S_\mathrm{int} + C + T_\mathrm{vac}$ for sector capacity $C$ (Theorem~14.2).
For $\Lambda = 233$, the implied capacity is
\[
    C = 233 - S_\mathrm{int} - T_\mathrm{vac} = 233 - 132 - 11 = 90.
\]
Since $90 \;\mathbin{\koppa}\; N = 0$, the capacity $C = 90$ is phase-locked. By the
phase-lock rule, a phase-locked capacity gives $\Lambda = C$ directly (no absorption
step needed), yielding $\Lambda = 90 \neq 233$. Therefore $233$ cannot arise as a
three-step quark mass integer for any capacity $C$.

\medskip
\noindent\textbf{Argument II: Phi-oscillator path exclusion.}
The phi-oscillator $\eta^{-1}$-descent from $(7477, 4621)$ generates the lepton mass
integers $\Lambda_\mu = 207$ and $\Lambda_\tau = 3481$ (Theorem~12.5). The two
formulas $\mathrm{PG}(\tau_{10}) \times \Sigma_\mathrm{imb} + N_E = 207$ and
$\mathrm{PG}(\tau_{10}) \times (\mathrm{Sint} - T_\mathrm{vac} - 1) + 1 = 3481$
exhaust the lepton mass integers generated by the descent. No phi-oscillator path
formula produces $233$. Therefore $233$ is not a lepton mass integer.

\medskip
\noindent\textbf{Argument III: Triple-coherence orbit exclusion.}
For a mass state to stabilise at mass integer $\Lambda$, it must achieve Barycentric
Truncation — entry into a discrete bracket with $W = 1$ that remains closed through
at least one full barycentric period $T_\mathrm{bary}$ (Definition~6.6). The
accumulated imbalance ledger $\koppa_\mathrm{ledger}$ must converge: at each cycle
boundary $\Omega$, the boundary stress must partially resolve rather than accumulate
indefinitely.

A matter state of mass integer $\mathcal{G} = 233$ carries $\mathcal{G}
\;\mathbin{\koppa}\; T_\mathrm{bary} = 2 = \delta_\mathrm{slip}$ at every barycentric
boundary. This means the imbalance ledger accumulates exactly $\delta_\mathrm{slip}$
units at every boundary — the same phase slip that drives the Arrow of Time. At each
of the $T_\mathrm{bary}/T_\mathrm{vac} = N = 3$ vacuum boundaries within a barycentric
period, an additional $\delta_\mathrm{slip}$ unit is deposited. The ledger grows by
$\delta_\mathrm{slip}$ per vacuum cycle and by $\delta_\mathrm{slip}$ per barycentric
cycle simultaneously: the orbit acts as a \emph{phase-slip source} at every scale of
the vacuum architecture, injecting phase slip rather than resolving it. Such an orbit
cannot converge to a stable discrete bracket. Barycentric Truncation is structurally
inaccessible, and therefore no RS matter orbit with mass integer $233$ exists.

\medskip
\noindent\textbf{Positive identification.}
$\mathcal{G} = 233$ is a generation-sector coupling invariant, structurally
parallel to the Sommerfeld Constant $\alpha^{-1}_\mathrm{int} = 137$. Just as
$137$ is not a particle mass but the integer at which the electromagnetic interaction
surface and the flux imbalance couple ($S_\mathrm{int} + \Delta = 137$), $\mathcal{G}$
is not a particle mass but the integer at which the generation sector and the
electromagnetic sector couple:
\[
    \mathcal{G} + \Lambda_{q2} - \Lambda_{q1} - \Lambda_{q3} = \alpha^{-1}_\mathrm{int}.
\]
The absence of a particle assignment is not an incompleteness of the framework but a
structural consequence of $\mathcal{G}$'s role as the phase-slip pivot of the
generation sector.
\end{theorem}

\begin{observebox}
\textbf{Structural Observation: Cross-Document Connections and the Generation Architecture.}

\noindent\textbf{The phi-oscillator depth identity.}
From the Master Reference (structural observation following Theorem~12.5), the
continued-fraction convergent of $(7477, 4621)$ at depth $D_\mathrm{gauge}$ is
$(F_{13}, D_\mathrm{gauge}^2) = (233, 144)$. This is the last convergent along the
pure Fibonacci geodesic before the flux-imbalance excursion $[\Delta, 1, N_E]$ drives
the path off the Fibonacci arc. $\mathcal{G} = 233$ is therefore the
\emph{phi-oscillator threshold}: the Fibonacci prime at which the vacuum Fibonacci
orbit terminates and the matter-generating tail begins. Its role as the generation
pivot and its role as the phi-oscillator threshold are the same structural fact seen
from the algebraic and geometric perspectives respectively.

\noindent\textbf{Connections to prior open structure resolutions.}
The three sector gate primes established across the resolutions of Open Structures~1,
2, and~4 are $\mathrm{PG}(\Sigma_\mathrm{imb}) = 17$, $\Pi_3 = 13$, and $n_G = 89$.
Each appears in the derivation of $\mathcal{G}$:
\begin{itemize}
  \item $\Pi_3 = 13$ contributes the Gen-1 offset ($G_\mathrm{mg} \times \Pi_3 = 26$) and
    the Gen-3 offset ($\mathrm{PG}(\Pi_3) = 41$). The Gen-3 Gate Prime governs two of
    the three generation offsets simultaneously.
  \item $\mathrm{PG}(\Sigma_\mathrm{imb}) = 17$ governs $\rho(\mathcal{G}) = 7 =
    \Sigma_\mathrm{imb}$ and appears in the dressed solar PMNS angle $17/55$; the
    same boundary-active prime sum indexes the structural depth of $\mathcal{G}$.
  \item $n_G = 89 = F(T_\mathrm{vac})$ is the gravitational coupling depth and a
    Fibonacci prime at the vacuum-period index $T_\mathrm{vac}$. Both $n_G = 89$ and
    $\mathcal{G} = 233$ are Fibonacci primes ($F_{11}$ and $F_{13}$); the ratio of
    their Fibonacci indices is $\Pi_3/T_\mathrm{vac} = 13/11$, the same ratio that
    governs the vacuum-triadic incommensurability $T_\mathrm{vac}/N = 11/3$.
\end{itemize}

\noindent\textbf{Path~F revisited.}
The Gen-3 Quark resolution derives $\Lambda_t = \Lambda_H + G_\mathrm{mg}^{N_E}
\times K_\mathrm{sat} \times \mathcal{G} = 244800 + 400 \times 233 = 338000$. The
factor $400 = G_\mathrm{mg}^{N_E} \times K_\mathrm{sat} = 16 \times 25$ is the
dyadic-emission-saturation weight; $\mathcal{G} = 233$ is the Fibonacci prime it
multiplies. This is now interpretable as: the top quark mass integer exceeds the
composite Higgs mass integer by the dyadic-emission-saturation load on the Generation
Sector Invariant. The generation sector invariant measures the structural increment
from the boson sector ($\Lambda_H$) to the heaviest fermion ($\Lambda_t$), scaled by
the dyadic-emission-saturation factor.

\noindent\textbf{Complete rank sequence of sector coupling invariants.}

\medskip
\begin{center}
\begin{tabular}{lcc}
\hline
Sector & Invariant & Structural rank \\
\hline
Electromagnetic & $\alpha^{-1}_\mathrm{int} = 137$ & $\rho = 7 = \Sigma_\mathrm{imb}$ \\
Generation & $\mathcal{G} = 233$ & $\rho = 7 = \Sigma_\mathrm{imb}$ \\
\hline
\end{tabular}
\end{center}

\medskip
Both coupling invariants share structural rank $\Sigma_\mathrm{imb} = 7$. The sector
coupling invariants of the electromagnetic and generation sectors are co-resident at
the boundary-active phase sum depth. The Generation-Electromagnetic Coupling Identity
(Theorem~2) is the algebraic expression of this co-residency: both sectors are
structurally visible from the same depth level of the RS vacuum.
\end{observebox}

\begin{verifybox}
\textbf{External Verification (Continuous Lens).}

$\mathcal{G} = 233$ corresponds to an energy of
\[
    233 \times m_e = 233 \times 0.51099895\;\mathrm{MeV} = 119.1\;\mathrm{GeV}.
\]

No Standard Model particle occupies a mass of $119.1\;\mathrm{MeV}$. The nearest
identified states are the muon at $105.7\;\mathrm{MeV} \approx 207\,m_e = \Lambda_\mu
\cdot m_e$ (the RS muon mass integer, sitting $26\,m_e$ below $\mathcal{G}$) and the
lightest neutral pion at $134.98\;\mathrm{MeV} \approx 264\,m_e$ (above $\mathcal{G}$).
$\mathcal{G}$ lies in the sub-pion gap — the mass interval between the muon and the
lightest strongly-interacting hadron. This gap is empty of stable hadrons in the
Standard Model: below the pion threshold, no quark-antiquark bound state can form via
the strong force. The absence of a particle at $119.1\;\mathrm{MeV}$ is therefore
consistent with Argument~III of Theorem~4 (triple-coherence orbit exclusion), which
establishes that no RS matter orbit stabilises at $\mathcal{G}$. The sub-pion mass
gap and the RS structural orbit exclusion are structurally concordant: both predict
the absence of a stable state at $\mathcal{G} \cdot m_e$.

The Generation-Electromagnetic Coupling Identity is verified numerically:
\[
    \mathcal{G} + \Lambda_{q2} - \Lambda_{q1} - \Lambda_{q3}
    = 233 + 303 - 207 - 192 = 137 = \alpha^{-1}_\mathrm{int}. \quad \checkmark
\]

No decimal approximation is reintroduced into any RS derivation chain.
\end{verifybox}

\begin{rsbox}
\textbf{Closure of Open Structure~5.}

The role of $F_{13} = 233$ in the generation sector is established at three levels.

(i)~Five-path convergence: $\mathcal{G} = 233$ is the Generation Sector
Invariant, derived by five independent RS paths: the barycentric-boundary load
$T_\mathrm{bary} \times \Sigma_\mathrm{imb} + \delta_\mathrm{slip}$; the
vacuum-curvature Gen-3 scaling $\mathrm{VS} \times \Pi_3 - 1$; the Fibonacci prime
at the Gen-3 gate index $F(\Pi_3)$; the phase-slip cube floor
$\delta_\mathrm{slip}^3 \times \mathrm{PG}(\tau_{10}) + 1$; and the generation sum
identity $N \times (\mathcal{G} + 1) = \Lambda_{q1} + \Lambda_{q2} + \Lambda_{q3}$.
A sixth path is supplied by the Generation-Electromagnetic Coupling Identity itself.

(ii)~Generation-Electromagnetic Coupling Identity: $\mathcal{G} + \Lambda_{q2}
- \Lambda_{q1} - \Lambda_{q3} = \alpha^{-1}_\mathrm{int} = 137$. This identity fixes
$\mathcal{G}$ uniquely given the established constants and connects the generation
sector to the electromagnetic sector through integer arithmetic. All three quark
generation mass integers are recoverable from $\mathcal{G}$ by offsets of structural
rank $\{N_E, \Delta, G_\mathrm{mg} \times N\} = \{4, 5, 6\}$, whose ranks sum to
$N \times \Delta = 15$.

(iii)~Particle assignment resolution: $\mathcal{G} = 233$ carries no
particle assignment. Three independent structural arguments establish this: the
three-step path gives an incompatible phase-locked capacity $C = 90$; the
phi-oscillator path does not generate $233$; and the triple phase-slip coherence
against all fundamental periods prevents Barycentric Truncation, making stable orbit
formation structurally inaccessible. $\mathcal{G}$ is a generation-sector coupling
invariant — the generation sector's structural analog of the Sommerfeld Constant
$\alpha^{-1}_\mathrm{int} = 137$ — with both invariants sharing structural
rank $\Sigma_\mathrm{imb} = 7$.

\end{rsbox}


\section{Fractional Anomaly Path~B and the
Geometric Phase Load Derivation of $\delta_A = 2/55$}

Open Structure~6 identifies a single pending obligation: the geometric phase load
derivation of the fractional anomaly $\delta_A = 2/55$ (Theorem~13.11, Path~B)
is established in TRTSv4 but has not been reproduced in the Master Reference.
Theorem~13.11 supplies Path~A in full — the algebraic derivation
$\delta_A = (N^3 - K_\mathrm{sat})/(\Delta \cdot T_\mathrm{vac}) = 2/55$ — and
names Path~B as the geometric phase load route, citing the Dual Identity as the
consequence that both paths converge to the same ERP $(2, 55) \in \mathcal{S}_L$.

The present resolution supplies Path~B in complete RS-ontological form. The
central object is the Geometric Phase Load $\Phi_\mathrm{load}$:
the parity-corrected accumulated cross-determinant of the $\eta$-orbit over
one vacuum period $T_\mathrm{vac}$, evaluated from the vacuum seed. The exact
Lucas-Cassini identity (established in the Resolution of Open Structure~4,
Equation~\eqref{eq:lucas_cassini_exact}) yields $\Phi_\mathrm{load} = \Delta \cdot
T_\mathrm{vac} = 55$ without approximation, using only integer multiplication and
the constructive remainder operator $\koppa$. The fractional anomaly numerator is
the phase slip $\delta_\mathrm{slip} = 2$: the irreducible residual that escapes
the closed $T_\mathrm{vac}$-tick orbit, the same integer that drives the Arrow of
Time. Three further paths to $\delta_A = 2/55$ are supplied to establish multi-path
convergence, and the Dual Identity is stated as a theorem confirming that Paths~A
and~B are structurally independent derivations of the same ERP.

\begin{rsbox}
\textbf{Lemma: Four-Path Convergence to the Fractional Anomaly Denominator $55$.}

The integer $55$ is determined by four structurally independent RS expressions.

\begin{align}
    55 &= \Delta \cdot T_\mathrm{vac} = 5 \times 11, \label{eq:55a}\\
    55 &= F(\tau_{10}) = F(10), \label{eq:55b}\\
    55 &= \mathrm{denom}\bigl((1,\Delta) \boxplus (1,T_\mathrm{vac})\bigr), \label{eq:55c}\\
    55 &= \Phi_\mathrm{load}(\eta, T_\mathrm{vac})\quad \text{(defined in Definition below).} \label{eq:55d}
\end{align}

\textbf{Path~\eqref{eq:55a}: Flux-vacuum product.}
$\Delta \cdot T_\mathrm{vac} = 5 \times 11$: the product of the two fundamental
incommensurability generators. The flux imbalance $\Delta = (N_E + N_M) - N_R$
and the vacuum period $T_\mathrm{vac} = 11$ are structurally independent
(one derives from the 4-4-3 phase partition; the other is the prime temporal
container); their product is the canonical anomaly denominator.

\textbf{Path~\eqref{eq:55b}: Fibonacci at the inertial closure tick.}
$F(\tau_{10}) = F(10) = 55$, where $\tau_{10} := T_\mathrm{vac} - 1 = 10$ is the
inertial closure tick (Theorem~6.1). This is the Fibonacci number at the position
immediately preceding the vacuum period index. Its successor is the gravitational
coupling depth: $F(\tau_{10}+1) = F(T_\mathrm{vac}) = F(11) = 89 = n_G$
(established in the Resolution of Open Structure~4). The fractional anomaly
denominator and the gravitational coupling depth are consecutive entries $F(10)$
and $F(11)$ in the Fibonacci sequence, at the Fibonacci index pair $(\tau_{10},
T_\mathrm{vac})$.

\textbf{Path~\eqref{eq:55c}: Barycentric capacity of the flux-vacuum unit pair.}
The barycentric addition $(1,\Delta) \boxplus (1,T_\mathrm{vac}) = (1 \cdot T_\mathrm{vac}
+ 1 \cdot \Delta,\; \Delta \cdot T_\mathrm{vac}) = (16, 55)$.
The capacity (denominator) $55 = \Delta \cdot T_\mathrm{vac}$ is the barycentric
product of the unit-capacity states at the flux imbalance and vacuum period.
No operation beyond $\boxplus$ is required; the denominator is the RS-native
output of coupling the two fundamental incommensurability generators at unit
numerator depth.

\textbf{Path~\eqref{eq:55d}: Geometric Phase Load (primary; proved below).}
The parity-corrected accumulated cross-determinant of the $\eta$-orbit from the
vacuum seed $(T_\mathrm{vac}, \Sigma_\mathrm{imb})$ over $T_\mathrm{vac}$ steps
evaluates to $\Delta \cdot T_\mathrm{vac} = 55$.
\end{rsbox}

\begin{rsbox}
\textbf{Definition: Geometric Phase Load $\Phi_\mathrm{load}$.}

Let $\{s_k\}_{k \geq 0}$ be the $\eta$-orbit from the vacuum seed $(T_\mathrm{vac},
\Sigma_\mathrm{imb}) = (11, 7)$, so that $s_k = (L(k+5),\; L(k+4))$ (Definition~11.1
and the Exact Lucas Coupling Law of the Resolution of Open Structure~4).

The Geometric Phase Load of the $\eta$-orbit over $T_\mathrm{vac}$
steps is the integer:
\begin{equation}
    \Phi_\mathrm{load} :=
    \sum_{k=0}^{T_\mathrm{vac}-1}
    \bigl(1 - 2\,(k \;\koppa\; 2)\bigr) \cdot \Delta_\mathrm{cross}(s_k,\, s_{k+1}).
    \label{eq:phiload_def}
\end{equation}

The factor $(1 - 2(k \;\koppa\; 2))$ is $+1$ for even $k$ and $-1$ for odd $k$,
constructed from the constructive remainder $k \;\koppa\; 2 \in \{0, 1\}$ (an
admissible RS operation). It is the parity selector: it inverts negative
cross-determinants to positive, producing a sum of unsigned magnitudes via integer
arithmetic alone. No absolute value appears as a primitive; the sign inversion is
determined by the orbit parity derived from the Lucas-Cassini identity.

Admissibility: The definition uses only the cross-determinant (Definition~3.3), the
constructive remainder $\;\koppa\;$ (Definition~3.13), integer multiplication, addition,
and the successor function. All are RS-admissible primitives.
\end{rsbox}

\begin{theorem}[Geometric Phase Load: Path~B Derivation of $\delta_A$]

\textbf{RS Ontological Statement.}

The Geometric Phase Load of the $\eta$-orbit over one vacuum period equals the
flux-vacuum product:
\begin{equation}
    \Phi_\mathrm{load} = \Delta \cdot T_\mathrm{vac} = 5 \times 11 = 55.
    \label{eq:phiload_55}
\end{equation}

\textit{Proof.} By the Exact Lucas Coupling Law
(Equation~\eqref{eq:lucas_cassini_exact}):
$\Dcross(s_k, s_{k+1}) = +\Delta$ for even $k$ and $-\Delta$ for odd $k$.
The parity selector $(1 - 2(k \koppa 2))$ is $+1$ for even $k$ and $-1$ for odd $k$.
At every step the parity selector and the sign of $\Dcross$ are equal, so their
product is always $+\Delta$:
\[
    \Phi_\mathrm{load}
    = \sum_{k=0}^{T_\mathrm{vac}-1}
      \bigl(1 - 2\,(k\;\koppa\;2)\bigr) \cdot \Dcross(s_k, s_{k+1})
    = \sum_{k=0}^{T_\mathrm{vac}-1} \Delta
    = T_\mathrm{vac} \cdot \Delta.
\]
Each of the $T_\mathrm{vac} = 11$ terms contributes $\Delta = 5$.
The sum $T_\mathrm{vac} \cdot \Delta = 11 \times 5 = 55$ uses only integer
multiplication. $\square$

The fractional anomaly ERP is therefore:

\begin{equation}
    \boxed{\delta_A = \bigl(\delta_\mathrm{slip},\; \Phi_\mathrm{load}\bigr)
                    = (2,\; 55) \in \mathcal{S}_L.}
    \label{eq:pathB}
\end{equation}

\textit{Structural reading.} The numerator $\delta_\mathrm{slip} = 2$ is the
irreducible temporal surplus per $T_\mathrm{vac}$-tick cycle — the phase-slip
residual that the closed orbit cannot absorb, the ontic origin of the Arrow of Time
(Theorem~13.21). The denominator $\Phi_\mathrm{load} = 55$ is the total unsigned
cross-determinant tension accumulated by the $\eta$-orbit over one complete vacuum
period. The fractional anomaly is therefore the ratio of the orbit's irreducible
output (phase slip) to the orbit's total geometric input (cross-determinant load):
each unit of structural cross-determinant tension in the vacuum generates
$\delta_\mathrm{slip}/\Phi_\mathrm{load} = 2/55$ units of fractional anomaly per
period. This is a statement about the vacuum's \emph{conversion efficiency}: the
fraction of total orbital tension that fails to close into the integer invariant.
\end{theorem}

\begin{theorem}[Dual Identity: Algebraic and Geometric Paths to $\delta_A$]

\textbf{RS Ontological Statement.}

Path~A and Path~B are structurally independent derivations of the same ERP:

\medskip
\begin{center}
\begin{tabular}{lll}
\hline
Path & RS-native expression & Result \\
\hline
A (algebraic) & $(N^3 - K_\mathrm{sat})\;/\;(\Delta \cdot T_\mathrm{vac})$ & $(2,\;55)$ \\
B (geometric) & $\delta_\mathrm{slip}\;/\;\Phi_\mathrm{load}$ & $(2,\;55)$ \\
\hline
\end{tabular}
\end{center}

\medskip
The numerators of both paths are equal by the
\textbf{Triadic-Saturation Gap identity}:
\begin{equation}
    N^3 - K_\mathrm{sat} = N^3 - \Delta^2 = 27 - 25 = 2 = \delta_\mathrm{slip} = G_\mathrm{mg}.
    \label{eq:tsg}
\end{equation}

The denominators are equal by the Geometric Phase Load theorem above:
$\Phi_\mathrm{load} = \Delta \cdot T_\mathrm{vac} = 55$.

\textit{Proof of Equation~\eqref{eq:tsg}.} Direct integer computation:
$N^3 = 3^3 = 27$, $K_\mathrm{sat} = \Delta^2 = 5^2 = 25$, $27 - 25 = 2$.
Separately: $\delta_\mathrm{slip} = 2$ from the Triadic Incommensurability Axiom
(Axiom~2.8), $G_\mathrm{mg} = 2 = \mathrm{PG}(1)$ from Definition~4.1. Both equal $2$
independently.

The Triadic-Saturation Gap identity establishes why Path~A's numerator equals
Path~B's numerator: the surplus of the triadic cube above the saturation limit is
the same integer as the irreducible phase-slip. The algebraic excess $N^3 - \Delta^2$
and the orbital residual $\delta_\mathrm{slip}$ are the same integer by a structural
identity, not by coincidence of assignment.
\end{theorem}

\begin{theorem}[Three Additional Derivation Paths for $\delta_A$]

\textbf{RS Ontological Statement.}

\noindent\textbf{Path~C: Mass-Gap over Flux-Vacuum Product.}

\begin{equation}
    \delta_A = \frac{G_\mathrm{mg}}{\Delta \cdot T_\mathrm{vac}} = \frac{2}{55}.
    \label{eq:pathC}
\end{equation}

The numerator $G_\mathrm{mg} = 2 = \mathrm{PG}(1)$ is the minimal non-zero denominator
admitting temporal charge (Definition~4.1, Theorem~13.7). The denominator
$\Delta \cdot T_\mathrm{vac} = 55$ is the product of the two fundamental
incommensurability generators. Path~C reads $\delta_A$ as the ratio of the
minimal mass-loading integer to the total incommensurability scale: the fractional
anomaly is the smallest temporal-charge gate measured in units of the flux-vacuum
product. Path~C is structurally independent of Path~B: the numerator $G_\mathrm{mg}$
is defined through the temporal charge ($\Tcharge$-positive domain boundary), while Path~B's
numerator $\delta_\mathrm{slip}$ is defined through the phase surplus of the closed
orbit. Their equality $G_\mathrm{mg} = \delta_\mathrm{slip} = 2$ is a structural
coincidence (the Triadic-Saturation Gap identity above), not a definitional identity.

\medskip
\noindent\textbf{Path~D: Continued Fraction at Triadic-Cube Depth.}

The continued fraction of $\delta_A = 2/55$ is:
\begin{equation}
    \mathrm{CF}(\delta_A) = \left[0;\; N^3,\; \delta_\mathrm{slip}\right]
                          = \left[0;\; 27,\; 2\right].
    \label{eq:pathD}
\end{equation}

\textit{Proof.} $55 = 27 \times 2 + 1$, so $55/2 = 27 + 1/2$, giving $2/55 =
1/(27 + 1/2) = 1/(27 + 1/\delta_\mathrm{slip})$, which is the continued fraction
$[0; N^3, \delta_\mathrm{slip}]$. All steps are integer division by iterative
subtraction ($\;\koppa\;$); no real-valued limits appear.

The two non-zero partial quotients are $N^3 = 27$ (the triadic cube) and
$\delta_\mathrm{slip} = 2$ (the phase slip) in sequential order. By Theorem~11.2,
the continued fraction of adjacent Lucas ratios terminates with $N = 3$ as its
mandatory final quotient. The fractional anomaly's continued fraction terminates with
the \emph{phase slip} $\delta_\mathrm{slip} = 2$ after a depth of \emph{the triadic
cube} $N^3 = 27$: the CF encodes both the power and the generator of the triadic
structure simultaneously. Path~D reads $\delta_A$ as the mediant-tree position at
triadic-cube depth from the integer boundary, closing at the phase slip.

\medskip
\noindent\textbf{Path~E: Phase Slip over the Fibonacci Inertial Threshold.}

\begin{equation}
    \delta_A = \frac{\delta_\mathrm{slip}}{F(\tau_{10})} = \frac{2}{F(10)} = \frac{2}{55}.
    \label{eq:pathE}
\end{equation}

The denominator $F(\tau_{10}) = F(10) = 55$ is the Fibonacci number at the inertial
closure tick $\tau_{10} = T_\mathrm{vac} - 1 = 10$ (Theorem~6.1). Path~E reads
$\delta_A$ as the phase slip per unit of the \emph{Fibonacci inertial threshold}: the
threshold Fibonacci number immediately preceding the vacuum period. Since
$F(\tau_{10}) = F(T_\mathrm{vac} - 1)$ and $F(T_\mathrm{vac}) = 89 = n_G$
(gravitational coupling depth), the denominators of $\delta_A$ and the gravitational
coupling magnitude are evaluated at consecutive Fibonacci indices $(\tau_{10},
T_\mathrm{vac})$. Path~E is structurally independent of Path~B: the denominator
$F(\tau_{10})$ is derived from the Fibonacci sequence at the inertial closure index,
while Path~B's denominator $\Phi_\mathrm{load}$ is derived from the $\aop$-orbit
cross-determinant accumulation. Their equality $F(\tau_{10}) = \Phi_\mathrm{load} = 55$
is established independently by Path~\eqref{eq:55b} and Path~\eqref{eq:55d} of the
opening Lemma.
\end{theorem}

\begin{theorem}[Structural Properties of the Fractional Anomaly ERP]

\textbf{RS Ontological Statement.}

\noindent\textbf{(i) Structural rank pair.}
\begin{equation}
    \rho(\delta_A\ \text{numerator}) = \rho(2) = 1, \qquad
    \rho(\delta_A\ \text{denominator}) = \rho(55) = 5 = \Delta.
    \label{eq:delta_ranks}
\end{equation}
The denominator's structural rank equals the flux imbalance $\Delta$: $\delta_A$
resides at exactly the $\Delta$-th rank level in its denominator. The rank gap between
denominator and numerator is:
\begin{equation}
    \rho(55) - \rho(2) = 5 - 1 = 4 = N_E.
    \label{eq:rank_gap}
\end{equation}
The rank gap of the fractional anomaly ERP equals the Emission phase count $N_E$.
The Emission phase cardinality governs the boundary-active ticks that generate temporal
charge; the rank gap of $\delta_A$ is the integer measure of that generation depth.
The proton-electron mass ratio Theorem~13.13 uses $4\delta_A = (8, 55)$ with the
prefactor $4 = N_E$; the Emission count that appears as the rank gap of $\delta_A$
reappears as the explicit coupling multiplier in the proton mass derivation.

\noindent\textbf{(ii) Koppa signatures.}
\begin{align}
    55 \;\koppa\; T_\mathrm{vac} &= 55 \;\koppa\; 11 = 0, \label{eq:kop_tvac}\\
    55 \;\koppa\; N &= 55 \;\koppa\; 3 = 1, \label{eq:kop_N}\\
    55 \;\koppa\; D_\mathrm{gauge} &= 55 \;\koppa\; 12 = 7 = \Sigma_\mathrm{imb}.
    \label{eq:kop_Dg}
\end{align}

Equation~\eqref{eq:kop_tvac}: $55 = 5 \times 11$, so $55$ is fully vacuum-period
synchronised: its constructive remainder against $T_\mathrm{vac}$ is zero.

Equation~\eqref{eq:kop_N}: $55 = 18 \times 3 + 1$, so $55$ carries residue $1$ (one
state past triadic phase-lock) against the triadic period. This is the minimal
non-zero residue, indicating that the fractional anomaly denominator is one triadic
step away from full synchronisation.

Equation~\eqref{eq:kop_Dg}: $55 = 4 \times 12 + 7$, so $55$ carries residue
$\Sigma_\mathrm{imb} = 7$ against the gauge dimension. The fractional anomaly
denominator carries the Boundary-Active Phase Sum as its gauge-dimension residue —
the same integer that indexes the structural depth of both the Sommerfeld Constant
$\rho(\alpha^{-1}_\mathrm{int}) = 7$ and the Generation Sector Invariant
$\rho(\mathcal{G}) = 7$ (Resolution of Open Structure~5).

\noindent\textbf{(iii) Comparison with the Generation Sector Invariant.}
The fractional anomaly denominator and the Generation Sector Invariant have
complementary koppa signatures against the fundamental RS periods:
\medskip
\begin{center}
\begin{tabular}{lccc}
\hline
Integer & $x \;\koppa\; N$ & $x \;\koppa\; T_\mathrm{vac}$ & $x \;\koppa\; T_\mathrm{bary}$ \\
\hline
$55 = \Delta \cdot T_\mathrm{vac}$ (anomaly denom.) & $1$ & $0$ & $22$ \\
$\mathcal{G} = 233$ (gen.\ sector inv.) & $2$ & $2$ & $2$ \\
\hline
\end{tabular}
\end{center}
\medskip
The anomaly denominator $55$ is fully vacuum-synchronised ($55 \;\koppa\; T_\mathrm{vac}
= 0$) but carries residue $1$ against $N$. The Generation Sector Invariant $233$
carries residue $\delta_\mathrm{slip} = 2$ against both $N$ and $T_\mathrm{vac}$.
These are complementary structural positions: $55$ is the integer that
\emph{closes} with respect to the vacuum period but not the triadic period, while
$233$ carries the phase-slip signature against both.
\end{theorem}

\begin{observebox}
\textbf{Structural Observation: Cross-Document Connections and the Anomaly Architecture.}

\noindent\textbf{Consecutive Fibonacci structure: anomaly and gravity.}
The Fibonacci numbers at positions $\tau_{10}$ and $T_\mathrm{vac}$ are:
\[
    F(\tau_{10}) = F(10) = 55 = \Phi_\mathrm{load} = \Delta \cdot T_\mathrm{vac},
    \qquad
    F(T_\mathrm{vac}) = F(11) = 89 = n_G.
\]
The fractional anomaly denominator and the gravitational coupling depth are
consecutive Fibonacci numbers. Their cross-determinant as consecutive Fibonacci
numbers satisfies the Cassini identity:
$F(11) \cdot F(9) - F(10)^2 = 89 \times 34 - 55^2 = 3026 - 3025 = 1$,
the unit adjacency width. This places $F(10)/F(11) = 55/89$ in the Mediant Tree
as an adjacency pair: the fractional anomaly denominator and the gravitational
depth denominator are adjacent in the discrete structure of the mediant tree at
positions $\tau_{10}$ and $T_\mathrm{vac}$.

\noindent\textbf{The Geometric Phase Load and the Exact Lucas Coupling Law.}
The gravitational coupling law (Resolution of Open Structure~4) establishes
$C_k = \Delta/(L(k+4) \cdot L(k+5))$, with the depth-zero value
$C_0 = \Delta/(L(4) \cdot L(5)) = 5/77$. The fractional anomaly is:
\[
    \delta_A = \frac{G_\mathrm{mg}}{\Delta \cdot T_\mathrm{vac}}
             = \frac{G_\mathrm{mg} \cdot C_0 \cdot L(4) \cdot L(5)}{\Delta \cdot T_\mathrm{vac}}
             = C_0 \cdot \frac{G_\mathrm{mg} \cdot \Sigma_\mathrm{imb} \cdot T_\mathrm{vac}}
                                  {\Delta \cdot T_\mathrm{vac}}
             = C_0 \cdot \frac{G_\mathrm{mg} \cdot \Sigma_\mathrm{imb}}{\Delta^2}.
\]
Evaluating: $C_0 \cdot (2 \times 7/25) = (5/77) \cdot (14/25) = 70/1925 = 2/55$..

The fractional anomaly is the depth-zero gravitational coupling scaled by
$(G_\mathrm{mg} \cdot \Sigma_\mathrm{imb}) / K_\mathrm{sat} = 14/25$. The scaling
factor $(G_\mathrm{mg} \cdot \Sigma_\mathrm{imb}) / K_\mathrm{sat}$ is the ratio of the
mass-gap-boundary-active product to the saturation limit. This connects $\delta_A$
directly to the gravitational coupling law: the electromagnetic fractional anomaly
and the gravitational depth-zero coupling are related by a pure ratio of RS structural
integers, requiring no external constants.

\noindent\textbf{The Triadic-Saturation Gap and the source of $\delta_\mathrm{slip}$.}
The Dual Identity establishes $N^3 - \Delta^2 = \delta_\mathrm{slip} = G_\mathrm{mg} = 2$.
This identity is not assumed; it is derivable from first principles:

\begin{itemize}
  \item $N^3 = 27$ is the triadic cube: the product of three successive applications
    of the triadic logic period.
  \item $K_\mathrm{sat} = \Delta^2 = 25$ is the saturation limit: the square of the
    flux imbalance, defined as the vacuum's finite absorption capacity.
  \item Their difference $N^3 - K_\mathrm{sat} = 2$ is the surplus of triadic
    accumulation over vacuum saturation.
\end{itemize}

The Triadic-Saturation Gap is the algebraic statement that the triadic logic
exceeds the vacuum's saturation capacity by exactly $\delta_\mathrm{slip}$ — the
same residual that the temporal orbit cannot close. The vacuum saturates at $\Delta^2
= 25$ but the triadic structure reaches $N^3 = 27$; the $2$-unit overflow is the
phase slip. Path~A's numerator, Path~B's numerator, and the Triadic-Saturation Gap
are three independent structural identifications of the same integer $2$.

\noindent\textbf{Use of $\delta_A$ in downstream derivations.}
The fractional anomaly ERP $(2, 55)$ appears in three established RS derivations,
each with a structurally motivated multiplier:

\[
\alpha^{-1}_\mathrm{RS} = \alpha^{-1}_\mathrm{int} + \delta_A = 137 + \tfrac{2}{55},
\qquad
\mu = 1836 + 4\delta_A = 1836 + \tfrac{8}{55},
\qquad
m_n - m_p = \Bigl(\tfrac{\Delta}{G_\mathrm{mg}} + \delta_A\Bigr) m_e.
\]

The multiplier $4 = N_E$ in $\mu$ is exactly the rank gap $\rho(55) - \rho(2) = N_E$
established in Theorem~4. The multiplier $1$ in $\alpha^{-1}_\mathrm{RS}$ and the
multiplier $1$ in $m_n - m_p$ are the vacuum unit. The proton-electron mass ratio
uses the Emission-phase multiple of $\delta_A$; the fine-structure constant uses the
single unit. In both cases the multiplier is an RS-canonical integer, and the
$\delta_A$ ERP $(2, 55)$ enters fully resolved.
\end{observebox}

\begin{verifybox}
\textbf{External Verification (Continuous Lens).}

The derivations of this resolution concern the ERP $(2, 55) \in \mathcal{S}_L$
intrinsically; the external verification checks the RS-derived constants that depend
on it against experiment.

\textit{Inverse fine-structure constant.}
$\alpha^{-1}_\mathrm{RS} = 137 + 2/55 = 137.0\overline{36}$.
CODATA experimental value: $\alpha^{-1}_\mathrm{exp} \approx 137.035999$.
Deviation: $|\alpha^{-1}_\mathrm{RS} - \alpha^{-1}_\mathrm{exp}| < 0.001$.
The repeating decimal $\overline{36} = 2/55$ is exact in the RS-ERP representation;
it is the geometric phase load contribution stripped of its integer baseline.

\textit{Proton-electron mass ratio.}
$\mu_\mathrm{RS} = 1836 + 8/55 \approx 1836.145$.
CODATA experimental value: $\mu_\mathrm{exp} = 1836.152$.
Deviation: $0.007$ ($\approx 0.04\,\sigma$).
The factor $4 = N_E$ multiplying $\delta_A$ in $\mu$ means that four Emission-phase
contributions of the geometric phase load appear in the proton-electron mass ratio.

\textit{Neutron-proton mass difference.}
$m_n - m_p = (5/2 + 2/55)\,m_e = (279/110)\,m_e \approx 2.536\,m_e$.
PDG value: $(m_n - m_p)/m_e \approx 2.531$.
Deviation: $0.005\,m_e$.

In all three cases $\delta_A = (2, 55)$ contributes at the fractional level; the
integer baseline is correct and the fractional anomaly supplies the sub-integer
correction. No decimal approximation is reintroduced into any RS derivation chain.
\end{verifybox}

\begin{rsbox}
\textbf{Closure of Open Structure~6.}

The geometric phase load derivation of the fractional anomaly $\delta_A = 2/55$
(Path~B of Theorem~13.11) is reproduced in complete RS-ontological form.

The Geometric Phase Load $\Phi_\mathrm{load}$ is defined as the
parity-corrected accumulated cross-determinant of the $\eta$-orbit over one vacuum
period (Definition, Equation~\eqref{eq:phiload_def}). Its value is:
\[
    \Phi_\mathrm{load} = T_\mathrm{vac} \cdot \Delta = 55,
\]
proved from the Lucas-Cassini identity without approximation. The fractional anomaly
ERP is then $\delta_A = (\delta_\mathrm{slip}, \Phi_\mathrm{load}) = (2, 55)$: the
irreducible orbital surplus $\delta_\mathrm{slip}$ relative to the total geometric
phase load of the vacuum.

The Dual Identity (Theorem~2) confirms that Path~A (algebraic) and Path~B
(geometric) are structurally independent derivations of the same ERP, connected by
the Triadic-Saturation Gap identity $N^3 - \Delta^2 = \delta_\mathrm{slip} = G_\mathrm{mg} = 2$.

Three further independent paths are supplied: Path~C (mass gap over flux-vacuum
product $G_\mathrm{mg}/(\Delta \cdot T_\mathrm{vac})$); Path~D (continued fraction
$[0; N^3, \delta_\mathrm{slip}]$); and Path~E (phase slip over Fibonacci inertial
threshold $\delta_\mathrm{slip}/F(\tau_{10})$). Together with Path~A (established in
Theorem~13.11) and Path~B (proved here), five structurally independent derivations of
$\delta_A = (2, 55)$ are available.

\end{rsbox}

\begin{observebox}
\textbf{Causal reading of the fractional anomaly $2/55$.}

The fractional part of $\alpha_{RS}^{-1} = 137 + 2/55$ encodes the ratio of
the phase slip $\delta_{\mathrm{slip}} = 2$ to the $\koppa$-ledger accumulation
of one complete vacuum cycle from the vacuum seed under $\psi$:
$\koppa_{\mathrm{ledger}}(T_{\mathrm{vac}}\text{ steps}) = 55
= \Delta \cdot T_{\mathrm{vac}}$.

The fine structure constant fractional anomaly is the ratio of temporal phase
drift to vacuum kinematic depth.  The denominator $55$ is the scale at which
the solar neutrino mixing channel operates; both facts reflect the same
underlying vacuum koppa structure.

The proton-electron mass ratio carries a fractional anomaly of $8/55$:
\[
\frac{8}{55} = N_E \cdot \frac{2}{55} = \psi\text{-period} \cdot \delta_A.
\]
The proton fractional correction encodes $N_E = 4$ full $\psi$-cycles of
phase slip relative to the single-quantum (electron) scale.  The denominator
$55$ is shared by both fractional anomalies; the numerator ratio
$8/2 = N_E$ is the $\psi$-period.
\end{observebox}
\newpage


\section{Glossary of Terms}\label{sec:glossary}

All terms are defined within the RigbySpace formal system.  No term carries
automatic empirical correspondence to external physics unless an explicit
mapping is derived.

\bigskip

\noindent\textbf{Aggregate Transform ($\eta$, $\aop$).}
$\aop(n,d)=(n+d,n)$.  Drives convergence toward the golden ratio attractor
$\phi$ as barycentric limit.  From the vacuum seed $(T_{\mathrm{vac}},\Sigimb)=(11,7)$,
generates the Lucas sequence.  Corresponds to the continued fraction
operation $[1;\ldots]$.

\medskip\noindent\textbf{Arrow of Time.}
The $+1$ net temporal surplus per 11-microtick cycle.  Forced by the
incommensurability $T_{\mathrm{vac}}=11$, $N=3$, derived from the 4-4-3
partition (Theorem~\ref{thm:arrow-443}).  Irreversibly crystallised at the
inertial closure tick $\tau_{10}$ (Theorem~\ref{thm:tau10}).

\medskip\noindent\textbf{Barycentric Addition ($\boxplus$).}
Engine of rank growth and mass generation.  Linear form:
$(n_1,d_1)\bplus(n_2,d_2)=(n_1 d_2+n_2 d_1,d_1 d_2)$.

\medskip\noindent\textbf{Barycentric Period ($T_{\mathrm{bary}}=33$).}
Full period $T_{\mathrm{vac}}\times N=11\times 3=33$ required for phase and
microtick counters to resynchronise.

\medskip\noindent\textbf{Barycentric Structure.}
Stable limit object of a triadic system: a cyclic set of ERPs
$\{L_E,L_M,L_R\}$ defining the oscillation's effective centre.  Not a
real-valued scalar limit; defined by entry into a discrete bracket with $W=1$.

\medskip\noindent\textbf{Boundary-Active Phase Sum ($\Sigimb=7$).}
$\Sigimb:=N_E+N_R=4+3=7$; equivalently $\Delta+\delta_{\mathrm{slip}}=7$.
The sum of the two phases that generate boundary stress and temporal charge.
The Modulation phase is excluded.

\medskip\noindent\textbf{Boundary Stress.}
Geometric asymmetry at cycle closure from $N_R=3<N_E=N_M=4$.  Accumulation
site for $\koppa_{\mathrm{ledger}}$.  Origin of restorative tension.

\medskip\noindent\textbf{Bridge Constant ($\mathcal{B}=1303/22$).}
Ratio $\delta_\alpha/\delta_Q$.  RS-constructive: $22=\Tbary-T_{\mathrm{vac}}$;
$1303$ is prime; $1303\equiv 5\pmod{11}$ and $1303\equiv 5\pmod{22}$ carry
the flux imbalance signature.

\medskip\noindent\textbf{Causal Step.}
One application of a vacuum or matter operator; the fundamental unit of
causal time $t$.

\medskip\noindent\textbf{Charge Conjugation ($\mathcal{C}$).}
The ontic origin of charge conjugation is $\psi^2$: the full ERP-pair swap
that inverts the cross-determinant sign, mapping matter chirality to
antimatter chirality (Theorem~\ref{thm:cc}).

\medskip\noindent\textbf{Chiral Stability.}
Condition $\Dcross=+1$ selecting matter states; $-1$ for antimatter; $0$ for
collapsed identity.  Derives matter-antimatter distinction from integer
ordering without an external norm.


\medskip\noindent\textbf{Constraint Vacuum.}
State $n/0$, $n\neq 0$.  Suspension state; structural history $n$ preserved
in numerator; metric basis absent.  Resolved by $\psi$ paired with an active
partner; history redistributed to the kinematic and structural entropy
ledgers at $\Omega$ (Definition~\ref{def:zero-resolution}).

\medskip\noindent\textbf{Constructive Remainder.}
$x\koppa L$: iterative subtraction; no division; replaces modular arithmetic
as ontological primitive.

\medskip\noindent\textbf{Cross-Determinant Deviation ($\Dcross$).}
$\Dcross(a,b)=p_b q_a-p_a q_b$.  Replaces the absolute value metric.
Encodes ordering (sign) and deviation (magnitude).

\medskip\noindent\textbf{Cross-Ordering ($\prec$).}
$a\prec b$ iff $\Dcross(a,b)>0$; discrete total order on $\SL/{\sim}$.

\medskip\noindent\textbf{Cross-Synchronisation Lift.}
Embedding $L(\upsilon,\beta)=(n_\upsilon d_\beta,n_\beta d_\upsilon,
d_\upsilon d_\beta)$ carrying two linear ERPs into $\SPs$.

\medskip\noindent\textbf{Cycle Boundary Event ($\Omega$).}
The reset $\succ(11)=1$; the locus where accumulated flux imbalances are
reconciled and $\koppa_{\mathrm{ledger}}$ is updated.

\medskip\noindent\textbf{Cycle Discriminant.}
$C:=\mathrm{tr}(M)^2-4\det(M)$ for $M\in GL_2(\Z)$; identified with
inertial mass for minimal orbits under the Discriminant Mass Law
(Theorem~\ref{thm:disc-mass}).  Equals geometric capacity $d_{\mathrm{stable}}$
at stabilisation.

\medskip\noindent\textbf{Democratic Axis.}
The equal-weight direction $(1,1,1)/\sqrt{3}$ in lepton mass space; the
Koide condition places the mass vector at $45^\circ$ to this axis,
corresponding to $Q=2/3$.  A continuous-lens concept; see External
Verification in Theorem~\ref{thm:koide}.

\medskip\noindent\textbf{Discrete Bracket.}
Ordered pair $(L,U)$ with $L\prec U$; bounded interval in $\SL$ with integer
width $W(L,U)$.

\medskip\noindent\textbf{Discrete Convergence.}
A sequence $\{x_n\}\subset\SL$ converges if there exist nested brackets
$[L_n,U_n]$ with $x_n\in[L_n,U_n]$ and entry into a $W=1$ bracket in finite
steps.  No real-valued limit invoked.

\medskip\noindent\textbf{Discrete L-Function ($\LQ(N)$).}
$(R_N,T_N)\in\SL$; orientation-bias summary of an elliptic orbit; integer
analogue of the Hasse-Weil L-function.

\medskip\noindent\textbf{Discriminant Mass Law.}
$m_{\mathrm{obs}}=C=d_{\mathrm{stable}}$ for minimal-orbit generation cycles.
Mass is the orbit discriminant, not a background parameter.
(Theorem~\ref{thm:disc-mass}.)

\medskip\noindent\textbf{Doubling Generator (Case B).}
Projective matter generator with $P_1=P_2$; cubic rank growth; drives the
Arrow of Time.  Correct fully-projective formulas:
$W_B:=3X_1^2+aZ_1^2$, $S_B:=Y_1Z_1$, $B_B:=X_1Y_1S_B$,
$H_B:=W_B^2-8B_B$; $X_3:=2H_BS_B$, $Y_3:=W_B(4B_B-H_B)-8Y_1^2S_B^2$,
$Z_3:=8S_B^3$.

\medskip\noindent\textbf{Dyadic Closure Class 64 ($\mathcal{R}_D$).}
Resolution class of capacity $C=64=2^6$; phase defective ($64\equiv 1\pmod 3$);
quark sector.

\medskip\noindent\textbf{Emission Phase ($E$).}
Active output phase; $N_E=4$ ticks $\{1,4,7,10\}$; temporal contribution
$+1$ per tick.

\medskip\noindent\textbf{Explicit Rational Pair (ERP).}
Ordered pair $(n,d)\in\Z\times(\Z\setminus\{0\})$; unreduced state; the
fundamental object of the framework.

\medskip\noindent\textbf{Extended Dyadic Class 128 ($\mathcal{R}_E$).}
Resolution class $C=128=2^7$; phase defective ($128\equiv 2\pmod 3$);
identified with the Tau lepton.

\medskip\noindent\textbf{4-4-3 Partition.}
Derived partition of the 11-microtick cycle: $E$ ($N_E=4$), $M$ ($N_M=4$),
$R$ ($N_R=3$); not postulated.

\medskip\noindent\textbf{Fibonacci Prime.}
An integer that is both prime and a member of the denominator sequence of
the $\aop$-orbit from the unit active seed $(1,1)$ (the Fibonacci sequence).
The first six are $\{2,3,5,13,89,233\}$.
The first three equal $\{G_{\mathrm{mg}},N,\Delta\}$ exactly
(Theorem~\ref{thm:fib-primes}).

\medskip\noindent\textbf{Flux Imbalance ($\Delta=5$).}
$(N_E+N_M)-N_R=5$; the fundamental zero-point potential of the vacuum engine;
derived from graph topology.

\medskip\noindent\textbf{Fractional Anomaly ($\delta_A=2/55$).}
Correction to integer fine-structure constant; $\delta_A=\delta_{\mathrm{slip}}
/(\Delta\cdot T_{\mathrm{vac}})=(N^3-K_{\mathrm{sat}})/55$; derived via two
independent paths (Theorem~\ref{thm:dual-anomaly}).

\medskip\noindent\textbf{Frictionless Propagation Principle.}
A prime-capacity ERP cannot shed structural load into sub-period oscillators;
it must engage the full vacuum architecture or not at all.  The structural
mechanism of mass precipitation at prime-capacity states.

\medskip\noindent\textbf{Gauge Manifold Dimension ($\Dg=12$).}
$1+3+8=12$; also $\Delta+\Sigimb=5+7$ (Goldbach path, Theorem~\ref{thm:goldbach-dg});
also $L(2)\times L(3)=N\times N_E$ (Lucas path, Theorem~\ref{thm:lucas-encoding}).

\medskip\noindent\textbf{Geometric Capacity.}
Denominator $d$ of an active oscillator; determines local vacuum resolution
density.

\medskip\noindent\textbf{Height Function ($h(m)$).}
Unique $k\geq 0$ with $2^{2k}\leq m^2<2^{2k+2}$; sign-independent surrogate
for the magnitude of $m$.

\medskip\noindent\textbf{Imbalance Ledger ($\koppa_{\mathrm{ledger}}$).}
Accumulated structural residue at the cycle boundary; integer whose structural
rank equals rest mass.

\medskip\noindent\textbf{Imbalance Operator ($\koppa$).}
$x\koppa L$: constructive remainder; iterative subtraction only; no division.

\medskip\noindent\textbf{Incommensurability, Triadic.}
$11\equiv 2\pmod 3$: after 11 successor operations on a 3-state cycle, net
displacement is $+2$ states.  The engine of all non-equilibrium evolution.

\medskip\noindent\textbf{Inertial Closure Tick ($\tau_{10}=10$).}
$\tau_{10}:=T_{\mathrm{vac}}-1=10$; the unique microtick at which the
temporal Arrow becomes irreversible within the cycle.  Structural identity:
$\tau_{10}+\delta_{\mathrm{slip}}=\Dg$ (Theorem~\ref{thm:tau10}).

\medskip\noindent\textbf{Inertial Mass.}
$m_I\equiv\rank(\koppa_{\mathrm{ledger}})$.  Equal to rest mass by
single-origin definition, not by measurement.

\medskip\noindent\textbf{Integer.}
The sole ontological primitive; element of $\Z$.  The integer is substance,
not data.

\medskip\noindent\textbf{Interaction Surface ($\Sint=132$).}
$T_{\mathrm{vac}}\times\Dg=11\times 12=132$; minimal count of states required
to synchronise vacuum cycle with gauge manifold.  Two independent
derivations: additive sum and barycentric $\bplus$ path.

\medskip\noindent\textbf{Interval Width ($W$).}
$W(L,U):=\Dcross(L,U)\in\Z_{>0}$ for $L\prec U$; discrete distance;
saturation at $W=1$.

\medskip\noindent\textbf{Koppa ($\koppa$).}
See \textit{Imbalance Operator} and \textit{Imbalance Ledger}.  Symbol:
archaic Greek koppa.

\medskip\noindent\textbf{Koide Correspondence ($Q\approx\varepsilon=2/3$).}
The phase slip efficiency $\varepsilon=2/3$ is the RS-structural result.  The
empirical Koide ratio $Q$ of the lepton masses equals $\varepsilon$; this is
an external correspondence recorded in the continuous lens.

\medskip\noindent\textbf{Lattice Coupling Constant ($\Lalpha$).}
$(\mathrm{count}(\Omega_\psi),\mathrm{count}(\Omega_{33}))\in\SL$; discrete
fine-structure analogue; asymptotically stabilises to $\alpha^{-1}\approx 137$.

\medskip\noindent\textbf{Linear Transform ($\lambda$, $\lop$).}
$\lop(n,d)=(n+d,d)$; vacuum propagation; increments magnitude without
altering potential; matrix $\mathbf{L}$.

\medskip\noindent\textbf{Lucas Number ($L(n)$).}
$L(0)=2$, $L(1)=1$, $L(n)=L(n-1)+L(n-2)$.  Generated by $\aop$ from the
vacuum seed $(11,7)$.  The foundational RS constants $\{2,3,4,7,11\}$ are
Lucas numbers at positions $\{0,2,3,4,5\}$ (Theorem~\ref{thm:lucas-encoding}).

\medskip\noindent\textbf{Luminal State.}
$\Tcharge=0$ ($n=d$); no proper time accumulates.

\medskip\noindent\textbf{Manin-Drinfeld Relations.}
$\{A,B\}+\{B,C\}=\{A,C\}$; $\{A,A\}=0$; $\{A,B\}+\{B,A\}=0$.

\medskip\noindent\textbf{Mask Postulate.}
Continuous constants are phenomenological masks of high-frequency rational
oscillations; not ontological primitives.  All appear exclusively in External
Verification boxes.

\medskip\noindent\textbf{Mass Gap ($G_{\mathrm{mg}}$).}
$G_{\mathrm{mg}} := T_{\mathrm{vac}}\koppa N = 11\koppa 3 = 2$.
Derived, not primitive: it is the per-cycle phase slip, the constructive
remainder when the prime vacuum period $T_{\mathrm{vac}}$ is measured
against the triadic period $N$.  Because $T_{\mathrm{vac}}$ is prime and
does not divide by $N$, this residual is irresolvable and permanent.
It is the deposit made by $\Omega$ into the temporal ledger at each cycle
boundary, and the base unit of the RS structural rank.  See
Definition~\ref{def:Gmg} and Definition~\ref{def:Omega}.

\medskip\noindent\textbf{Mass-Gap Postulate.}
Two-phase axiom: vacuum (additive rank growth) and matter (multiplicative
rank growth).

\medskip\noindent\textbf{Matter Generator.}
Projective polynomial operator $G:\SP\to\SP$; Cases A (identity), B
(doubling - corrected CMO formulas), C (addition); no division; all outputs
in $\Z$.

\medskip\noindent\textbf{Mediant Geodesic ($\gamma(A,B)$).}
Unique path from $A$ to $B$ in the Mediant Tree; length equals continued
fraction partial quotient sum.

\medskip\noindent\textbf{Mediant Sum ($\oplus$, $\med$).}
$(a,b)\med(c,d):=(a+c,b+d)$; preserves cross-determinant width.

\medskip\noindent\textbf{Mediant Tree ($\mathcal{T}$).}
Binary tree of integer pairs generated by iterative mediant insertion from
roots $(0,1)$ and $(1,0)$.

\medskip\noindent\textbf{Memory Phase ($M$).}
Near-luminal phase; $N_M=4$ ticks; contributes $0$ to temporal surplus.
Carries structural history between cycle boundaries.

\medskip\noindent\textbf{Microtick.}
Unit of discrete temporal evolution; element of $\Tcyc=\{1,\dots,11\}$.

\medskip\noindent\textbf{Minimal Uniform Structural Tension.}
State of capacity $4621=2\times\#(T_{\mathrm{vac}})+1$: carries exactly
$+1$ unresolved unit of structural history against every fundamental RS cycle
simultaneously (Theorem~\ref{thm:primorial}).

\medskip\noindent\textbf{Null State / Empty State.}
The label $0/0$ denotes an empty lattice position in a mediant propagation
context.  No structural content; not a persistent algebraic object.  The
same emptiness as $0/1$, encountered from a different propagation context.
See Definition~\ref{def:null}.

\medskip\noindent\textbf{Numerical ZERO / Zero Magnitude.}
State $0/d$, $d>0$: collapses to empty state.  Not a valid physical state.
Zero magnitude means the lattice position is empty.
See Definition~\ref{def:num-zero}.

\medskip\noindent\textbf{Observational Equivalence ($\sim$).}
$(n_1,d_1)\sim(n_2,d_2)$ iff $n_1 d_2=n_2 d_1$; sole admissible equivalence.

\medskip\noindent\textbf{Operator Algebra.}
The collective action of $\{\lop,\aop,\psi,\bplus,\succ,\koppa,\PG\}$
on the state spaces.  The engine of the framework.

\medskip\noindent\textbf{Orientation Component ($\Theta(s)$).}
$\Theta(s):=\tau\in\Tcyc$; modulates interaction eligibility.

\medskip\noindent\textbf{Pell Equation.}
$a^2-nb^2=1$; source of cycle matrix coefficients $(a,b)$ for attractor
$\sqrt{n}$; the integer foundation of all quadratic irrational convergence.

\medskip\noindent\textbf{Phase Evolution ($\phi(t)$).}
$\phi(0):=1$; $\phi(t+1):=\succ(\phi(t))$; integer-valued; no division.

\medskip\noindent\textbf{Phase Precession.}
$+2$ states per 11-tick cycle; driven by triadic incommensurability;
prevents closure into a static loop.

\medskip\noindent\textbf{Phase Slip ($\delta_{\mathrm{slip}}=2$).}
After 11 successor operations on a 3-state cycle, net displacement is
$+2$ states.

\medskip\noindent\textbf{Phase Slip Efficiency ($\varepsilon=2/3$).}
$\varepsilon:=\delta_{\mathrm{slip}}/N=2/3$; the RS structural result;
empirically equals the Koide ratio $Q$ (Theorem~\ref{thm:koide}).

\medskip\noindent\textbf{Pi-Analog ERP ($\piRS=(355,113)$).}
The RS-native rational pair reached by four legs of the mediant tree path
consuming the partial quotients $N$, $\Sigimb$, $N\times\Delta$, and the
vacuum unit in sequence, each consecutive convergent pair adjacent with $W=1$.
Total koppa load $\Gmg\times\Pithree=26$.  Both components carry $N=3$
against $\Tvac$.  The $\pi$-remainder after barycentric truncation is
structurally identified as the source of the Koide correction to
$\varepsilon=2/3$ (Definition~\ref{def:pi-analog},
Theorem~\ref{thm:pi-koppa}).

\medskip\noindent\textbf{Phase Space ($\Phi_N$).}
Finite set $\{0,\dots,N-1\}$ with successor map $\succ_N$; discrete
replacement for modular arithmetic in phase-tracking.

\medskip\noindent\textbf{Phenomenological Mask.}
A continuous constant that appears as the barycentric average of a
high-frequency integer orbit; not ontologically primitive.

\medskip\noindent\textbf{Phi-Oscillator.}
The ERP $(7477,4621)$: the unique minimal prime-prime approximant to $\phi$
accessible from the Lucas seed $(11,7)$ via the Mediant Tree.  Its CF
$[1^{12},5,1,4]$ encodes $\Dg$, $\Delta$, $N_E$.  Its descent path generates
the full lepton mass hierarchy as intermediate products.

\medskip\noindent\textbf{Prime Gate ($\PG(k)$).}
The $k$-th prime in natural order; an admissible RS primitive, detectable
exclusively through $\koppa$, comparison, and $\succ$
(Theorem~\ref{thm:pg-admissible}).

\medskip\noindent\textbf{Prime-11 Time Cycle ($\Tcyc$).}
$\{1,\dots,11\}$ with successor $\succ$; primality of 11 ensures
11 does not divide $\Dg=12$, so $\Sint=T_{\mathrm{vac}}\times\Dg=132$.

\medskip\noindent\textbf{Primorial ($\#(p)$).}
Product of all primes $\leq p$; $\#(T_{\mathrm{vac}})=2310$.

\medskip\noindent\textbf{Prime Pair, RS Sommerfeld ($(\pa,\qa)=(115331,15804499)$).}
The ordered pair of primes generated by the CF convergent recurrence from
$\Da=[\aSC;\,N^3,1,N,\dslip,N\times\Sigimb,1,1,\tauX]$.  The Euclidean
descent of $(\qa,\pa)$ recovers $\Da$ exactly and terminates at $\tauX$.
Every quotient and remainder of the descent is expressible as a product of
RS structural constants (Theorem~\ref{thm:descent-closed}).  Both primes
carry $\Sigimb=7$ against $\Tvac$.  On the Verification Path,
$\pa/\qa$ lies within $8.26\times10^{-11}$ of the CODATA fine structure
constant.

\medskip\noindent\textbf{Projective Identity ($\IP$).}
$(0,1,0)$; identity element of the elliptic group law; Case A of the matter
generator.

\medskip\noindent\textbf{Projective Lifting.}
See \textit{Cross-Synchronisation Lift}.

\medskip\noindent\textbf{PMNS Bare Mixing Angles.}
$\sin^2\theta_{13}^{\mathrm{bare}}=1/45$,
$\sin^2\theta_{12}^{\mathrm{bare}}=3/10$,
$\sin^2\theta_{23}^{\mathrm{bare}}=7/12$.
Derived from RS structural integers; bare values systematically exceed
experiment (Theorem~\ref{thm:pmns}).

\medskip\noindent\textbf{Quark Mass Integers.}
$\{\Lambda_{q1},\Lambda_{q2},\Lambda_{q3}\}=\{207,303,192\}$; derived via the
three-step path $\Sint+C+T_{\mathrm{vac}}$ for generation capacities
$\{64,160,192\}$ (Theorem~\ref{thm:quark-mass}).

\medskip\noindent\textbf{Rank Function ($\rho$).}
$\rank(m)$: unique $k\geq 0$ such that $G_{\mathrm{mg}}^k\leq m 
G_{\mathrm{mg}}^{k+1}$, where $G_{\mathrm{mg}} = T_{\mathrm{vac}}\koppa N = 2$
is the mass gap (Definition~\ref{def:Gmg}).  Computed RS-natively by
initialising $\mathit{power}:=1$ and repeatedly multiplying by
$G_{\mathrm{mg}}$ while $\mathit{power}\times G_{\mathrm{mg}}\leq m$,
counting increments; no floor function, logarithm, or division appears.
$\rank(0):=0$.  Measures causal depth in units of the mass gap: the number
of wave/particle crossings required to reach structural depth $m$.
See also \textit{Prime-Event Depth ($\rho_\Omega$)}.

\medskip\noindent\textbf{Prime-Event Depth ($\rho_\Omega$).}
$\rho_\Omega(m)$: total number of prime factors of $m$ counted with
multiplicity (the arithmetic function $\Omega(m)$).  Computed RS-natively
by repeated $\koppa$-primality testing and iterative subtraction of each
detected prime factor; no division operator appears.  $\rho_\Omega(0) :=
\rho_\Omega(1) := 0$.  Measures the number of prime-gated $\psi$-events
required to fully decompose state $m$ into irreducible structural
components.  Distinguished from the Rank Function $\rank(m)$: rank measures
binary causal depth; prime-event depth measures koppa-decomposition depth.
Key values: $\rho_\Omega(\mathcal{V}_S\!=\!18)=N=3$ (generation count);
$\rho_\Omega(\Lambda_\mu\!=\!207)=N=3$; $\rho_\Omega(\Lambda_\tau\!=\!3481)=2$.

\medskip\noindent\textbf{Rational Lorentz Boost.}
Boost transformation $A_Q$ on events; interval-preserving; zero numerical
drift; proved via integer identity.

\medskip\noindent\textbf{Resolution Class ($\mathcal{R}_C$).}
Set of history trajectories stabilising at geometric capacity $d=C$; defines
a particle species.

\medskip\noindent\textbf{Rest Mass.}
$m_{\mathrm{rest}}\equiv\rank(\koppa_{\mathrm{ledger}})$; accumulated
historical imbalance measured by structural depth of the imbalance ledger.

\medskip\noindent\textbf{Return Phase ($R$).}
Cycle closure phase; $N_R=3$ ticks $\{3,6,9\}$; each $R$-tick contributes
$-1$ to temporal surplus; contraction ($3<4$) generates Boundary Stress.
Ticks $\{3,6,9\}$ are the arrowless instants (zero net surplus).

\medskip\noindent\textbf{Rigby-Tate Correspondence.}
Isomorphism: temporal unfolding in RS takes $t = m^2$ causal steps,
equal to the quadratic growth of unreduced height $\hat{h}(mP) = m^2 \cdot \hat{h}(P)$.
Structural friction $\Phi(t)$ accumulates over $t$ steps.

\medskip\noindent\textbf{Saturation.}
$W(L,U)=1$; minimum discrete bracket width; the adjacency condition;
physical realisation criterion.

\medskip\noindent\textbf{Saturation Limit ($K_{\mathrm{sat}}=25$).}
$\Delta^2=25$; the vacuum's finite absorption capacity.

\medskip\noindent\textbf{Secondary Stability Integer ($\Lambda_\mu=207$).}
$\Sint+\mathcal{R}_D+T_{\mathrm{vac}}=132+64+11=207$; governs muon mass
hierarchy; shared with Gen-1 quark mass integer.

\medskip\noindent\textbf{Sign Function ($\sgn$).}
$\sgn(m)\in\{-1,0,+1\}$; integer-only; used in orientation-bias sums.

\medskip\noindent\textbf{Sommerfeld Constant ($\aSC=137$).}
$\Sint+\Delta=132+5=137$; the 33rd prime $=\PG(\Tbary)$; integer baseline
for $\alpha^{-1}$; electromagnetic coupling stability is a consequence of
its primality.  The leading partial quotient of the CF descent of the RS
Sommerfeld ERP (Theorem~\ref{thm:sommerfeld}).

\medskip\noindent\textbf{Spacetime Event ($E$).}
$((n_t,d_t),(n_x,d_x))\in\SL\times\SL$; temporal and spatial ERP pair.

\medskip\noindent\textbf{Structural Entropy ($H(s)$).}
$\rank(d)$; causal history of the state; cost of the state's metric basis.

\medskip\noindent\textbf{Structural Friction ($\Phi(t)$).}
Cumulative rank accumulation over $t$ steps.
Discrete manifestation of resistance to orbit growth.

\medskip\noindent\textbf{Structural Rank.}
See \textit{Rank Function}.

\medskip\noindent\textbf{Structural Symbol ($\{A,B\}$).}
Formal sum of mediant geodesic edges; satisfies Manin-Drinfeld relations.

\medskip\noindent\textbf{Successor with Reset ($\succ$).}
$\succ(\tau)=\tau+1$ ($\tau<11$); $\succ(11)=1$; total and deterministic; no
arithmetic reduction.

\medskip\noindent\textbf{Superluminal State.}
$\Tcharge<0$ ($n>d$); structurally excluded by Chiral Stability.

\medskip\noindent\textbf{Suspended State.}
State containing a Constraint Vacuum component $n/0$; evolution frozen
until $\psi$ resolves it by acting paired with an active partner,
redistributing structural history to the kinematic and entropy ledgers
(Definition~\ref{def:zero-resolution}).

\medskip\noindent\textbf{Temporal Charge ($\Gamma$).}
$\Tcharge(n,d):=d^2-n^2$; cross-determinant of a state with its own swap;
determines luminal regime.

\medskip\noindent\textbf{Threshold Semiprime ($1309$).}
$\Sigimb\times T_{\mathrm{vac}}\times(\Dg+\Delta)=7\times 11\times 17=1309$;
the convergent of $(7477,4621)$ at depth $\Dg$; the lattice must transit this
state before a prime-prime $\phi$-approximant can emerge.

\medskip\noindent\textbf{Topological State Vector ($\Psistate$).}
$(M,H,\Theta)$: projection of an ERP onto the physical manifold; Magnitude,
History, Orientation.

\medskip\noindent\textbf{Transformative Reciprocal ($\psi$, $\psi$).}
$\psi\!\left(\tfrac{a}{b},\tfrac{c}{d}\right):=\left(\tfrac{d}{a},\tfrac{b}{c}\right)$;
denominators cross to become numerators of the other state; numerators become
new denominators; coupled they cannot act on one state in isolation; order~4;
$\psi^4=\mathrm{Id}$; $\psi^2$ is charge conjugation; mechanism for spin,
chirality, Constraint Vacuum resolution, and charge conjugation.
Period~4 equals the Emission cardinality $N_E$
(Proposition~\ref{prop:psi-ne}).

\medskip\noindent\textbf{Triadic Closure Class 72 ($\mathcal{R}_T$).}
Resolution class $C=72=8\times 9$; $72\equiv 0\pmod 3$ (phase-locked);
identified with the Electron; the shell $72-64=8$ accounts for its
exceptional stability.

\medskip\noindent\textbf{Triadic Convergence.}
The process by which a system settles into a stable 3-phase cycle $\{E,M,R\}$
rather than a single scalar limit.

\medskip\noindent\textbf{Triadic Incommensurability.}
See \textit{Incommensurability, Triadic}.

\medskip\noindent\textbf{Unreduced Height ($\hat{h}(P)$).}
$\max(\rank(n),\rank(d))$ for point $P$ with coordinate $n/d$; grows as
$\Theta(m^2)$ for $mP$.

\medskip\noindent\textbf{Vacuum.}
The Rank-0/2 phase of evolution; governed by $\{\lop,\aop,\psi\}$; rank
growth additive.

\medskip\noindent\textbf{Vacuum Line ($\Vac$).}
$\{(p,1):p\in\Z\}\subset\SL$; unit-denominator ERPs.

\medskip\noindent\textbf{Vacuum Resolution Frequency ($\Omega_{\mathrm{vac}}$).}
Integer count of vacuum generator applications required to balance a tension;
discrete surrogate for energy.

\medskip\noindent\textbf{Vacuum Seed ($S_{\mathrm{vac}}=(1,11)$).}
Unitary state governed by the vacuum period $T=11$.

\medskip\noindent\textbf{Viscosity Sum ($\VS$).}
Accumulated phase drift of a system over a closed cycle; $\VS=18$ for the
standard vacuum orbit; invariant under cyclic permutation of $\tau_0$.

\medskip\noindent\textbf{Width ($W$).}
See \textit{Interval Width}.

\newpage

\section{Notation Quick Reference (Single-Page Summary)}\label{sec:quickref}

\begin{longtable}{L{2.2cm} L{5.5cm} L{5.5cm}}
\toprule
\textbf{Symbol} & \textbf{Name} & \textbf{One-Line Definition} \\
\midrule\endhead\bottomrule\endfoot

$\Z$ & Integers & Sole ontological substrate \\
$(n,d)$ & ERP & Ordered pair in $\Z\times(\Z\setminus\{0\})$; never reduced \\
$\SL$ & Linear state space & $\{(n,d)\in\Z^2:d\neq 0\}$ \\
$\SPs$ & Active matter space & $\{(X,Y,Z)\in\Z^3:XYZ\neq 0\}$ \\
$\IP$ & Projective identity & $(0,1,0)$ \\
$\lop$ & Linear transform & $(n,d)\mapsto(n+d,d)$ \\
$\aop$ & Aggregate transform & $(n,d)\mapsto(n+d,n)$; Lucas generator from $(11,7)$ \\
$\psi$ & Transformative reciprocal &
  $\left(\tfrac{a}{b},\tfrac{c}{d}\right)\mapsto\left(\tfrac{d}{a},\tfrac{b}{c}\right)$;
  coupled; order~4; $\psi^2$ = charge conjugation \\
$\bplus$ & Barycentric addition &
  $(n_1,d_1)\bplus(n_2,d_2)=(n_1 d_2+n_2 d_1,d_1 d_2)$ \\
$\med$ & Mediant sum & $(a,b)\med(c,d)=(a+c,b+d)$ \\
$\succ$ & Successor with reset & $\succ(\tau)=\tau+1$; $\succ(11)=1$ \\
$\koppa$ & Imbalance operator/ledger &
  Constructive remainder; $m_{\mathrm{rest}}=\rank(\koppa_{\mathrm{ledger}})$ \\
$\PG(k)$ & Prime Gate & $k$-th prime; admissible RS primitive \\
$\rank(m)$ & Structural rank &
  Unique $k\geq 0$: $G_{\mathrm{mg}}^k\leq m < G_{\mathrm{mg}}^{k+1}$;
  computed by counting $G_{\mathrm{mg}}$-doublings; $\rank(0):=0$ \\
$\hfn(m)$ & Height function &
  Unique $k$: $2^{2k}\leq m^2<2^{2k+2}$; sign-independent \\
$\Dcross(a,b)$ & Cross-determinant & $p_b q_a-p_a q_b$; replaces metric distance \\
$W(L,U)$ & Interval width & $\Dcross(L,U)$; saturation at $W=1$ \\
$\Tcharge$ & Temporal charge & $d^2-n^2$; $>0$ massive, $=0$ luminal, $<0$ excluded \\
$\VS$ & Viscosity sum & Accumulated phase drift; $\VS=18$ for standard orbit \\
$\Tcyc$ & Prime-11 cycle & $\{1,\dots,11\}$ with $\succ$; 11 prime, does not divide 12 \\
$\Tbary$ & Barycentric period & $11\times 3=33$ \\
$\tau_{10}$ & Inertial closure tick &
  $T_{\mathrm{vac}}-1=10$; Arrow irreversibly fixed; $\tau_{10}+\delta_{\mathrm{slip}}=\Dg$ \\
$\delta_{\mathrm{slip}}$ & Phase slip & $+2$ states after 11 ticks on 3-cycle \\
$\varepsilon=Q$ & Phase slip efficiency / Koide & $2/3$ \\
$\Sigimb$ & Boundary-active phase sum & $N_E+N_R=4+3=7$; also $\Delta+\delta_{\mathrm{slip}}$ \\
$\Sint$ & Interaction surface & $T_{\mathrm{vac}}\times\Dg=11\times 12=132$ \\
$\Dg$ & Gauge dimension &
  $1+3+8=\Delta+\Sigimb=L(2)\times L(3)=12$ \\
$\alpha^{-1}_{RS}$ & Inverse fine-structure & $137+2/55$ \\
$\mu$ & Proton-electron mass ratio & $1836+8/55$ \\
$\Lambda_\mu$ & Muon mass integer & $132+64+11=207$ \\
$\Lambda_\tau$ & Tau mass integer & $\PG(\Dg+\Delta)^2=59^2=3481$ \\
$\Bridge$ & Generation bridge & $1303/22$ \\
$\mathcal{R}_D$ & Dyadic class & $C=64$; quark sector \\
$\mathcal{R}_T$ & Triadic class & $C=72$; Electron \\
$\mathcal{R}_E$ & Extended dyadic class & $C=128$; Tau \\
$L(n)$ & Lucas number &
  $L(0)=2$, $L(1)=1$, $L(n)=L(n-1)+L(n-2)$; generated by $\aop$ from $(11,7)$ \\
$\#(p)$ & Primorial &
  Product of primes $\leq p$; $\#(T_{\mathrm{vac}})=2310$ \\
$\mathbf{L},\mathbf{H},\mathbf{S}$ & Vacuum matrices &
  $\bigl(\begin{smallmatrix}1&1\\0&1\end{smallmatrix}\bigr)$,
  $\bigl(\begin{smallmatrix}1&1\\1&0\end{smallmatrix}\bigr)$,
  $\bigl(\begin{smallmatrix}0&1\\1&0\end{smallmatrix}\bigr)$ \\
$\LQ(N)$ & Discrete L-function &
  $(R_N,T_N)\in\SL$; orbit orientation bias \\
$\hat{h}(mP)$ & Unreduced height & $=\Theta(m^2)$; Rigby-Tate growth \\
$\Phi(t)$ & Structural friction & cumulative rank accumulation \\
$\mathcal{T}$ & Mediant tree &
  Binary tree; roots $(0,1)$, $(1,0)$; mediant insertion \\
$\gamma(A,B)$ & Mediant geodesic &
  Unique path in $\mathcal{T}$; length $=$ CF partial quotient sum \\
$\{A,B\}$ & Structural symbol &
  Formal geodesic sum; Manin-Drinfeld relations \\
$L(\upsilon,\beta)$ & Cross-sync lift &
  $(n_\upsilon d_\beta,n_\beta d_\upsilon,d_\upsilon d_\beta)$;
  $\SL\times\SL\to\SPs$ \\
$\piRS$ & RS $\pi$-Analog ERP &
  $(355,113)$; four-leg mediant truncation;
  koppa load $\Gmg\times\Pithree=26$ \\
$(\pa,\qa)$ & RS Sommerfeld Prime Pair &
  $(115331,15804499)$; CF descent $\Da$ terminates at $\tauX$ \\
$\Da$ & CF Descent Sequence &
  $[\aSC;\,N^3,1,N,\dslip,N\times\Sigimb,1,1,\tauX]$ \\
\end{longtable}

\newpage


\appendix
\section{Alternative Derivation Paths}\label{app:alt-paths}

This appendix records all secondary and tertiary derivation paths for
canonical RS constants.  The primary derivation of each constant appears in
the main body.  Every path here is fully RS-axiomatic.

\subsection{Barycentric Derivation of $\Sint=132$}\label{app:sint-bary}

\noindent\textit{Primary path:} additive sum
$\Sint=\sum_{i=1}^{T_{\mathrm{vac}}}\Dg=132$
(Theorem~\ref{thm:sint-additive}).

\medskip
\noindent\textbf{Tertiary path - Barycentric $\bplus$ derivation.}

\begin{rsbox}
\[
(1,T_{\mathrm{vac}})\bplus(1,\Dg)
= (1,11)\bplus(1,12)
= (1\cdot 12+1\cdot 11,\;11\cdot 12)
= (23,\;132).
\]
The interaction surface $\Sint=132$ is the barycentric denominator produced
when the vacuum seed ERP $(1,T_{\mathrm{vac}})$ couples to the unit-gauge
state $(1,\Dg)$ via a single application of $\bplus$.  The numerator
$23=\PG(N^2)=\PG(9)$ is the Prime Gate of the triadic square.  This path
derives $\Sint$ through the fundamental matter-addition operator without
reference to $\mathrm{lcm}$ or an additive iteration.
\end{rsbox}

Structural reading: $(1,11)$ is the vacuum seed $S_{\mathrm{vac}}$;
$(1,12)=(1,\Dg)$ is the unit-gauge state.  Their barycentric product encodes
the minimal synchronisation requirement in the denominator and the triadic
square prime in the numerator simultaneously.

\medskip\noindent Adjoining derivation: $(1,G_{\mathrm{mg}})\bplus(1,T_{\mathrm{vac}})
=(1,2)\bplus(1,11)=(13,22)$.  The Bridge denominator $22=\Tbary-T_{\mathrm{vac}}$
is the barycentric product of the mass gap and the vacuum period; the
numerator $13=\PG(\Sigimb)$ is the 4th Fibonacci prime.

\subsection{Prime Gate Derivation of $\alpha^{-1}_{\mathrm{int}}=137$}
\label{app:alpha-primegate}

\noindent\textit{Primary path:}
$\alpha^{-1}_{\mathrm{int}}=\Sint+\Delta=132+5=137$
(Theorem~\ref{thm:alpha-int}).\\
\textit{Secondary path:} Uniqueness conditions of Theorem~\ref{thm:Dg-unique},
Condition~(5).

\medskip
\noindent\textbf{Tertiary path - Prime Gate at $T_{\mathrm{bary}}$.}

\begin{rsbox}
\[
\PG(T_{\mathrm{bary}}) = \PG(33) = 137 = \alpha^{-1}_{\mathrm{int}}.
\]
The Sommerfeld Constant is the prime at position $T_{\mathrm{bary}}=33$
in the natural ordering of primes.  This path accesses 137 via the
Frictionless Propagation structure: the vacuum barycentric period identifies
the prime-ordinal position, and $\PG$ retrieves the integer.  No reference
to $\Sint$ or $\Delta$ is made.
\end{rsbox}

The three paths of additive interaction surface, uniqueness conditions, and
Prime Gate at $T_{\mathrm{bary}}$ all converge to 137 from three structurally
independent starting points.  Flux signature: $137\equiv 5\pmod{11}$ and
$137\equiv 5\pmod{22}$ carry $\Delta=5$ in both residues.

\subsection{Phi-Oscillator Path to $\Lambda_\mu=207$}\label{app:lmu-phiosc}

\noindent\textit{Primary paths (TRTSv4):} four derivation paths for
$\Lambda_\mu=207$, including $\Sint+\mathcal{R}_D+T_{\mathrm{vac}}
=132+64+11=207$ and $N^2\times\PG(N^2)=9\times 23=207$.

\medskip
\noindent\textbf{Fifth path - Phi-oscillator descent.}

\begin{rsbox}
The penultimate state of the $\aop^{-1}$ descent from $(7477,4621)$ at
step~12 is $(29,\Delta)=(\PG(\tau_{10}),5)$.  From this state:
\[
\Lambda_\mu = \PG(\tau_{10})\times\Sigimb+N_E = 29\times 7+4 = 207.
\]
\end{rsbox}

Structural reading: $\PG(\tau_{10})=\PG(10)=29=\Tbary-N_E$ is the prime at
the inertial closure index; $\Sigimb=7$ is the boundary-active phase sum;
$N_E=4$ is the Emission phase count.  This path derives $\Lambda_\mu$ from
the phi-oscillator trajectory with no reference to $\Sint$, $\mathcal{R}_D$,
or $T_{\mathrm{vac}}$.

\subsection{Phi-Oscillator Path to $\Lambda_\tau=3481$}\label{app:ltau-phiosc}

\noindent\textit{Primary path:}
$\Lambda_\tau=\PG(\Dg+\Delta)^2=P_{17}^2=59^2=3481$
(Theorem~\ref{thm:pg-identities}, Eq.~\eqref{eq:pg-tau}).

\medskip
\noindent\textbf{Secondary path - Phi-oscillator descent.}

\begin{rsbox}
From the penultimate descent state
$(\PG(\tau_{10}),\Delta)=(29,5)$:
\[
\Lambda_\tau
= \PG(\tau_{10})\times(\Sint-T_{\mathrm{vac}}-1)+1
= 29\times 120+1
= 3481.
\]
where $\Sint-T_{\mathrm{vac}}-1=132-11-1=120=\delta_{\mathrm{slip}}^3
\times(T_{\mathrm{vac}}+N_E)=8\times 15$.
\end{rsbox}

Structural reading: the factor $120$ encodes the interaction surface
contracted by the vacuum period and a unit gap; the factor
$\PG(\tau_{10})=29$ is the inertial-closure prime; the unit offset $+1$
is the minimal structural tension.  This path is independent of $P_{17}^2$
and derives $\Lambda_\tau$ from the phi-oscillator and $\Sint$ directly.

\bigskip\noindent The two paths for $\Lambda_\tau$ converge to $3481=59^2$
from orthogonal directions: the primary path accesses it as the square of the
Structural Sum Prime; the secondary path accesses it via the phi-oscillator
descent and the interaction surface.  Their agreement constitutes a two-path
confirmation of the tau mass integer.


\subsection{Python implementation of the canonical PMNS sector}
\label{app:pmns-python}

\begin{verbatim}
from fractions import Fraction

N = 3
N_E = 4
N_R = 3
Delta = 5
T_vac = 11
tau10 = 10
Sigma_imb = 7
D_gauge = 12
G_mg = 2
Pi3 = 13
PG_Sigma_imb = 17
T_bary = 33

def frac(a, b):
    return Fraction(a, b)

def cross_det(x, y):
    a, b = x.numerator, x.denominator
    c, d = y.numerator, y.denominator
    return c * b - a * d

theta23 = frac(N * N, N_E * N_E)
theta12 = frac(N_R, tau10)
theta13 = frac(1, Delta * N * N)

aux23 = frac(N * N, N_E * N_E)
aux12 = frac(PG_Sigma_imb, Delta * T_vac)
aux13 = frac(Pi3, G_mg * N**3 * T_vac)

assert theta23 == frac(9, 16)
assert theta12 == frac(3, 10)
assert theta13 == frac(1, 45)

assert theta23 == frac(Sigma_imb + G_mg, D_gauge + N_E)
assert theta12 == frac(N, G_mg * Delta)
assert theta13 == frac(1, T_bary + D_gauge)

b23 = frac(Sigma_imb, D_gauge)
b12 = frac(N_R, tau10)
b13 = frac(1, Delta * N * N)

assert cross_det(b13, aux13) == -N * N
assert cross_det(b12, aux12) == Delta
assert cross_det(b23, aux23) == -N_E

pdg_theta12 = 0.307
pdg_theta23 = 0.561
pdg_theta13 = 0.02195

residuals = {
    "theta12": abs(float(theta12) - pdg_theta12),
    "theta23": abs(float(theta23) - pdg_theta23),
    "theta13": abs(float(theta13) - pdg_theta13),
}

print("canonical")
print(theta23, float(theta23))
print(theta12, float(theta12))
print(theta13, float(theta13))
print("auxiliary")
print(aux23, float(aux23))
print(aux12, float(aux12))
print(aux13, float(aux13))
print("cross determinants")
print(cross_det(b13, aux13))
print(cross_det(b12, aux12))
print(cross_det(b23, aux23))
print("residuals")
print(residuals)
\end{verbatim}

Executing this code returns the canonical fractions
\[
\frac{9}{16}, \qquad \frac{3}{10}, \qquad \frac{1}{45},
\]
verifies the three auxiliary cross-determinant identities, and computes the
external residuals recorded in the verification box.


\end{document}